% Traduction : expression de l'accord et du désaccord — Espagnol Tle A — Cours signé OnBuch+
% Programme officiel MINESEC. Compiler avec : tectonic E15-acuerdo-y-desacuerdo.tex
\newcommand{\DOCMATIERE}{Espagnol}
\newcommand{\DOCNIVEAU}{Tle A}
\newcommand{\DOCMODULE}{Módulo 3 : Vie citoyenne et ouverture au monde}
\newcommand{\DOCLECON}{E15}
\newcommand{\DOCTITRE}{Traduction : expression de l'accord et du désaccord}
% OnBuch+ — préambule « cours signés » ENRICHI (compilé avec tectonic / XeLaTeX)
% Système visuel « Pop » d'OnBuch+ : fond crème (jamais blanc), bordures encre
% épaisses, ombres « dures » décalées, titres Archivo Black en capitales.
% Version MATHÉMATIQUES : graphes (pgfplots/tikz), boîtes pédagogiques
% (définition, propriété, méthode, attention, le savais-tu, exercice résolu,
% exercice, TP, bilan), page de garde et signature OnBuch+ ancrée en bas.
%
% Macros attendues dans main.tex AVANT \input{preamble} :
%   \DOCMATIERE  \DOCNIVEAU  \DOCMODULE  \DOCLECON (numéro)  \DOCTITRE
\documentclass[11pt]{article}

\usepackage{fontspec}
\usepackage[spanish,french]{babel} % français = langue principale ; l'espagnol via \esp{..} / espagnol
\usepackage[a4paper,margin=2cm,top=3.4cm,bottom=2.8cm,headheight=1.4cm,footskip=1.3cm]{geometry}
\usepackage{amsmath,amssymb,amsfonts}
\usepackage{mathtools} % \xmapsto, \coloneqq... souvent utilisés par les modèles
\usepackage{cancel} % simplifications barrées (bilans de forces)
% \abs{z} (module, valeur absolue) et \norm{u} (norme), utilisés par les modèles
\providecommand{\abs}{}\let\abs\relax\DeclarePairedDelimiter{\abs}{\lvert}{\rvert}
\providecommand{\norm}{}\let\norm\relax\DeclarePairedDelimiter{\norm}{\lVert}{\rVert}
\usepackage[table]{xcolor} % [table] : fournit aussi \rowcolors (alternance de lignes)
\usepackage{pagecolor}
\usepackage{tikz}
\usetikzlibrary{arrows.meta,positioning,calc,shapes.geometric,shapes.misc,shapes.arrows,decorations.pathmorphing,decorations.pathreplacing,decorations.markings,patterns,shadows,mindmap,trees,fit,backgrounds,matrix,intersections,angles,quotes}
\usepackage{pgfplots}
\usepackage[version=4]{mhchem} % équations nucléaires \ce{^{235}_{92}U}
\usepackage[european,siunitx]{circuitikz} % schémas électriques
\pgfplotsset{compat=1.18}
\usepgfplotslibrary{fillbetween,groupplots} % aires entre courbes (calcul intégral)
\usepackage{enumitem}
\usepackage{fancyhdr}
% [most] charge toutes les bibliothèques tcolorbox (listings, minted, raster,
% documentation...) et gonfle inutilement la mémoire TeX sur un document long
% (~300 boîtes) : on ne charge que ce qui est réellement utilisé.
\usepackage[skins,breakable]{tcolorbox}
\usepackage{graphicx}
\usepackage{colortbl}
\usepackage{booktabs}
\usepackage{tabularx}
\usepackage{diagbox} % case d'en-tête diagonale des tableaux à double entrée
% Colonnes extensibles centrée / gauche / droite (C, L, R), souvent utilisées par les modèles
\newcolumntype{C}{>{\centering\arraybackslash}X}
\newcolumntype{L}{>{\raggedright\arraybackslash}X}
\newcolumntype{R}{>{\raggedleft\arraybackslash}X}
\usepackage{array}
\usepackage{multicol}
\usepackage{siunitx}
% parse-numbers=false : accepte aussi \SI{2,5 \times 10^{-3}}{...}
\sisetup{locale=FR,output-decimal-marker={,},group-minimum-digits=4,parse-numbers=false}
\DeclareSIUnit\ohm{\ensuremath{\symup{\Omega}}} % U+2126 absent de la police
\DeclareSIUnit\division{div} % réglages d'oscilloscope (ms/div, V/div)
\DeclareSIUnit\tr{tr} % tours (vitesse de rotation, ex. \SI{1500}{\tr\per\minute})
\DeclareSIUnit\molar{M} % concentration molaire (ex. \SI{3}{\molar}, \SI{1}{\micro\molar})
\DeclareSIUnit\an{an} % année (le modèle confond parfois avec \year, un compteur TeX) — ex. \SI{5}{\kilo\meter\per\an}
\DeclareSIUnit\FCFA{FCFA} % franc CFA (ex. \SI{500}{\FCFA\per\kilo\gram})
\DeclareSIUnit\degreeBrix{\textdegree Brix} % teneur en sucre d'un sirop/jus (réfractométrie)
\DeclareSIUnit\month{mois} % le modèle utilise parfois \month, un compteur TeX, comme unité de durée
\usepackage{qrcode}
% \degree : commande fréquemment utilisée par les modèles pour un angle ou une
% température en dehors de \SI{}{} (non fournie par les paquets chargés).
\providecommand{\degree}{^{\circ}}
% \qed (fin de démonstration), utilisé par les modèles sans amsthm
\providecommand{\qed}{\hfill$\square$}
% \barre : double barre des valeurs interdites dans un tableau de variations (tkz-tab)
\providecommand{\barre}{\ensuremath{\|}}
% environnement proof (amsthm non chargé) utilisé par les modèles
\ifdefined\proof\else\newenvironment{proof}[1][Démonstration]{\par\noindent\textit{#1.}\ }{\qed\par}\fi
\usepackage{lastpage}
\usepackage{hyperref}
\hypersetup{hidelinks}

% ============================================================
% Palette « Pop » OnBuch+ — mêmes hex que la classe P (plus_ui.dart)
% ============================================================
\definecolor{poporange}{HTML}{F59321}
\definecolor{popdark}{HTML}{E07A0C}
\definecolor{popcream}{HTML}{FFF6E8}
\definecolor{popink}{HTML}{1C1714}
\definecolor{poppurple}{HTML}{7A5AE0}
\definecolor{popgreen}{HTML}{2E9E6A}
\definecolor{poppink}{HTML}{E0457B}
\definecolor{popblue}{HTML}{2F7DE1}
\definecolor{popgold}{HTML}{C98A12}
\definecolor{popmuted}{HTML}{8A7F76}
\definecolor{popline}{HTML}{EADFD2}
\definecolor{popgray}{HTML}{8A7F76}
\colorlet{popgrey}{popgray}
% Alias tolérants (le modèle nomme parfois les couleurs différemment)
\colorlet{popwhite}{white}
\colorlet{popblack}{popink}
\colorlet{popred}{poppink}
\colorlet{popyellow}{popgold}
% Teintes claires (fonds de tableaux/encadrés) — le modèle nomme parfois
% les couleurs avec un suffixe L (ex. popblueL) sans les avoir définies.
\colorlet{poporangeL}{poporange!20}
\colorlet{popdarkL}{popdark!20}
\colorlet{poppurpleL}{poppurple!20}
\colorlet{popgreenL}{popgreen!20}
\colorlet{poppinkL}{poppink!20}
\colorlet{popblueL}{popblue!20}
\colorlet{popgoldL}{popgold!20}
\colorlet{popredL}{popred!20}
\colorlet{popgrayL}{popgray!20}
% Teintes claires pour fonds de tableaux / zones
\colorlet{popblueL}{popblue!10!white}
\colorlet{poporangeL}{poporange!14!white}
\colorlet{popgreenL}{popgreen!12!white}
\colorlet{poppurpleL}{poppurple!12!white}
\colorlet{poppinkL}{poppink!12!white}
\colorlet{popgoldL}{popgold!14!white}

\pagecolor{popcream}

% ============================================================
% Typographie
% ============================================================
\defaultfontfeatures{Ligatures=TeX}
\setmainfont{PlusJakartaSans-Medium.ttf}[Path=../../../fonts/,BoldFont=PlusJakartaSans-Bold.ttf]
% Maths en sans-serif (Fira Math) avec chiffres et lettres droites de Plus
% Jakarta, pour que les formules se fondent dans le texte.
\usepackage[math-style=ISO,bold-style=ISO]{unicode-math}
\setmathfont{FiraMath-Regular.otf}[Path=../../../fonts/]
\setmathfont{PlusJakartaSans-Medium.ttf}[Path=../../../fonts/,range={up/num,up/{Latin,latin},\mathup}]
\usepackage{icomma} % virgule décimale sans espace en mode math
% Caractères mathématiques Unicode absents des polices de texte (Plus Jakarta,
% Archivo Black) : rendus par leur équivalent mathématique, en texte comme en
% formule — sinon ils s'affichent en carré vide (ex. « Arithmétique dans ℤ »).
\usepackage{newunicodechar}
\newunicodechar{ℝ}{\ensuremath{\mathbb{R}}}
\newunicodechar{ℤ}{\ensuremath{\mathbb{Z}}}
\newunicodechar{ℂ}{\ensuremath{\mathbb{C}}}
\newunicodechar{ℕ}{\ensuremath{\mathbb{N}}}
\newunicodechar{ℚ}{\ensuremath{\mathbb{Q}}}
\newunicodechar{𝔻}{\ensuremath{\mathbb{D}}}
\newunicodechar{α}{\ensuremath{\alpha}}
\newunicodechar{β}{\ensuremath{\beta}}
\newunicodechar{γ}{\ensuremath{\gamma}}
\newunicodechar{δ}{\ensuremath{\delta}}
\newunicodechar{ε}{\ensuremath{\varepsilon}}
\newunicodechar{θ}{\ensuremath{\theta}}
\newunicodechar{λ}{\ensuremath{\lambda}}
\newunicodechar{μ}{\ensuremath{\mu}}
\newunicodechar{σ}{\ensuremath{\sigma}}
\newunicodechar{τ}{\ensuremath{\tau}}
\newunicodechar{φ}{\ensuremath{\varphi}}
\newunicodechar{ψ}{\ensuremath{\psi}}
\newunicodechar{ω}{\ensuremath{\omega}}
\newunicodechar{Σ}{\ensuremath{\Sigma}}
% majuscules produites par \MakeUppercase dans les titres (α-aminés -> Α)
\newunicodechar{Α}{\ensuremath{\alpha}}
\newunicodechar{Β}{\ensuremath{\beta}}
\newunicodechar{Γ}{\ensuremath{\Gamma}}
\newunicodechar{Δ}{\ensuremath{\Delta}}
\newunicodechar{Ε}{\ensuremath{\varepsilon}}
\newunicodechar{Θ}{\ensuremath{\Theta}}
\newunicodechar{Λ}{\ensuremath{\Lambda}}
\newunicodechar{Μ}{\ensuremath{\mu}}
\newunicodechar{Τ}{\ensuremath{\tau}}
\newunicodechar{Φ}{\ensuremath{\Phi}}
\newunicodechar{Ψ}{\ensuremath{\Psi}}
\newunicodechar{Ω}{\ensuremath{\Omega}}
\newunicodechar{↦}{\ensuremath{\mapsto}}
\newunicodechar{∈}{\ensuremath{\in}}
\newunicodechar{∉}{\ensuremath{\notin}}
\newunicodechar{∧}{\ensuremath{\wedge}}
\newunicodechar{⇔}{\ensuremath{\Leftrightarrow}}
\newunicodechar{⟺}{\ensuremath{\Longleftrightarrow}}
\newunicodechar{⇒}{\ensuremath{\Rightarrow}}
\newunicodechar{⟹}{\ensuremath{\Longrightarrow}}
\newunicodechar{∘}{\ensuremath{\circ}}
\newunicodechar{∪}{\ensuremath{\cup}}
\newunicodechar{∩}{\ensuremath{\cap}}
\newunicodechar{≡}{\ensuremath{\equiv}}
\newunicodechar{⊂}{\ensuremath{\subset}}
\newunicodechar{⊥}{\ensuremath{\perp}}
\newunicodechar{∃}{\ensuremath{\exists}}
\newunicodechar{∀}{\ensuremath{\forall}}
\newunicodechar{∛}{\ensuremath{\sqrt[3]{\,}}}
\newunicodechar{∜}{\ensuremath{\sqrt[4]{\,}}}
\newunicodechar{⃗}{\ensuremath{{}^{\to}}}
\newunicodechar{ⁿ}{\ifmmode^{n}\else\textsuperscript{\ensuremath{n}}\fi}
\newunicodechar{ˣ}{\ifmmode^{x}\else\textsuperscript{\ensuremath{x}}\fi}
\newunicodechar{ᵇ}{\ifmmode^{b}\else\textsuperscript{\ensuremath{b}}\fi}
\newunicodechar{ʳ}{\ifmmode^{r}\else\textsuperscript{\ensuremath{r}}\fi}
\newunicodechar{ᵘ}{\ifmmode^{u}\else\textsuperscript{\ensuremath{u}}\fi}
\newunicodechar{ʸ}{\ifmmode^{y}\else\textsuperscript{\ensuremath{y}}\fi}
\newunicodechar{ᵃ}{\ifmmode^{a}\else\textsuperscript{\ensuremath{a}}\fi}
\newunicodechar{ᵉ}{\ifmmode^{e}\else\textsuperscript{\ensuremath{e}}\fi}
\newunicodechar{ᵏ}{\ifmmode^{k}\else\textsuperscript{\ensuremath{k}}\fi}
\newunicodechar{ᵖ}{\ifmmode^{p}\else\textsuperscript{\ensuremath{p}}\fi}
\newunicodechar{ᴺ}{\ifmmode^{N}\else\textsuperscript{\ensuremath{N}}\fi}
\newunicodechar{ᶜ}{\ifmmode^{c}\else\textsuperscript{\ensuremath{c}}\fi}
\newunicodechar{ʰ}{\ifmmode^{h}\else\textsuperscript{\ensuremath{h}}\fi}
\newunicodechar{ᵠ}{\ifmmode^{q}\else\textsuperscript{\ensuremath{q}}\fi}
\newunicodechar{ₙ}{\ifmmode_{n}\else\textsubscript{\ensuremath{n}}\fi}
\newunicodechar{ₐ}{\ifmmode_{a}\else\textsubscript{\ensuremath{a}}\fi}
\newunicodechar{ᵢ}{\ifmmode_{i}\else\textsubscript{\ensuremath{i}}\fi}
\newunicodechar{ᵣ}{\ifmmode_{r}\else\textsubscript{\ensuremath{r}}\fi}
\newunicodechar{ₖ}{\ifmmode_{k}\else\textsubscript{\ensuremath{k}}\fi}
\newunicodechar{ᵦ}{\ifmmode_{\beta}\else\textsubscript{\ensuremath{\beta}}\fi}
\newunicodechar{ₓ}{\ifmmode_{x}\else\textsubscript{\ensuremath{x}}\fi}
\newfontfamily\popdisplay{ArchivoBlack-Regular.ttf}[Path=../../../fonts/]
\newfontfamily\poplogo{SpaceGrotesk-Bold.ttf}[Path=../../../fonts/]
\newfontfamily\popsemi{PlusJakartaSans-SemiBold.ttf}[Path=../../../fonts/]
\newfontfamily\popxbold{PlusJakartaSans-ExtraBold.ttf}[Path=../../../fonts/]
\newfontfamily\popxboldls{PlusJakartaSans-ExtraBold.ttf}[Path=../../../fonts/,LetterSpace=3.0]

\setlist[enumerate]{leftmargin=1.7em,itemsep=5pt,topsep=4pt}
\setlist[itemize]{leftmargin=1.7em,itemsep=4pt,topsep=4pt,label={\color{poporange}\textbullet}}
\setlength{\parindent}{0pt}
\setlength{\parskip}{5pt}
\renewcommand{\arraystretch}{1.25}

% pgfplots : style Pop par défaut
\pgfplotsset{
  every axis/.append style={axis line style={popink,line width=1pt},tick style={popink},
    grid style={popline,line width=0.5pt},label style={font=\small},tick label style={font=\footnotesize}},
  popcurve/.style={poporange,line width=1.8pt,no markers,smooth},
}
% Même style disponible dans un \draw TikZ simple (hors pgfplots)
\tikzset{popcurve/.style={poporange,line width=1.8pt}}
\tikzset{popgrid/.style={popline,very thin}}
\colorlet{popcurve}{poporange} % le nom du style est parfois utilisé comme couleur

% ============================================================
% Pill / squiggle
% ============================================================
\newcommand{\poppill}[3]{%
  \tikz[baseline=-0.55ex]\node[draw=popink,line width=1.2pt,rounded corners=2.6pt,%
    fill=#2,inner xsep=6.5pt,inner ysep=3pt]{%
    \popxboldls\fontsize{7}{7}\selectfont\color{#3}\MakeUppercase{#1}};%
}
\newcommand{\popsquiggle}[1]{%
  \tikz[baseline=-0.4ex]{
    \draw[#1,line width=2.6pt,line cap=round]
      plot [smooth] coordinates {(0,0.09) (0.34,0.01) (0.68,0.21) (1.02,0.01) (1.36,0.10)};
  }%
}
% Mot-clé surligné (marqueur orange sous le texte)
\newcommand{\cle}[1]{{\popxbold\color{popdark}#1}}

% ============================================================
% En-tête / pied
% ============================================================
\pagestyle{fancy}
\fancyhf{}
\renewcommand{\headrulewidth}{0pt}
\renewcommand{\footrulewidth}{0pt}
\lhead{\raisebox{-0.15ex}{\poplogo\fontsize{13}{13}\selectfont\color{popink}OnBuch\color{poporange}+}}
\rhead{\poppill{\DOCMATIERE\ \textbullet\ \DOCNIVEAU\ \textbullet\ Leçon \DOCLECON}{white}{popink}}
\lfoot{\popxbold\fontsize{7.6}{7.6}\selectfont\color{popmuted}\MakeUppercase{Cours signé OnBuch+}\ \textbullet\ {\popsemi\DOCTITRE}}
\rfoot{\tikz[baseline=-0.6ex]\node[rounded corners=8.5pt,fill=popink,minimum height=17pt,minimum width=17pt,inner xsep=5pt,inner sep=0]{\popxbold\fontsize{8}{8}\selectfont\color{popcream}\thepage\,/\,\pageref*{LastPage}};}

% ============================================================
% Titres
% ============================================================
\newcommand{\chaptitre}[2]{%
  \begin{tcolorbox}[enhanced,colback=poporange,colframe=popink,boxrule=2.6pt,arc=11pt,
      drop shadow=popink,left=15pt,right=15pt,top=12pt,bottom=11pt,before skip=3pt,after skip=18pt]
    \poppill{#2}{popink}{popcream}\par\vspace{9pt}
    {\raggedright\hyphenpenalty=10000\exhyphenpenalty=10000\popdisplay\fontsize{22}{24}\selectfont\color{popcream}\MakeUppercase{#1}\par}
  \end{tcolorbox}
}
% Section numérotée : pastille orange avec le numéro + titre Archivo
\newcounter{popsec}
\newcommand{\coursec}[1]{%
  \stepcounter{popsec}\setcounter{popsub}{0}%
  \par\vspace{16pt}\noindent
  \begin{minipage}{\linewidth}
  \tikz[baseline=-0.7ex]\node[draw=popink,line width=1.6pt,fill=poporange,rounded corners=4pt,minimum size=22pt,inner sep=0]{\popdisplay\fontsize{12}{12}\selectfont\color{popcream}\thepopsec};%
  \hspace{8pt}{\popdisplay\fontsize{15}{17}\selectfont\color{popink}\MakeUppercase{#1}}\par
  \vspace{4pt}\noindent{\color{poporange}\rule{36pt}{3.6pt}}
  \end{minipage}\par\nopagebreak\vspace{8pt}
}
\newcounter{popsub}
\newcommand{\courssub}[1]{%
  \stepcounter{popsub}%
  \par\vspace{9pt}\noindent{\popxbold\fontsize{11.5}{13}\selectfont\color{popink}{\color{poporange}\thepopsec.\thepopsub}\hspace{6pt}#1}\par\nopagebreak\vspace{4pt}
}

% ============================================================
% Boîtes pédagogiques — même recette : fond blanc, bordure épaisse colorée,
% ombre dure encre, badge plein. Titre optionnel [#1].
% ============================================================
\tcbset{popbase/.style={breakable,enhanced,colback=white,boxrule=2.3pt,arc=9pt,
  shadow={3pt}{-3pt}{0pt}{popink},left=14pt,right=14pt,top=11pt,bottom=11pt,before skip=11pt,after skip=15pt,
  fontupper=\popsemi\fontsize{10.6}{14.5}\selectfont\color{popink},
  fontlower=\popsemi\fontsize{10.6}{14.5}\selectfont\color{popink}}}
% Coche dessinée (le glyphe ✓ n'existe pas dans les polices du document)
\newcommand{\popcheck}[1]{\tikz[baseline=-0.5ex]\draw[#1,line width=1.6pt,line cap=round,line join=round](-0.13,0.0)--(-0.04,-0.09)--(0.13,0.1);}
\providecommand{\coche}{\popcheck{popgreen}}
\providecommand{\resultat}[1]{\par\smallskip\noindent\fcolorbox{popgreen}{popgreenL}{\parbox{\dimexpr\linewidth-2\fboxsep-2\fboxrule\relax}{\textbf{Résultat.} #1}}\par\smallskip}
\newcommand{\popboxhead}[3]{\poppill{#1}{#2}{white}\ifx\relax#3\relax\else\hspace{6pt}{\popxbold\fontsize{10.6}{12}\selectfont\color{popink}#3}\fi\par\vspace{7pt}}

\newtcolorbox{aretenir}[1][]{popbase,colframe=popgreen,before upper={\popboxhead{À retenir}{popgreen}{#1}}}
% Texte en espagnol (document de lecture, dialogue, poème, extrait) : cadre
% d'étude. Un titre optionnel (source, auteur) s'affiche dans l'en-tête.
\newtcolorbox{textebox}[1][]{popbase,colframe=popblue,before upper={\popboxhead{Texto}{popblue}{#1}\begin{otherlanguage}{spanish}},after upper={\end{otherlanguage}}}
% Marques ✓ / ✗ (forme correcte / fautive) souvent utilisées par les modèles
\providecommand{\cross}{\textcolor{poppink}{\ensuremath{\times}}}
\providecommand{\xmark}{\cross}\providecommand{\crossmark}{\cross}\providecommand{\cmark}{\ensuremath{\checkmark}}
% Ligne de réponse / justification à compléter (inventée par les modèles)
\providecommand{\justification}[1]{\par\noindent\textit{Justification~:} \dotfill\par}
% \textipa (package tipa non chargé) : la police ne couvre pas l'API -> simple passage
\providecommand{\textipa}[1]{#1}
% Soulignement (ulem non chargé) : \uline, \ul
\providecommand{\uline}[1]{\underline{#1}}\providecommand{\ul}[1]{\underline{#1}}
% Mot ou phrase en espagnol à l'intérieur d'un paragraphe français
\newcommand{\esp}[1]{\ifmmode\text{\textit{#1}}\else\begingroup\selectlanguage{spanish}\itshape #1\endgroup\fi}
\newtcolorbox{exemplebox}[1][]{popbase,colframe=poporange,before upper={\popboxhead{Exemple}{poporange}{#1}}}
\newtcolorbox{definition}[1][]{popbase,colframe=popblue,before upper={\popboxhead{Définition}{popblue}{#1}}}
\newtcolorbox{propriete}[1][]{popbase,colframe=poppurple,before upper={\popboxhead{Propriété}{poppurple}{#1}}}
\newtcolorbox{methode}[1][]{popbase,colframe=popgold,before upper={\popboxhead{Méthode}{popgold}{#1}}}
\newtcolorbox{attention}[1][]{popbase,colframe=poppink,before upper={\popboxhead{Attention}{poppink}{#1}}}
\newtcolorbox{experience}[1][]{popbase,colframe=popblue,colback=popblueL,before upper={\popboxhead{Expérience}{popblue}{#1}}}
% « Le savais-tu ? » : carte violette pleine (contexte, culture, Cameroun)
\newtcolorbox{savaistu}[1][]{popbase,colback=poppurple,colframe=popink,
  fontupper=\popsemi\fontsize{10.4}{14.2}\selectfont\color{white},
  before upper={\poppill{Le savais-tu\,?}{white}{poppurple}\ifx\relax#1\relax\else\hspace{6pt}{\popxbold\color{white}#1}\fi\par\vspace{7pt}}}
% Exercice résolu : énoncé en haut, \tcblower puis la solution sur fond crème
\newcounter{popexo}\newcounter{popexr}
\newtcolorbox{exoresolu}[1][]{popbase,colframe=poporange,colbacklower=popcream,
  segmentation style={popink,line width=1.2pt,dashed},
  before upper={\stepcounter{popexr}\popboxhead{Exercice résolu \thepopexr}{poporange}{#1}},
  before lower={\poppill{Solution}{popink}{popcream}\par\vspace{6pt}}}
% Exercice (non résolu en place) — la correction est dans la partie Corrigés
\newtcolorbox{exercice}[1][]{popbase,colframe=popink,
  before upper={\stepcounter{popexo}\popboxhead{Exercice \thepopexo}{popink}{#1}}}
% Corrigé
\newtcolorbox{corrige}[1][]{popbase,colframe=popgreen,colback=popgreenL,
  before upper={\popboxhead{Corrigé}{popgreen}{#1}}}
% Objectifs / prérequis / bilan
\newtcolorbox{objectifs}{popbase,colframe=popink,colback=poporangeL,
  before upper={\popboxhead{Objectifs de la leçon}{popink}{}}}
\newtcolorbox{prerequis}{popbase,colframe=popink,colback=popblueL,
  before upper={\popboxhead{Prérequis}{popblue}{}}}
\newtcolorbox{bilan}{popbase,colframe=popink,colback=white,boxrule=2.8pt,
  before upper={\popboxhead{Fiche bilan}{poporange}{L'essentiel à connaître pour l'examen}}}
% Figure encadrée avec légende
\newtcolorbox{popfigure}[1][]{enhanced,breakable=false,colback=white,colframe=popline,boxrule=1.4pt,arc=8pt,
  left=10pt,right=10pt,top=10pt,bottom=8pt,before skip=10pt,after skip=12pt,halign=center,
  lower separated=false,halign lower=center,
  fontlower=\popsemi\fontsize{9}{11.5}\selectfont\color{popmuted},
  #1}
\newcommand{\legende}[1]{\par\vspace{6pt}{\popsemi\fontsize{9}{11.5}\selectfont\color{popmuted}#1}}

% ============================================================
% Signature OnBuch+
% ============================================================
\newcommand{\poplogoinline}[1]{{\poplogo\fontsize{#1}{#1}\selectfont\color{popink}OnBuch\color{poporange}+}}
\newcommand{\signatureonbuch}{%
  \begin{tcolorbox}[enhanced,colback=popink,colframe=popink,boxrule=2.2pt,arc=10pt,
      drop shadow=poporange,left=14pt,right=14pt,top=10pt,bottom=10pt,before skip=0pt,after skip=0pt]
    \tikz[baseline=-0.7ex]\node[circle,fill=popgreen,draw=popcream,line width=1.3pt,minimum size=22pt,inner sep=0]{\popcheck{white}};%
    \hspace{9pt}\begin{minipage}[c]{0.62\linewidth}
      {\popxbold\fontsize{12}{13}\selectfont\color{popcream}Signé\ \textbullet\ {\poplogo OnBuch\color{poporange}+}}\par\vspace{2pt}
      {\raggedright\popsemi\fontsize{8}{10.5}\selectfont\color{popline}\MakeUppercase{Équipe pédagogique OnBuch+}\\\MakeUppercase{Conforme au programme officiel MINESEC}\par}
    \end{minipage}\hfill
    {\poplogo\fontsize{22}{22}\selectfont\color{popcream}OnBuch\color{poporange}+}
  \end{tcolorbox}%
}

% ============================================================
% Page de garde
% #1 titre · #2 matière · #3 niveau · #4 module · #5 n° leçon · #6 durée ·
% #7 description / objectifs (texte ou itemize)
% ============================================================
\newcommand{\pagegarde}[7]{%
  \thispagestyle{empty}
  \newgeometry{a4paper,margin=1.7cm,top=1.6cm,bottom=1.4cm}%
  \begin{tikzpicture}[remember picture,overlay]
    \fill[poporange!18] ([xshift=-3.2cm,yshift=-3.6cm]current page.north east) circle (4.6cm);
    \fill[poppurple!14] ([xshift=2.4cm,yshift=7.6cm]current page.south west) circle (3.8cm);
  \end{tikzpicture}%
  {\poplogo\fontsize{30}{30}\selectfont\color{popink}OnBuch\color{poporange}+}\hfill
  \raisebox{0.6ex}{\poppill{Cours signé \textbullet\ Premium}{popink}{popcream}}\par
  \vspace{1pt}\popsquiggle{poppurple}\par
  \vspace{8pt}
  \begin{tcolorbox}[enhanced,colback=poporange,colframe=popink,boxrule=2.8pt,arc=13pt,
      drop shadow=popink,left=18pt,right=18pt,top=12pt,bottom=13pt,after skip=10pt]
    \poppill{#2 \textbullet\ #3}{popink}{popcream}\hspace{5pt}\poppill{#4}{white}{popink}\par\vspace{8pt}
    {\popdisplay\fontsize{14}{14}\selectfont\color{popink}LEÇON #5}\par\vspace{4pt}
    {\raggedright\hyphenpenalty=10000\exhyphenpenalty=10000\popdisplay\fontsize{25}{27}\selectfont\color{popcream}\MakeUppercase{#1}\par}
    \vspace{8pt}
    {\popxbold\fontsize{9.5}{9.5}\selectfont\color{popink}\MakeUppercase{Durée conseillée : #6}}
  \end{tcolorbox}
  \begin{tcolorbox}[enhanced,colback=white,colframe=popink,boxrule=2.2pt,arc=10pt,
      drop shadow=popink,left=16pt,right=16pt,top=10pt,bottom=10pt,after skip=8pt]
    \poppill{Dans cette leçon}{poppurple}{white}\par\vspace{5pt}
    \popsemi\fontsize{9.3}{12.6}\selectfont\color{popink}#7
  \end{tcolorbox}
  \noindent\begin{minipage}[t]{0.487\linewidth}
    \begin{tcolorbox}[enhanced,colback=popgreen,colframe=popink,boxrule=2.2pt,arc=10pt,
        drop shadow=popink,left=11pt,right=11pt,top=10pt,bottom=10pt,height=3.1cm,valign=center]
      \tcbox[colback=white,colframe=popink,boxrule=1.3pt,arc=5pt,boxsep=0pt,nobeforeafter,
        left=3pt,right=3pt,top=3pt,bottom=3pt,valign=center]{\qrcode[height=1.9cm]{https://play.google.com/store/apps/details?id=com.luvvix.onbuch&hl=fr}}%
      \hspace{8pt}\begin{minipage}[c]{0.5\linewidth}
        {\popdisplay\fontsize{11}{12.5}\selectfont\color{white}\MakeUppercase{Télécharge OnBuch}}\par\vspace{4pt}
        {\raggedright\popsemi\fontsize{8.4}{11}\selectfont\color{white}Gratuit \textbullet\ Google Play\\Tuteur IA, annales, concours, résultats\par}
      \end{minipage}
    \end{tcolorbox}
  \end{minipage}\hfill
  \begin{minipage}[t]{0.487\linewidth}
    \begin{tcolorbox}[enhanced,colback=poppurple,colframe=popink,boxrule=2.2pt,arc=10pt,
        drop shadow=popink,left=11pt,right=11pt,top=10pt,bottom=10pt,height=3.1cm,valign=center]
      \tikz[baseline=-0.7ex]\node[circle,fill=white,draw=popink,line width=1.3pt,minimum size=1.6cm,inner sep=0]{\popdisplay\fontsize{24}{24}\selectfont\color{poppurple}+};%
      \hspace{8pt}\begin{minipage}[c]{0.6\linewidth}
        {\popdisplay\fontsize{11}{12.5}\selectfont\color{white}\MakeUppercase{Passe à OnBuch+}}\par\vspace{4pt}
        {\raggedright\popsemi\fontsize{8.4}{11}\selectfont\color{white}Tous les cours signés, Tuteur IA illimité, fiches PDF — onglet OnBuch+ de l'app\par}
      \end{minipage}
    \end{tcolorbox}
  \end{minipage}
  \vspace{8pt}
  \popcontenu
  \vfill
  \signatureonbuch
  \restoregeometry
  \newpage
}

% Rangée « Ce cours contient » (page de garde)
\newcommand{\poptile}[3]{% #1 couleur #2 gros texte #3 légende
  \begin{tcolorbox}[enhanced,colback=#1,colframe=popink,boxrule=2pt,arc=9pt,shadow={2.5pt}{-2.5pt}{0pt}{popink},
      left=6pt,right=6pt,top=9pt,bottom=9pt,halign=center,valign=center,height=1.75cm,width=\linewidth,nobeforeafter]
    {\popdisplay\fontsize{12.5}{13}\selectfont\color{white}\MakeUppercase{#2}}\par\vspace{3pt}
    {\centering\popsemi\fontsize{7.8}{9.5}\selectfont\color{white}#3\par}
  \end{tcolorbox}}
\newcommand{\popcontenu}{%
  \par\noindent{\popxboldls\fontsize{7.5}{7.5}\selectfont\color{popmuted}CE COURS SIGNÉ CONTIENT}\par\vspace{7pt}
  \noindent\begin{minipage}[t]{0.235\linewidth}\poptile{popblue}{Cours}{complet, illustré, pas à pas}\end{minipage}\hfill
  \begin{minipage}[t]{0.235\linewidth}\poptile{poporange}{Méthodes}{et exercices résolus}\end{minipage}\hfill
  \begin{minipage}[t]{0.235\linewidth}\poptile{poppink}{Exercices}{type examen + corrigés}\end{minipage}\hfill
  \begin{minipage}[t]{0.235\linewidth}\poptile{popgreen}{Fiche bilan}{l'essentiel à retenir}\end{minipage}\par}

% Sommaire
\newtcolorbox{sommaire}{enhanced,colback=white,colframe=popink,boxrule=2.2pt,arc=10pt,
  drop shadow=popink,left=18pt,right=18pt,top=13pt,bottom=13pt,before skip=6pt,after skip=20pt,
  before upper={\poppill{Sommaire}{poppurple}{white}\par\vspace{9pt}\popsemi\fontsize{10.6}{14.5}\selectfont\color{popink}}}

% Signature de fin de cours — poussée tout en bas de la dernière page
\newcommand{\findecours}{%
  \par\vfill\nopagebreak
  \begin{center}\popsquiggle{poporange}\end{center}\vspace{4pt}
  \signatureonbuch
}
\AtBeginDocument{\def\llbracket{\mathopen{[\mkern-3.5mu[}}\def\rrbracket{\mathclose{]\mkern-3.5mu]}}} % intervalles d'entiers (glyphe ⟦ absent de la police)
% Glyphes absents de Fira Math : redéfinis à partir de glyphes disponibles
\AtBeginDocument{%
  \def\setminus{\mathbin{\backslash}}\let\smallsetminus\setminus
  \def\vdots{\mathord{\vbox{\baselineskip3.2pt\lineskiplimit0pt\kern4pt\hbox{.}\hbox{.}\hbox{.}}}}%
  \def\ddots{\mathinner{\mkern1mu\raise6.5pt\hbox{.}\mkern2mu\raise3.6pt\hbox{.}\mkern2mu\raise0.7pt\hbox{.}\mkern1mu}}%
  \def\triangle{\mathord{\scalebox{0.9}{$\symup{\Delta}$}}}%
  \def\nabla{\mathord{\raisebox{\depth}{\scalebox{1}[-1]{$\symup{\Delta}$}}}}%
  \def\checkmark{\popcheck{popdark}}%
  \def\star{\mathbin{\ast}}\def\bigstar{\mathord{\ast}}%
  \def\diamond{\mathbin{\scalebox{0.6}{$\Diamond$}}}%
  \def\bigcup{\mathop{\mathchoice{\raisebox{-0.3em}{\scalebox{1.7}{$\cup$}}}{\scalebox{1.25}{$\cup$}}{\cup}{\cup}}\displaylimits}%
  \def\bigcap{\mathop{\mathchoice{\raisebox{-0.3em}{\scalebox{1.7}{$\cap$}}}{\scalebox{1.25}{$\cap$}}{\cap}{\cap}}\displaylimits}%
  \def\lesseqqgtr{\mathrel{\lessgtr}}\def\bigoplus{\mathop{\oplus}\displaylimits}%
}
\providecommand{\ssi}{si et seulement si} % abréviation fréquente
\AtBeginDocument{\providecommand{\wideparen}{\overparen}} % arc de cercle

%% FIGFIX-BEGIN v3 — correctifs globaux des figures (docs/onbuch-generation/tools/figfix)
\usepackage{adjustbox}\usepackage{etoolbox}
% toute figure TikZ/pgfplots plus large que la ligne est réduite à la largeur disponible
\BeforeBeginEnvironment{tikzpicture}{\begin{adjustbox}{max width=\linewidth}}
\AfterEndEnvironment{tikzpicture}{\end{adjustbox}}
% légendes pgfplots lisibles par-dessus les courbes
\pgfplotsset{every axis/.append style={legend style={fill=white,fill opacity=0.92,text opacity=1,draw=black!15,inner sep=3pt}}}
% étiquettes d'arêtes (syntaxe edge["..."]) avec fond blanc
\tikzset{every edge quotes/.append style={fill=white,inner sep=1.2pt}}
% bulles de cartes mentales : texte à la ligne dans la bulle, bulle un peu plus grande
\tikzset{every concept/.append style={text width=2.0cm,align=center,minimum size=2.3cm,inner sep=2pt}}
%% FIGFIX-END
\begin{document}

\pagegarde{Traduction : expression de l'accord et du désaccord}{Espagnol}{Tle A}{Módulo 3 : Vie citoyenne et ouverture au monde}{E15}{2 h}{Leçon de traduction consacrée à l'expression de l'accord et du désaccord en espagnol : formules figées, constructions verbales, modalisation et concession. Elle prépare au thème et à la version du Bac par l'observation d'exemples, des règles contrastives français/espagnol et des exercices de traduction de phrases d'opinion.
\vspace{4pt}
\begin{multicols}{2}\small
\begin{itemize}
\item À la fin de cette leçon, je sais traduire et employer les formules d'accord : estar de acuerdo, tener razón, por supuesto, efectivamente, claro (que sí)
\item Je sais traduire et employer les formules de désaccord : no estar de acuerdo, de ninguna manera, al contrario, oponerse a, ni mucho menos
\item Je sais conjuguer et construire correctamente estar de acuerdo con / en que et oponerse a + infinitif ou nom
\item Je sais nuancer une opinion grâce aux tournures concessives (aunque, es verdad que... pero, puede que + subjonctif)
\item Je sais atténuer une opinion avec me parece que, creo que, quizás, en cierto modo
\item Je sais repérer et éviter les pièges de traduction (faux amis, constructions à préposition, subjonctif après no creer que)
\end{itemize}\end{multicols}}

\begin{sommaire}
\begin{multicols}{2}
\begin{enumerate}
\item Observation : l'accord et le désaccord en contexte
\item Les formules d'accord : formes et emplois
\item Les formules de désaccord : formes et emplois
\item Modalisation : atténuer et concéder
\item Contrastif français/espagnol et pièges de traduction
\item Entraînement à la traduction : thème et version
\item Activité d'intégration
\item Méthodes et exercices résolus
\item Exercices
\item Corrigés des exercices
\item Fiche bilan
\end{enumerate}
\end{multicols}
\end{sommaire}

\begin{objectifs}
\begin{itemize}
\item À la fin de cette leçon, je sais traduire et employer les formules d'accord : estar de acuerdo, tener razón, por supuesto, efectivamente, claro (que sí)
\item Je sais traduire et employer les formules de désaccord : no estar de acuerdo, de ninguna manera, al contrario, oponerse a, ni mucho menos
\item Je sais conjuguer et construire correctamente estar de acuerdo con / en que et oponerse a + infinitif ou nom
\item Je sais nuancer une opinion grâce aux tournures concessives (aunque, es verdad que... pero, puede que + subjonctif)
\item Je sais atténuer une opinion avec me parece que, creo que, quizás, en cierto modo
\item Je sais repérer et éviter les pièges de traduction (faux amis, constructions à préposition, subjonctif après no creer que)
\item Je sais traduire du français vers l'espagnol et inversement des phrases d'opinion complètes
\item Je sais réagir à l'oral et à l'écrit dans un débat en exprimant accord, désaccord et nuance
\end{itemize}
\end{objectifs}

\begin{prerequis}
\begin{itemize}
\item Présent de l'indicatif des verbes réguliers et irréguliers (estar, tener, oponerse)
\item Le subjonctif présent après les verbes d'opinion niés (no creer que + subjuntivo)
\item Les verbes pronominaux (oponerse, parecerse)
\item Les pronoms compléments et la construction gustar / parecer (me parece que)
\item Les connecteurs simples (pero, sin embargo, aunque) vus en Première
\end{itemize}
\end{prerequis}

\newpage

\coursec{Observation : l'accord et le désaccord en contexte}

En classe de Terminale à Yaoundé, un débat s'engage : « Faut-il interdire les téléphones portables au lycée ? » Certains élèves approuvent, d'autres protestent vivement. Comment dire en espagnol « je suis d'accord », « tu as raison », « pas du tout » ou « je m'y oppose » sans se tromper de construction ? Un débat télévisé espagnol ou un forum de discussion latino-américain exige de nuancer : dire « je suis d'accord, mais... ». Cette leçon vous donne toutes les clés pour traduire et manier l'accord et le désaccord comme un hispanophone.

\textbf{Problématique :} quelles formules l'espagnol offre-t-il pour exprimer l'accord et le désaccord, avec quelles constructions, et comment choisir la bonne formule selon le degré d'adhésion ou de refus ?

\courssub{Lecture : un débat entre deux lycéens}

Avant toute règle, observons la langue en action. Voici un dialogue original entre deux élèves espagnols, Lucía et Marcos, qui discutent dans la cour du lycée.

\begin{textebox}[Un debate en el patio del instituto]
— Lucía, ¿crees que los alumnos deben usar el móvil en clase?\\
— ¡De ninguna manera! Al contrario, distrae mucho. Los alumnos miran las pantallas y no escuchan al profesor.\\
— Pues yo no estoy de acuerdo contigo: el móvil puede ser una herramienta útil. Por ejemplo, sirve para buscar información o traducir palabras.\\
— Tienes razón en parte, pero hay que controlarlo. Si no, los chicos solo juegan y mandan mensajes.\\
— Efectivamente, eso es un problema. Pero yo me opongo a una prohibición total: sería mejor enseñar a usarlo bien.\\
— Claro que sí, la educación es importante. Exacto: el móvil no es el enemigo, el problema es el mal uso.\\
— ¡Eso es! Entonces, ¿estás de acuerdo en que debemos proponer reglas claras al director?\\
— Por supuesto. Estoy de acuerdo con esa idea. Podemos escribir una carta juntos.\\
— ¡Qué va! ¿Una carta? Mejor un debate con toda la clase.\\
— Bueno, no estoy de acuerdo del todo, pero acepto: primero el debate, después la carta.
\end{textebox}

\textbf{Glossaire :}
\begin{itemize}
\item \esp{distrae} : il distrait (verbe \esp{distraer}, irrégulier comme \esp{traer} : \emph{yo distraigo, tú distraes}...)
\item \esp{una herramienta útil} : un outil utile
\item \esp{hay que controlarlo} : il faut le contrôler (\esp{hay que + infinitif} = il faut..., impersonnel)
\item \esp{mandar mensajes} : envoyer des messages
\item \esp{una prohibición total} : une interdiction totale
\item \esp{el mal uso} : le mauvais usage
\item \esp{proponer reglas claras} : proposer des règles claires
\item \esp{no estoy de acuerdo del todo} : je ne suis pas entièrement d'accord
\item \esp{sería mejor} : ce serait mieux (conditionnel de \esp{ser})
\end{itemize}

\textbf{Questions de compréhension :}
\begin{enumerate}
\item ¿Qué piensa Lucía del móvil en clase al principio del diálogo?
\item ¿Está Marcos de acuerdo con ella? ¿Qué argumento da?
\item ¿Sobre qué punto se ponen de acuerdo al final?
\end{enumerate}

\begin{corrige}[Questions de compréhension]
\begin{enumerate}
\item Lucía piensa que los alumnos no deben usar el móvil en clase porque distrae mucho.
\item No, Marcos no está de acuerdo: piensa que el móvil puede ser una herramienta útil para buscar información.
\item Al final se ponen de acuerdo en organizar un debate con la clase y luego escribir una carta al director.
\end{enumerate}
\end{corrige}

\courssub{Repérage : les formules d'accord}

Relevons dans le dialogue toutes les formules qui expriment l'\cle{accord} :

\begin{definition}[Les formules d'accord]
\begin{itemize}
\item \esp{Estoy de acuerdo (contigo / con esa idea)} : je suis d'accord (avec toi / avec cette idée). C'est la formule de base, neutre, utilisable partout.
\item \esp{Tienes razón} : tu as raison. Très fréquent à l'oral ; on peut nuancer : \esp{tienes razón en parte} (tu as raison en partie).
\item \esp{Por supuesto} : bien sûr, évidemment. Accord total et enthousiaste.
\item \esp{Efectivamente} : effectivement, en effet. Confirme ce que vient de dire l'autre.
\item \esp{Claro que sí} : bien sûr que oui, évidemment. Familier et chaleureux.
\item \esp{Exacto} / \esp{¡Eso es!} : exactement ! / c'est ça ! Accord vif à l'oral.
\end{itemize}
\end{definition}

\begin{exemplebox}[Exemples traduits]
\esp{Estoy de acuerdo contigo.} « Je suis d'accord avec toi. »\\
\esp{Tienes razón: el móvil distrae.} « Tu as raison : le portable distrait. »\\
\esp{— ¿Vienes al debate? — Por supuesto.} « — Tu viens au débat ? — Bien sûr. »\\
\esp{Efectivamente, hay que controlarlo.} « Effectivement, il faut le contrôler. »\\
\esp{— ¿Estás de acuerdo? — Claro que sí.} « — Tu es d'accord ? — Bien sûr que oui. »
\end{exemplebox}

\courssub{Repérage : les formules de désaccord}

\begin{definition}[Les formules de désaccord]
\begin{itemize}
\item \esp{No estoy de acuerdo} : je ne suis pas d'accord. Formule neutre et polie.
\item \esp{De ninguna manera} : en aucune façon, pas question. Refus catégorique.
\item \esp{Al contrario} : au contraire. Sert à renverser l'argument de l'autre.
\item \esp{Me opongo (a eso / a esa idea)} : je m'oppose (à cela / à cette idée). Désaccord ferme, soutenu.
\item \esp{Ni mucho menos} : loin de là, pas du tout. Familier.
\item \esp{¡Qué va!} : mais non ! / n'importe quoi ! Très familier, à réserver aux conversations entre amis.
\end{itemize}
\end{definition}

\begin{exemplebox}[Exemples traduits]
\esp{No estoy de acuerdo con esa propuesta.} « Je ne suis pas d'accord avec cette proposition. »\\
\esp{— ¿Debemos prohibir el móvil? — ¡De ninguna manera!} « — Faut-il interdire le portable ? — Pas question ! »\\
\esp{Al contrario, el móvil puede ayudar.} « Au contraire, le portable peut aider. »\\
\esp{Me opongo a esa idea.} « Je m'oppose à cette idée. »\\
\esp{— ¿Te gusta la idea? — ¡Qué va!} « — L'idée te plaît ? — Mais non ! »
\end{exemplebox}

\courssub{Classement selon le degré d'accord}

Toutes ces formules ne se valent pas : un bon hispanophone choisit selon l'\cle{intensité} de son adhésion ou de son refus.

\begin{center}
\begin{tabularx}{\linewidth}{|l|X|X|}
\hline
\rowcolor{popblueL} \textbf{Degré} & \textbf{Formules} & \textbf{Français} \\
\hline
Accord total & \esp{por supuesto, claro que sí, exacto, eso es, efectivamente} & bien sûr, évidemment, exactement, effectivement \\
\hline
Accord nuancé / partiel & \esp{estoy de acuerdo, tienes razón (en parte), hasta cierto punto} & je suis d'accord, tu as raison (en partie), jusqu'à un certain point \\
\hline
Désaccord poli & \esp{no estoy de acuerdo, no lo creo, no estoy de acuerdo del todo} & je ne suis pas d'accord, je ne le crois pas, pas tout à fait d'accord \\
\hline
Désaccord catégorique & \esp{de ninguna manera, ni mucho menos, me opongo, ¡qué va!} & pas question, loin de là, je m'y oppose, mais non ! \\
\hline
\end{tabularx}
\end{center}

\begin{popfigure}
\begin{tikzpicture}[mindmap, grow cyclic, every node/.style={concept}, root concept/.append style={concept color=popblue, font=\Large\bfseries}, level 1/.append style={level distance=4.2cm, sibling angle=180}, level 2/.append style={level distance=3.2cm, sibling angle=45}]
\node[root concept]{OPINIÓN}
  child[concept color=popgreen] { node {ACUERDO}
    child { node[align=center] {estoy de\\ acuerdo} }
    child { node[align=center] {tienes\\ razón} }
    child { node[align=center] {por\\ supuesto} }
    child { node {efectivamente} }
  }
  child[concept color=poporange] { node {DESACUERDO}
    child { node[align=center] {no estoy\\ de acuerdo} }
    child { node[align=center] {de ninguna\\ manera} }
    child { node {al contrario} }
    child { node {me opongo} }
  };
\end{tikzpicture}
\legende{Carte mentale : les deux familles de formules pour réagir à une opinion.}
\end{popfigure}

\begin{popfigure}
\begin{tikzpicture}
\draw[-{Stealth[length=3mm]}, line width=1.5pt, popmuted] (0,0) -- (12.5,0);
\foreach \x/\txt/\col in {0.6/{por supuesto\\ claro que sí}/popgreen, 4.2/{estoy de acuerdo\\ tienes razón en parte}/popgreen!60!popblue, 7.8/{no estoy de acuerdo\\ no lo creo}/poporange, 11.4/{de ninguna manera\\ me opongo ¡qué va!}/poporange!60!popink} {
  \fill[\col] (\x,0) circle (0.14);
  \node[align=center, above, font=\small] at (\x,0.2) {\txt};
}
\node[align=center, below, font=\small\bfseries, popgreen] at (0.6,-0.2) {accord total};
\node[align=center, below, font=\small\bfseries, popgreen!60!popblue] at (4.2,-0.2) {accord nuancé};
\node[align=center, below, font=\small\bfseries, poporange] at (7.8,-0.2) {désaccord poli};
\node[align=center, below, font=\small\bfseries, poporange!60!popink] at (11.4,-0.2) {désaccord total};
\end{tikzpicture}
\legende{Échelle graduée : de l'accord total au désaccord catégorique.}
\end{popfigure}

\courssub{Modalisation : atténuer et concéder}

Exprimer son opinion, ce n'est pas seulement dire « oui » ou « non ». En espagnol comme en français, on module son adhésion pour rester poli, diplomate ou persuasif.

\begin{propriete}[La concession : « c'est vrai, mais... »]
Pour reconnaître une part de vérité chez l'autre avant d'apporter sa propre idée, on utilise :
\begin{itemize}
\item \esp{Es cierto que... , pero...} / \esp{Es verdad que... , pero...} : « C'est vrai que... , mais... »
\item \esp{Admito que... , sin embargo...} : « J'admets que... , cependant... »
\item \esp{Puede ser, pero...} / \esp{Puede que tengas razón, pero...} : « C'est possible, mais... » / « Tu as peut-être raison, mais... » (\esp{puede que} + \textbf{subjonctif}).
\item \esp{Tienes razón en parte, pero...} : « Tu as raison en partie, mais... »
\end{itemize}
\end{propriete}

\begin{propriete}[L'atténuation : adoucir son désaccord]
Pour ne pas heurter, on évite le \esp{no} frontal et on préfère :
\begin{itemize}
\item \esp{No estoy muy de acuerdo...} / \esp{No estoy del todo de acuerdo...} : « Je ne suis pas très/entièrement d'accord... »
\item \esp{No lo veo muy claro...} : « Je ne le vois pas très clair... » (je ne suis pas convaincu).
\item \esp{Hasta cierto punto...} : « Jusqu'à un certain point... »
\item \esp{Me parece bien, aunque...} : « Ça me va, bien que... » (\esp{aunque} + \textbf{indicatif} pour fait connu, \textbf{subjonctif} pour hypothétique).
\end{itemize}
\end{propriete}

\begin{exemplebox}[Modalisation en contexte]
\esp{Es cierto que el móvil distrae, pero también es una herramienta de trabajo.} « C'est vrai que le portable distrait, mais c'est aussi un outil de travail. »\\
\esp{No estoy del todo de acuerdo con la prohibición total.} « Je ne suis pas entièrement d'accord avec l'interdiction totale. »\\
\esp{Puede que tengas razón, pero me opongo a que lo prohíban.} « Tu as peut-être raison, mais je m'oppose à ce qu'on l'interdise. » (\esp{me opongo a que} + subjonctif).
\end{exemplebox}

\courssub{Observation des constructions}

Regardons de près les structures grammaticales du dialogue : chaque verbe ou locution exige sa préposition.

\begin{propriete}[Les trois constructions fondamentales]
\begin{enumerate}
\item \esp{estar de acuerdo \textbf{con} + nom / pronom} : \esp{Estoy de acuerdo \textbf{con} Lucía / \textbf{contigo}.} (Je suis d'accord avec Lucía / avec toi.)
\item \esp{estar de acuerdo \textbf{en} + nom / \textbf{en que} + phrase} : \esp{Estamos de acuerdo \textbf{en} ese punto.} / \esp{Estamos de acuerdo \textbf{en que} hay que controlar el móvil.} (Nous sommes d'accord sur ce point / sur le fait qu'il faut contrôler le portable.)
\item \esp{oponerse \textbf{a} + nom / infinitif / \textbf{a que} + subjonctif} : \esp{Me opongo \textbf{a} esa idea.} / \esp{Me opongo \textbf{a} prohibir el móvil.} / \esp{Me opongo \textbf{a que} lo prohíban.} (Je m'oppose à cette idée / à interdire le portable / à ce qu'on l'interdise.)
\end{enumerate}
\end{propriete}

\begin{attention}[Le verbe ESTAR, pas SER !]
Ne traduis jamais « je suis d'accord » par *\esp{soy de acuerdo}. Le verbe correct est \cle{estar} : \esp{\textbf{estoy} de acuerdo}. Pourquoi ? Parce que l'accord est un \emph{état} momentané (je peux changer d'avis), non une caractéristique permanente de ma personnalité. C'est la même logique que \esp{estoy cansado} (je suis fatigué, en ce moment). De même : \esp{¿Estás de acuerdo?} et non *\esp{¿Eres de acuerdo?}.
\end{attention}

\begin{aretenir}[Constructions : français vs espagnol]
\begin{itemize}
\item \textbf{Accord avec quelqu'un} : FR « d'accord \textbf{avec} » $\rightarrow$ ES \esp{de acuerdo \textbf{con}}.
\item \textbf{Accord sur un fait / pour qu'il se passe quelque chose} : FR « d'accord \textbf{sur} / \textbf{pour que} » $\rightarrow$ ES \esp{de acuerdo \textbf{en} / \textbf{en que}} + \textbf{subjonctif} (volonté, projet).
\item \textbf{Opposition} : FR « s'opposer \textbf{à} » $\rightarrow$ ES \esp{oponerse \textbf{a}}. La préposition \esp{a} est \textbf{obligatoire} même devant un infinitif ou une proposition (\esp{a que} + subjonctif). Ne l'omettez pas (*\esp{me opongo prohibir} est faux).
\end{itemize}
\end{aretenir}

\begin{center}
\begin{tabularx}{\linewidth}{|X|X|X|}
\hline
\rowcolor{popblueL} \textbf{Français} & \textbf{Español} & \textbf{Remarque} \\
\hline
Je suis d'accord avec toi. & \esp{Estoy de acuerdo contigo.} & \esp{estar} + \esp{con} \\
\hline
Nous sommes d'accord sur ce point. & \esp{Estamos de acuerdo en ese punto.} & \esp{en} + nom \\
\hline
Je suis d'accord pour qu'il vienne. & \esp{Estoy de acuerdo en que venga.} & \esp{en que} + \textbf{subjonctif} \\
\hline
Je m'oppose à cette idée. & \esp{Me opongo a esa idea.} & verbe pronominal + \esp{a} \\
\hline
Je m'oppose à interdire le portable. & \esp{Me opongo a prohibir el móvil.} & \esp{a} + infinitif \\
\hline
Je m'oppose à ce qu'on l'interdise. & \esp{Me opongo a que lo prohíban.} & \esp{a que} + \textbf{subjonctif} \\
\hline
\end{tabularx}
\end{center}

\courssub{Contrastif français/espagnol et pièges de traduction}

\begin{attention}[Pièges fréquents pour francophones]
\begin{enumerate}
\item \textbf{*Soy de acuerdo} $\rightarrow$ \textbf{\esp{Estoy de acuerdo}}. (Voir boîte \emph{Attention} ci-dessus : état vs nature).
\item \textbf{*Estoy de acuerdo \emph{à} cette idée} $\rightarrow$ \textbf{\esp{Estoy de acuerdo \emph{con} esa idea}}. La préposition \esp{con} gouverne le nom/pronom. Pour une action, on utilise \esp{en que}.
\item \textbf{*Me opongo \emph{de} cette idée} $\rightarrow$ \textbf{\esp{Me opongo \emph{a} esa idea}}. Le verbe \esp{oponerse} régit \esp{a}, jamais \esp{de}.
\item \textbf{*Je m'y oppose} $\rightarrow$ \textbf{\esp{Me opongo a eso}}. En espagnol, la préposition \esp{a} ne tombe pas devant le pronom démonstratif \esp{eso} (contrairement au français « j'm'y oppose » où \emph{y} absorbe \emph{à}).
\item \textbf{*Tu as raison \emph{sur} ce point} $\rightarrow$ \textbf{\esp{Tienes razón \emph{en} ese punto}}. \esp{Tener razón} s'emploie avec \esp{en} + nom, ou seul.
\item \textbf{Faux ami : \esp{Realizar} $\neq$ Réaliser}. \esp{Realizar} = accomplir, effectuer, concrétiser. Pour « réaliser » (comprendre), on dit \esp{darse cuenta de}.
\end{enumerate}
\end{attention}

\begin{methode}[Réagir à une opinion en espagnol (étape par étape)]
\begin{enumerate}
\item \textbf{Identifier le degré} d'adhésion souhaité : total, partiel, refus poli, refus catégorique.
\item \textbf{Choisir la formule} adaptée (voir tableau p. 2).
\item \textbf{Construire la suite} :
      \begin{itemize}
      \item Accord avec personne : \esp{con} + pronom tonique (\esp{conmigo, contigo, con él...}).
      \item Accord sur action/projet : \esp{en que} + \textbf{subjonctif}.
      \item Opposition : \esp{oponerse a} + nom / infinitif / \esp{a que} + \textbf{subjonctif}.
      \end{itemize}
\item \textbf{Justifier} avec \esp{porque}, \esp{ya que}, \esp{puesto que} (cause) ou \esp{pero}, \esp{sin embargo}, \esp{no obstante} (concession/adversatif). Un désaccord nuancé sans justification sonne impoli.
\end{enumerate}
\end{methode}

\begin{exoresolu}[Traduction guidée (Thème)]
Énoncé : traduis en espagnol : « Je ne suis pas d'accord avec toi : au contraire, je m'oppose à cette interdiction, parce que le portable est un outil utile. »
\tcblower
Solution détaillée :
\begin{enumerate}
\item « Je ne suis pas d'accord avec toi » : verbe \esp{estar} (1ère pers. sg) + préposition \esp{con} + pronom tonique $\rightarrow$ \esp{No estoy de acuerdo contigo}.
\item « au contraire » : locution figée $\rightarrow$ \esp{al contrario}.
\item « je m'oppose à cette interdiction » : verbe pronominal \esp{oponerse} (conjugaison comme \esp{poner} : \esp{me opongo}) + préposition \esp{a} + nom $\rightarrow$ \esp{me opongo a esa prohibición}.
\item « parce que le portable est un outil utile » : \esp{porque el móvil es una herramienta útil}.
\end{enumerate}
Traduction complète : \esp{No estoy de acuerdo contigo: al contrario, me opongo a esa prohibición, porque el móvil es una herramienta útil.}
\end{exoresolu}

\begin{exoresolu}[Traduction guidée (Version)]
Énoncé : traduis en français : « \esp{Es cierto que los móviles distraen, pero no estoy de acuerdo en que se prohíban totalmente. Hasta cierto punto, tienes razón, aunque me opongo a una prohibición radical.} »
\tcblower
Solution détaillée :
\begin{enumerate}
\item \esp{Es cierto que} = « C'est vrai que » (concession). \esp{los móviles distraen} = « les portables distraient » (indicatif : fait constaté).
\item \esp{pero no estoy de acuerdo en que} = « mais je ne suis pas d'accord pour que / sur le fait que ». \esp{se prohíban} = \textbf{subjonctif présent} (3e pers. pl. de \esp{prohibir}) car \esp{estar de acuerdo en que} exprime une opinion/volonté sur une action future/hypothétique. $\rightarrow$ « qu'on les interdise ».
\item \esp{Hasta cierto punto, tienes razón} = « Jusqu'à un certain point, tu as raison » (accord partiel / atténuation).
\item \esp{although me opongo a una prohibición radical} = « bien que je m'oppose à une interdiction radicale » (\esp{aunque} + indicatif ici car le fait de s'opposer est réel/constaté ; \esp{me opongo a} + nom).
\end{enumerate}
Traduction : « C'est vrai que les portables distraient, mais je ne suis pas d'accord pour qu'on les interdise totalement. Jusqu'à un certain point, tu as raison, bien que je m'oppose à une interdiction radicale. »
\end{exoresolu}

\courssub{Entraînement à la traduction : thème et version}

\begin{exercice}[Thème : Français $\rightarrow$ Espagnol]
Traduis en espagnol en soignant les prépositions et les modes (indicatif/subjonctif).
\begin{enumerate}
\item Tu as raison en partie, mais je ne suis pas d'accord du tout.
\item — Est-ce que tu es d'accord avec cette règle ? — Pas question !
\item Nous nous opposons à une interdiction totale.
\item C'est vrai que c'est utile, mais je m'oppose à ce qu'on l'utilise en classe.
\item Je ne suis pas très convaincu, cependant j'accepte le débat.
\end{enumerate}
\end{exercice}

\begin{exercice}[Version : Espagnol $\rightarrow$ Français]
Traduis en français naturel.
\begin{enumerate}
\item \esp{— ¿Estás de acuerdo con la propuesta? — Por supuesto, es una idea excelente.}
\item \esp{Me opongo a que cambien el horario: no estoy de acuerdo en que perjudique a los alumnos.}
\item \esp{Admito que el móvil distrae, pero puede ser una herramienta muy útil si se controla.}
\item \esp{¡De ninguna manera! Ni mucho menos voy a aceptar esa norma.}
\item \esp{No lo veo muy claro, aunque hasta cierto punto tienes razón.}
\end{enumerate}
\end{exercice}

\begin{corrige}[Exercices : Thème et Version]
\textbf{Thème :}
\begin{enumerate}
\item \esp{Tienes razón en parte, pero no estoy de acuerdo del todo.}
\item \esp{— ¿Estás de acuerdo con esta regla? — ¡De ninguna manera!}
\item \esp{Nos oponemos a una prohibición total.}
\item \esp{Es cierto que es útil, pero me opongo a que se use en clase.} (\esp{se use} = subjonctif passif/impersonnel).
\item \esp{No lo veo muy claro, sin embargo acepto el debate.} (ou \esp{no estoy muy convencido}).
\end{enumerate}

\textbf{Version :}
\begin{enumerate}
\item « — Es-tu d'accord avec la proposition ? — Bien sûr, c'est une excellente idée. »
\item « Je m'oppose à ce qu'ils changent l'horaire : je ne suis pas d'accord pour que cela nuise aux élèves. »
\item « J'admets que le portable distrait, mais il peut être un outil très utile s'il est contrôlé. »
\item « Pas question ! Loin de là je vais accepter cette règle. » (ou « Je n'accepterai pas cette règle. »)
\item « Je ne le vois pas très clair, bien que jusqu'à un certain point tu as raison. »
\end{enumerate}
\end{corrige}

\begin{experience}[Débat oral guidé : « Le portable au lycée »]
\textbf{Consigne :} En binômes, rejouez le dialogue de la leçon en inversant les rôles, puis improvisez une suite : le proviseur arrive et demande votre avis final. Utilisez \textbf{au moins 3 formules d'accord}, \textbf{3 de désaccord} et \textbf{1 concession} (\esp{es cierto que... pero}). Le professeur notera la justesse des prépositions (\esp{con, en, a}) et l'emploi du subjonctif après \esp{en que / a que}.
\end{experience}

\begin{savaistu}[Le débat, une tradition hispanique]
Dans les lycées d'Espagne et d'Amérique latine, le \esp{debate escolar} est une activité très répandue, souvent encadrée par des ligues nationales (ex. \esp{Liga Española de Debate Universitario} adaptée au secondaire). Les élèves apprennent à défendre une opinion avec des formules codifiées : \esp{a mi juicio} (à mon avis), \esp{estoy a favor de} (je suis pour), \esp{me opongo firmemente a} (je m'oppose fermement à). En Guinée équatoriale, seul pays africain hispanophone, ces formules font partie du langage quotidien des élèves de \esp{bachillerato} et de la vie parlementaire. Les maîtriser, c'est aussi se rapprocher de vos voisins hispanophones du continent et réussir l'épreuve orale du Baccalauréat !
\end{savaistu}

\begin{textebox}[Observation : débat radiophonique « Téléphones en classe »]
\esp{— Presentador: Buenos días, amigos. Hoy debatimos: ¿teléfono móvil en el aula, sí o no? Con nosotros, Kevin desde Yaoundé, Amina desde Douala y Carlos desde Madrid. Kevin, ¿qué opinas?} \\
\esp{— Kevin: Buenas. \textbf{Estoy totalmente de acuerdo} con que el móvil es una herramienta pedagógica. \textbf{Por supuesto}, hay riesgos, pero \textbf{efectivamente} permite acceder a diccionarios, apps educativas...} \\
\esp{— Amina: \textbf{No estoy de acuerdo en parte}. \textbf{Tienes razón} en que es útil, \textbf{pero} en la práctica los chicos chatean en WhatsApp. \textbf{Me parece} que distrae más que ayuda.} \\
\esp{— Carlos: \textbf{Hasta cierto punto} os doy la razón a los dos. En mi instituto usamos \emph{Chromebooks}, no móviles. \textbf{Comparto la opinión} de Amina: el móvil es personal, cuesta controlarlo.} \\
\esp{— Presentador: \textbf{Exacto}, Carlos. \textbf{Desde luego}, la gestión del aula es clave. Kevin, ¿qué propones?} \\
\esp{— Kevin: \textbf{Estoy de acuerdo en} regular su uso. \textbf{Estamos de acuerdo en que} el profesor debe decidir. \textbf{¿Estás de acuerdo conmigo}, Amina?} \\
\esp{— Amina: \textbf{¡Claro que sí!} \textbf{De acuerdo}, pero con normas claras. \textbf{Sin duda}, eso funcionaría mejor.}
\par\medskip
\textbf{Traduction :} — Animateur : Bonjour les amis. Aujourd'hui on débat : téléphone portable en classe, oui ou non ? Avec nous, Kevin depuis Yaoundé, Amina depuis Douala et Carlos depuis Madrid. Kevin, ton avis ? — Kevin : Salut. \textbf{Je suis totalement d'accord} pour dire que le mobile est un outil pédagogique. \textbf{Bien sûr}, il y a des risques, mais \textbf{effectivement} il permet d'accéder à des dictionnaires, apps éducatives... — Amina : \textbf{Je ne suis pas tout à fait d'accord}. \textbf{Tu as raison} sur l'utilité, \textbf{mais} en pratique les gamins tchattent sur WhatsApp. \textbf{Il me semble} que ça distrait plus qu'autre chose. — Carlos : \textbf{Jusqu'à un certain point} je vous donne raison à tous les deux. Dans mon lycée on utilise des \emph{Chromebooks}, pas des mobiles. \textbf{Je partage l'avis} d'Amina : le mobile est personnel, difficile à contrôler. — Animateur : \textbf{Exact}, Carlos. \textbf{Bien entendu}, la gestion de classe est clé. Kevin, tu proposes quoi ? — Kevin : \textbf{Je suis d'accord pour} réguler son usage. \textbf{Nous sommes d'accord que} le prof doit décider. \textbf{Es-tu d'accord avec moi}, Amina ? — Amina : \textbf{Bien sûr que oui !} \textbf{D'accord}, mais avec des règles claires. \textbf{Sans aucun doute}, ça marcherait mieux.
\end{textebox}

\coursec{Les formules d'accord : formes et emplois}

Dans la section précédente, nous avons observé un débat radiophonique où deux lycéens de Yaoundé et leur correspondant espagnol échangeaient leurs points de vue sur l'usage du téléphone portable en classe. Nous y avons repéré des expressions comme \esp{estoy de acuerdo}, \esp{tienes razón} ou \esp{por supuesto}. Il est temps maintenant d'étudier ces formules une à une : comment se construisent-elles ? Comment se conjuguent-elles ? Quelles sont leurs nuances de registre ? C'est toute la grammaire de l'accord que nous allons systématiser ici, avant de passer, dans la section suivante, aux formules de désaccord.

\courssub{Estar de acuerdo : la formule de base}

\begin{definition}[Estar de acuerdo]
\esp{Estar de acuerdo} signifie « être d'accord ». C'est la formule d'accord la plus courante et la plus neutre en espagnol. Le mot \esp{acuerdo} est ici un nom invariable : on ne dit jamais \esp{*estoy de acuerda}, même si l'on est une femme.
\end{definition}

Observons d'abord la conjugaison complète au présent de l'indicatif :

\begin{center}
\begin{tabular}{|l|l|l|}
\hline
\rowcolor{popblueL} \textbf{Personne} & \textbf{Español} & \textbf{Français} \\
\hline
yo & estoy de acuerdo & je suis d'accord \\
\hline
tú & estás de acuerdo & tu es d'accord \\
\hline
él / ella / usted & está de acuerdo & il / elle / vous (formel) êtes d'accord \\
\hline
nosotros / nosotras & estamos de acuerdo & nous sommes d'accord \\
\hline
vosotros / vosotras & estáis de acuerdo & vous êtes d'accord (Espagne) \\
\hline
ellos / ellas / ustedes & están de acuerdo & ils / elles / vous (formel pl.) êtes d'accord \\
\hline
\end{tabular}
\end{center}

\begin{attention}[Le verbe, pas l'adjectif]
En français, « d'accord » ressemble à un adjectif, et le francophone est tenté de l'accorder. En espagnol, c'est le verbe \esp{estar} qui porte toute la flexion ; \esp{de acuerdo} ne bouge jamais : \esp{María está de acuerdo} « Marie est d'accord », \esp{las chicas están de acuerdo} « les filles sont d'accord ». Jamais \esp{*de acuerda}, jamais \esp{*de acuerdos}.
\end{attention}

\begin{savaistu}[Vosotros vs Ustedes au Cameroun]
Au Cameroun, comme dans toute l'Amérique hispanophone et en Guinée équatoriale, le pronom \esp{vosotros} (et sa conjugaison \esp{estáis}) est très peu utilisé à l'oral courant ; on lui préfère \esp{ustedes} (\esp{están}) pour le « vous » pluriel, quel que soit le registre. Pour le bac, il faut cependant connaître la forme \esp{vosotros} (programme officiel), mais en expression orale ou écrite libre, \esp{ustedes} est la norme standard attendue.
\end{savaistu}

\courssub{Les constructions de estar de acuerdo}

C'est ici que les francophones commettent le plus d'erreurs, car le français dit « être d'accord \emph{avec} quelqu'un », « être d'accord \emph{que}... », « être d'accord \emph{pour}... ». L'espagnol a ses propres prépositions, qu'il faut mémoriser.

\begin{propriete}[Les trois constructions de estar de acuerdo]
\begin{enumerate}
\item \cle{estar de acuerdo con + personne ou chose} : \esp{Estoy de acuerdo con tu hermano} « Je suis d'accord avec ton frère » ; \esp{Estoy de acuerdo con la decisión} « Je suis d'accord avec la décision ».
\item \cle{estar de acuerdo en que + indicatif} : \esp{Estoy de acuerdo en que hay que actuar} « Je suis d'accord qu'il faut agir ». Après \esp{en que}, le verbe est à l'\textbf{indicatif} car on constate un accord sur un fait présenté comme réel.
\item \cle{estar de acuerdo en + infinitif} : \esp{Estamos de acuerdo en organizar una fiesta} « Nous sommes d'accord pour organiser une fête ». Cette construction exprime un accord sur une action à faire.
\end{enumerate}
\end{propriete}

\begin{exemplebox}[Exemples traduits]
\esp{Estamos de acuerdo en que la educación es prioritaria.} « Nous sommes d'accord que l'éducation est prioritaire. »

\esp{¿Estás de acuerdo conmigo?} « Es-tu d'accord avec moi ? »

\esp{Los alumnos están de acuerdo en hacer una excursión a Kribi.} « Les élèves sont d'accord pour faire une excursion à Kribi. »

\esp{No estoy de acuerdo con esa ley.} « Je ne suis pas d'accord avec cette loi. »
\end{exemplebox}

\begin{attention}[Piège de la préposition]
Ne traduisez jamais « avec lui » par \esp{*con lo} : \esp{lo} ne renvoie pas à une personne. « Je suis d'accord avec lui » se dit \esp{estoy de acuerdo con él}. De même, attention à \esp{conmigo} « avec moi » et \esp{contigo} « avec toi », formes contractées obligatoires : jamais \esp{*con mí}, jamais \esp{*con ti}.
\end{attention}

\courssub{Tener razón : « avoir raison »}

\begin{definition}[Tener razón]
\esp{Tener razón} signifie « avoir raison ». Le nom \esp{razón} est invariable : quelle que soit la personne qui parle, on dit toujours \esp{razón} au singulier. On peut le renforcer : \esp{tener toda la razón} « avoir entièrement raison ». La variante \esp{llevar razón} est également correcte et très fréquente en Espagne péninsulaire.
\end{definition}

\begin{exemplebox}[Exemples traduits]
\esp{Tienes razón, el examen es mañana.} « Tu as raison, l'examen est demain. »

\esp{Mi profesor tiene toda la razón.} « Mon professeur a entièrement raison. »

\esp{No tienes razón, Douala no es la capital.} « Tu as tort, Douala n'est pas la capitale. »

\esp{Estás equivocado, el partido es el sábado.} « Tu te trompes, le match est samedi. »
\end{exemplebox}

\begin{attention}[Faux amis et constructions interdites]
\begin{itemize}
\item « Avoir raison » se dit \esp{tener razón}, jamais \esp{*tener razón de...} suivi d'un infinitif. Pour « tu as raison de protester », dites \esp{haces bien en protestar} ou \esp{tienes razón al protestar}.
\item « Tu as tort » se dit \esp{no tienes razón} ou \esp{estás equivocado}. Ne dites jamais \esp{*tienes tort} : le mot \esp{tort} n'existe pas en espagnol.
\item \esp{Razón} est féminin : \esp{toda la razón}, pas \esp{*todo el razón}.
\end{itemize}
\end{attention}

\courssub{Les formules figées d'approbation}

Au-delà des verbes, l'espagnol dispose d'une série de formules figées, très fréquentes à l'oral, qui valident immédiatement la parole de l'interlocuteur. Les voici, avec leurs équivalents français :

\begin{popfigure}
\begin{tikzpicture}[
  box/.style={draw=popblue, thick, fill=popblueL, rounded corners=2pt, align=center, text width=5.0cm, minimum height=0.85cm, font=\small},
  eq/.style={draw=poporange, thick, fill=poporangeL, rounded corners=2pt, align=center, text width=5.0cm, minimum height=0.85cm, font=\small},
  node distance=0.2cm and 1.0cm]
\node[box, fill=popblue, text=white, font=\small\bfseries] (h1) {Formule};
\node[eq, fill=poporange, text=white, font=\small\bfseries, right=of h1] (h2) {Équivalent français};
\node[box, below=of h1] (a1) {por supuesto};
\node[eq, right=of a1] {bien sûr, évidemment};
\node[box, below=of a1] (a2) {efectivamente};
\node[eq, right=of a2] {effectivement, en effet};
\node[box, below=of a2] (a3) {claro (que sí)};
\node[eq, right=of a3] {bien sûr (que oui)};
\node[box, below=of a3] (a4) {desde luego};
\node[eq, right=of a4] {bien entendu, évidemment};
\node[box, below=of a4] (a5) {exacto};
\node[eq, right=of a5] {exactement, c'est exact};
\node[box, below=of a5] (a6) {¡eso es!};
\node[eq, right=of a6] {c'est ça !, voilà !};
\node[box, below=of a6] (a7) {sin duda};
\node[eq, right=of a7] {sans aucun doute};
\node[box, below=of a7] (a8) {tienes razón};
\node[eq, right=of a8] {tu as raison};
\end{tikzpicture}
\legende{Les huit formules d'accord et leurs équivalents français}
\end{popfigure}

\begin{exemplebox}[Exemples traduits]
\esp{—¿Vendrás a la reunión? —Por supuesto que vendré.} « — Tu viendras à la réunion ? — Bien sûr que je viendrai. »

\esp{—El cacao es la riqueza de esta región. —Efectivamente.} « — Le cacao est la richesse de cette région. — Effectivement. »

\esp{—¿Te gusta el fútbol? —¡Claro que sí!} « — Tu aimes le football ? — Bien sûr que oui ! »

\esp{—Hay que proteger el medio ambiente. —Desde luego.} « — Il faut protéger l'environnement. — Bien entendu. »

\esp{—Entonces el tren sale a las ocho. —Exacto.} « — Donc le train part à huit heures. — Exact. »
\end{exemplebox}

\begin{attention}[Faux ami : efectivamente]
\esp{Efectivamente} ne signifie presque jamais « effectivement » au sens français de « en réalité, contrairement à ce qu'on croit ». Dans un débat, c'est un marqueur d'\textbf{accord} : « en effet, c'est bien vrai ». De même, \esp{actualmente} signifie « actuellement » et non « actuellement = en fait ». Méfiez-vous des adverbes en \esp{-mente} qui ressemblent au français.
\end{attention}

\begin{savaistu}[Efectivamente, le mot des débats télévisés]
Dans le monde hispanophone, \esp{efectivamente} est omniprésent dans les débats télévisés et les interviews politiques : l'orateur l'utilise pour valider l'argument de son interlocuteur avant de répondre ou de nuancer. C'est un véritable marqueur de \textbf{politesse débattive} : dire \esp{efectivamente} montre qu'on écoute l'autre, même si l'on va ensuite exprimer un désaccord. En Guinée équatoriale, seul pays hispanophone d'Afrique subsaharienne, on l'entend aussi bien dans les débats de Malabo que dans les conversations quotidiennes.
\end{savaistu}

\courssub{Accord renforcé et accord partiel}

Toutes les opinions ne se valent pas : parfois on approuve à cent pour cent, parfois seulement à moitié. L'espagnol possède des outils précis pour ces deux nuances.

\begin{propriete}[L'accord renforcé]
Pour intensifier l'accord, on place un adverbe devant \esp{de acuerdo} :
\begin{itemize}
\item \esp{estoy totalmente de acuerdo} « je suis totalement d'accord » ;
\item \esp{estoy completamente de acuerdo} « je suis complètement d'accord » ;
\item \esp{no puedo estar más de acuerdo} « je ne peux qu'être d'accord » (littéralement « je ne peux pas être plus d'accord »), formule très expressive.
\end{itemize}
\end{propriete}

\begin{propriete}[L'accord partiel]
Pour limiter l'accord, on utilise :
\begin{itemize}
\item \esp{estoy de acuerdo en parte} « je suis d'accord en partie » ;
\item \esp{hasta cierto punto} « jusqu'à un certain point » ;
\item \esp{puede ser, pero...} « peut-être, mais... », formule de concession qui prépare une objection.
\end{itemize}
\end{propriete}

\begin{exemplebox}[Exemples traduits]
\esp{No puedo estar más de acuerdo contigo: el deporte une a los pueblos.} « Je ne peux qu'être d'accord avec toi : le sport unit les peuples. »

\esp{Estoy de acuerdo en parte, pero el problema es más complejo.} « Je suis d'accord en partie, mais le problème est plus complexe. »

\esp{Hasta cierto punto tienes razón.} « Jusqu'à un certain point, tu as raison. »

\esp{Puede ser, pero hay que pensar en los jóvenes.} « Peut-être, mais il faut penser aux jeunes. »
\end{exemplebox}

\courssub{Les registres de langue : du soutenu au familier}

Comme en français, on ne parle pas de la même manière à son ami et au proviseur. Voici les principaux niveaux de registre pour exprimer l'accord :

\begin{center}
\begin{tabularx}{\linewidth}{|X|X|X|}
\hline
\rowcolor{popblueL} \textbf{Registre} & \textbf{Español} & \textbf{Français} \\
\hline
Soutenu & Comparto su opinión. & Je partage votre avis. \\
\hline
Soutenu & Coincido con usted. & Je suis de votre avis / Je coïncide avec vous. \\
\hline
Neutre & Estoy de acuerdo contigo. & Je suis d'accord avec toi. \\
\hline
Neutre & Tienes razón. & Tu as raison. \\
\hline
Familier & ¡Claro! / ¡Eso es! & Bien sûr ! / C'est ça ! \\
\hline
Familier & ¡Exacto! & Carrément ! / Exact ! \\
\hline
\end{tabularx}
\end{center}

\begin{exemplebox}[Exemples traduits]
\esp{Comparto tu opinión sobre la lectura.} « Je partage ton avis sur la lecture. »

\esp{—El barrio necesita una biblioteca. —¡Eso es!} « — Le quartier a besoin d'une bibliothèque. — C'est ça ! »

\esp{Coincido con usted, señor director.} « Je suis de votre avis, monsieur le directeur. »
\end{exemplebox}

\begin{methode}[Traduire « je suis tout à fait d'accord avec lui »]
\begin{enumerate}
\item Identifier la formule de base : \esp{estar de acuerdo}.
\item Conjuguer \esp{estar} à la première personne : \esp{estoy}.
\item Traduire l'intensité « tout à fait » par \esp{totalmente} ou \esp{completamente}, placé avant \esp{de acuerdo}.
\item Traduire « avec lui » par la préposition \esp{con} + pronom tonique \esp{él}.
\item Résultat : \esp{Estoy totalmente de acuerdo con él.} — et non \esp{*con lo}, qui ne désigne pas une personne.
\end{enumerate}
\end{methode}

\begin{exoresolu}[Thème : phrases d'opinion]
Traduisez en espagnol : « — Tu as entièrement raison : l'éducation est prioritaire. — Bien sûr, et je suis d'accord avec toi pour organiser un débat au lycée. »
\tcblower
\textbf{Solution détaillée.}

Première réplique : « tu as entièrement raison » se traduit par \esp{tienes toda la razón} (renforcement de \esp{tener razón} par \esp{toda la}, féminin singulier). « L'éducation est prioritaire » : \esp{la educación es prioritaria}.

Deuxième réplique : « bien sûr » = \esp{por supuesto} ou \esp{claro}. « Je suis d'accord avec toi pour organiser » : accord avec une personne (\esp{contigo}) + accord sur une action à faire, donc \esp{estar de acuerdo en + infinitif} : \esp{estoy de acuerdo contigo en organizar}.

\textbf{Traduction complète :} \esp{—Tienes toda la razón: la educación es prioritaria. —Por supuesto, y estoy de acuerdo contigo en organizar un debate en el liceo.}
\end{exoresolu}

\begin{exercice}[À vous de jouer]
Traduisez en espagnol :
\begin{enumerate}
\item Nous sommes d'accord que le sport est important pour la santé.
\item — Tu viens au match samedi ? — Bien sûr que oui !
\item Je partage votre avis, madame la professeure.
\item Je ne peux qu'être d'accord avec elle : il faut protéger la forêt.
\item Tu as raison jusqu'à un certain point, mais réfléchis encore.
\end{enumerate}
\end{exercice}

\begin{corrige}[Exercice « À vous de jouer »]
\begin{enumerate}
\item \esp{Estamos de acuerdo en que el deporte es importante para la salud.} (construction \esp{en que} + indicatif)
\item \esp{—¿Vienes al partido el sábado? —¡Claro que sí!}
\item \esp{Comparto su opinión, señora profesora.} (registre soutenu, \esp{su} et non \esp{tu})
\item \esp{No puedo estar más de acuerdo con ella: hay que proteger el bosque.} (\esp{con ella}, personne ; jamais \esp{*con la})
\item \esp{Hasta cierto punto tienes razón, pero reflexiona un poco más.} (ou \esp{piénsalo mejor})
\end{enumerate}
\end{corrige}

\begin{aretenir}[L'essentiel des formules d'accord]
\begin{itemize}
\item \esp{Estar de acuerdo} : seul le verbe \esp{estar} se conjugue ; \esp{de acuerdo} est invariable.
\item Trois constructions : \esp{con} + personne, \esp{en que} + indicatif, \esp{en} + infinitif.
\item \esp{Tener razón} (invariable), renforcé en \esp{tener toda la razón} ; « avoir tort » = \esp{no tener razón} / \esp{estar equivocado}. Variante \esp{llevar razón}.
\item Formules figées : \esp{por supuesto}, \esp{efectivamente}, \esp{claro que sí}, \esp{desde luego}, \esp{exacto}, \esp{¡eso es!}, \esp{sin duda}.
\item Intensité : \esp{totalmente de acuerdo}, \esp{no puedo estar más de acuerdo} ; nuance : \esp{en parte}, \esp{hasta cierto punto}, \esp{puede ser, pero...}.
\item Registre soutenu : \esp{comparto su opinión}, \esp{coincido con usted} ; registre familier : \esp{¡claro!}, \esp{¡eso es!}.
\end{itemize}
\end{aretenir}

\coursec{Les formules de désaccord : formes et emplois}

Dans la section précédente, nous avons appris à dire \emph{oui} : \esp{estoy de acuerdo}, \esp{tienes razón}, \esp{por supuesto}. Mais la vie citoyenne, le débat en classe, la discussion entre amis exigent aussi de savoir dire \emph{non} — et de le dire avec la bonne intensité : parfois poliment, parfois fermement, parfois catégoriquement. C'est l'objet de cette section : maîtriser toute la gamme du désaccord en espagnol, de la nuance discrète à l'opposition radicale, afin de pouvoir traduire sans faute des phrases comme « Je m'oppose à cette décision » ou « Pas question ! ».

\courssub{Observation : le désaccord en contexte}

Lisons d'abord un dialogue entre deux élèves d'un lycée de Yaoundé qui préparent un débat sur les réseaux sociaux. Soulignons mentalement toutes les formules qui expriment le désaccord.

\begin{textebox}[Dialogue : un débat au lycée]
— Yo creo que las redes sociales son muy útiles para los jóvenes. ¿Estás de acuerdo?

— \textbf{No estoy muy de acuerdo}. \textbf{Entiendo tu punto de vista, pero} pienso que también tienen peligros: la adicción, las noticias falsas...

— ¡Pero si gracias a Internet podemos aprender gratis!

— \textbf{Al contrario}, muchos alumnos pierden el tiempo en vez de estudiar. \textbf{Me opongo a} que los niños usen el teléfono antes de los quince años.

— ¿\textbf{De ninguna manera} quieres reconocer las ventajas?

— \textbf{No digo eso}. Reconozco las ventajas, pero \textbf{estoy en contra de} un uso sin control. \textbf{Rechazo totalmente} la idea de que todo \textbf{sea} positivo.

— Bueno, veo que \textbf{te opones firmemente} a mi propuesta...

— ¡\textbf{Qué va}! Solo quiero que hablemos de los dos lados del problema.
\end{textebox}

\textbf{Traduction.} — Je pense que les réseaux sociaux sont très utiles pour les jeunes. Tu es d'accord ? — Je ne suis pas très d'accord. Je comprends ton point de vue, mais je pense qu'ils ont aussi des dangers : la dépendance, les fausses informations... — Mais grâce à Internet, on peut apprendre gratuitement ! — Au contraire, beaucoup d'élèves perdent leur temps au lieu d'étudier. Je m'oppose à ce que les enfants utilisent le téléphone avant quinze ans. — Tu refuses absolument de reconnaître les avantages ? — Je ne dis pas cela. Je reconnais les avantages, mais je suis contre un usage sans contrôle. Je rejette totalement l'idée que tout soit positif. — Bon, je vois que tu t'opposes fermement à ma proposition... — Mais non ! Je veux seulement que nous parlions des deux côtés du problème.

\textbf{Questions d'observation :}
\begin{enumerate}
\item Relevez toutes les formules de désaccord du dialogue et classez-les de la plus polie à la plus catégorique.
\item Quelle formule le second élève utilise-t-il pour corriger une exagération de son camarade ?
\item Quels verbes expriment l'opposition ? Sont-ils suivis d'une préposition ?
\end{enumerate}

\courssub{La négation simple : \esp{no estar de acuerdo} et \esp{no tener razón}}

La forme la plus simple du désaccord consiste à nier les formules d'accord vues dans la section précédente. Il suffit de placer \esp{no} devant le verbe.

\begin{propriete}[Négation simple de l'accord]
\begin{itemize}
\item \esp{no estar de acuerdo} + \esp{con} + nom/personne : \esp{No estoy de acuerdo con esa decisión.} « Je ne suis pas d'accord avec cette décision. »
\item \esp{no estar de acuerdo} + \esp{en que} + subordonnée : \esp{No estoy de acuerdo en que haya un problema.} « Je ne suis pas d'accord sur le fait qu'il y ait un problème. »
\item \esp{no tener razón} : \esp{Lo siento, pero no tienes razón.} « Je suis désolé, mais tu n'as pas raison. »
\end{itemize}
\end{propriete}

\begin{exemplebox}[Exemples traduits]
\esp{No estoy de acuerdo contigo en este punto.} « Je ne suis pas d'accord avec toi sur ce point. »

\esp{El profesor no está de acuerdo con la nueva medida.} « Le professeur n'est pas d'accord avec la nouvelle mesure. »

\esp{Perdona, pero no tienes razón: el examen es el lunes.} « Pardon, mais tu n'as pas raison : l'examen est lundi. »
\end{exemplebox}

\begin{attention}[La négation espagnole reste simple]
En français on dit « je ne suis \emph{pas} d'accord » avec deux mots de négation ; en espagnol, un seul \esp{no} suffit : \esp{no estoy de acuerdo}. Ne calquez pas le « pas » français sur \esp{nada} : \esp{no estoy de acuerdo} suffit pour « je ne suis pas d'accord ». \esp{No estoy nada de acuerdo} signifie « je ne suis pas \emph{du tout} d'accord » (\esp{nada} renforce la négation).
\end{attention}

\courssub{Le désaccord catégorique : \esp{de ninguna manera}, \esp{en absoluto}, \esp{¡qué va!}}

Quand on veut refuser net, sans appel, l'espagnol dispose d'une série de locutions très fréquentes à l'oral comme à l'écrit.

\begin{definition}[Locutions de refus catégorique]
\begin{itemize}
\item \esp{de ninguna manera} / \esp{de ningún modo} = en aucune façon, pas question
\item \esp{en absoluto} = pas du tout, absolument pas
\item \esp{ni mucho menos} = loin de là, certainement pas
\item \esp{para nada} = pas du tout (très courant en Amérique latine)
\item \esp{¡qué va!} = mais non ! voyons ! (réaction spontanée, Espagne)
\end{itemize}
\end{definition}

\begin{exemplebox}[Exemples traduits]
\esp{¡De ninguna manera aceptaré eso!} « Je n'accepterai jamais cela ! »

\esp{¿Piensas renunciar? — De ningún modo.} « Tu penses abandonner ? — Pas question. »

\esp{¿Te molesta que abra la ventana? — En absoluto.} « Ça te dérange que j'ouvre la fenêtre ? — Pas du tout. »

\esp{No es tonto, ni mucho menos.} « Il n'est pas bête, loin de là. »

\esp{¿Vas a copiar en el examen? — ¡Qué va!} « Tu vas copier à l'examen ? — Mais non ! »
\end{exemplebox}

\begin{attention}[Faux ami : \esp{en absoluto} ne veut pas dire « absolument »]
C'est l'un des pièges les plus célèbres pour les francophones. \esp{En absoluto} s'emploie \textbf{uniquement dans un contexte négatif ou interrogatif} et signifie « pas du tout ». Pour dire « absolument » au sens affirmatif, on utilise \esp{absolutamente}, \esp{totalmente} ou \esp{desde luego}.

\esp{¿Te gusta el fútbol? — ¡Absolutamente!} « Tu aimes le football ? — Absolument ! » (et non \esp{en absoluto}, qui voudrait dire le contraire !)
\end{attention}

\courssub{L'opposition frontale : \esp{al contrario}, \esp{por el contrario}}

Ces locutions servent à opposer une idée à celle qui vient d'être dite, ou à corriger une affirmation fausse.

\begin{propriete}[Les trois « contraires »]
\begin{itemize}
\item \esp{al contrario} : corrige directement ce qui vient d'être dit. \esp{Al contrario, la situación mejora.} « Au contraire, la situation s'améliore. »
\item \esp{todo lo contrario} : renforcement. \esp{No empeora; todo lo contrario.} « Ça n'empire pas ; bien au contraire. »
\item \esp{por el contrario} : introduit une opposition entre deux idées (plus soutenu, très utile dans un essai). \esp{Unos prefieren la ciudad; por el contrario, yo prefiero el campo.} « Certains préfèrent la ville ; en revanche, je préfère la campagne. »
\end{itemize}
\end{propriete}

\begin{attention}[\esp{al contrario} et \esp{por el contrario} ne sont pas interchangeables]
\esp{Al contrario} réfute une affirmation (« tu dis que ça va mal ? au contraire, ça va bien ! »), tandis que \esp{por el contrario} oppose deux réalités (« lui travaille ; en revanche, elle se repose »). En version, « au contraire » de réfutation se traduit \esp{al contrario} ; « en revanche » se traduit \esp{por el contrario} ou \esp{en cambio}.
\end{attention}

\courssub{Les verbes d'opposition : \esp{oponerse a}, \esp{estar en contra de}, \esp{estar a favor de}}

Pour exprimer une opposition construite, argumentée — celle du débat citoyen et du commentaire de texte —, l'espagnol utilise des verbes et des locutions verbales qu'il faut conjuguer correctement.

\begin{propriete}[Le verbe \esp{oponerse a} : forme et constructions]
\esp{Oponerse} est un verbe \textbf{toujours pronominal} et \textbf{toujours suivi de la préposition \esp{a}}. Il admet trois constructions :
\begin{itemize}
\item \esp{oponerse a} + nom : \esp{Me opongo a la violencia.} « Je m'oppose à la violence. »
\item \esp{oponerse a} + infinitif : \esp{Me opongo a cerrar la biblioteca.} « Je m'oppose à fermer la bibliothèque. »
\item \esp{oponerse a que} + \textbf{subjonctif} : \esp{Me opongo a que vengas.} « Je m'oppose à ce que tu viennes. »
\end{itemize}
Antonymes utiles : \esp{estar a favor de} « être favorable à », \esp{apoyar} « soutenir ».
\end{propriete}

\begin{attention}[\esp{oponerse} est pronominal et exige \esp{a}]
Deux erreurs typiques des francophones : oublier le pronom réfléchi (\esp{opongo la idea}) ou oublier la préposition (\esp{me opongo la idea}). La seule forme correcte est \esp{me opongo a la idea}. Le français « je m'oppose cette idée » (sans « à ») ne doit pas être calqué.
\end{attention}

\begin{popfigure}
\begin{center}
\begin{tabular}{|l|l|}
\hline
\rowcolor{popblueL} \textbf{oponerse (présent)} & \textbf{Traduction} \\
\hline
yo \textbf{me opongo} & je m'oppose \\
\hline
tú \textbf{te opones} & tu t'opposes \\
\hline
él / ella / usted \textbf{se opone} & il / elle s'oppose \\
\hline
nosotros \textbf{nos oponemos} & nous nous opposons \\
\hline
vosotros \textbf{os oponéis} & vous vous opposez \\
\hline
ellos / ustedes \textbf{se oponen} & ils / elles s'opposent \\
\hline
\end{tabular}
\end{center}
\legende{Conjugaison du verbe pronominal \esp{oponerse} au présent de l'indicatif : irrégularité \esp{opongo} à la première personne (comme \esp{poner} : \esp{pongo}).}
\end{popfigure}

\begin{exemplebox}[Exemples traduits]
\esp{Estoy en contra de la pena de muerte.} « Je suis contre la peine de mort. »

\esp{Muchos ciudadanos se oponen a la nueva ley.} « Beaucoup de citoyens s'opposent à la nouvelle loi. »

\esp{Nos oponemos a que cierren el mercado de Mokolo.} « Nous nous opposons à ce qu'on ferme le marché de Mokolo. »

\esp{¿Estás a favor de los deberes el fin de semana?} « Tu es favorable aux devoirs le week-end ? »
\end{exemplebox}

\courssub{Le désaccord poli et l'intensité}

Dans un débat, surtout avec un professeur ou un adulte, le désaccord brutal est mal vu. L'espagnol, comme le français, adoucit l'opposition par des formules de politesse, ou au contraire la renforce par des adverbes d'intensité.

\begin{definition}[Atténuer ou renforcer le désaccord]
\textbf{Désaccord poli (atténuation) :}
\begin{itemize}
\item \esp{No estoy muy de acuerdo.} « Je ne suis pas très d'accord. »
\item \esp{Lo siento, pero no comparto tu opinión.} « Je suis désolé, mais je ne partage pas ton opinion. »
\item \esp{Entiendo tu punto de vista, pero...} « Je comprends ton point de vue, mais... »
\end{itemize}
\textbf{Désaccord intense (renforcement) :}
\begin{itemize}
\item \esp{Me opongo firmemente / radicalmente a esa medida.} « Je m'oppose fermement / radicalement à cette mesure. »
\item \esp{Rechazo totalmente esa idea.} « Je rejette totalement cette idée. »
\item \esp{Estoy completamente en desacuerdo.} « Je suis en complet désaccord. »
\end{itemize}
\end{definition}

\begin{popfigure}
\begin{center}
\begin{tikzpicture}[scale=1]
\draw[-{Stealth}, very thick, popblue] (0,0) -- (12,0);
\node[above, align=center, fill=popgreenL, rounded corners, inner sep=4pt, font=\small] at (1.5,0.3) {no estoy muy\\de acuerdo};
\node[above, align=center, fill=popgoldL, rounded corners, inner sep=4pt, font=\small] at (5,0.3) {no estoy\\de acuerdo};
\node[above, align=center, fill=poporangeL, rounded corners, inner sep=4pt, font=\small] at (8.5,0.3) {de ninguna\\manera};
\node[above, align=center, fill=poppinkL, rounded corners, inner sep=4pt, font=\small] at (11,0.3) {me opongo\\radicalmente};
\node[below, font=\small\itshape, text=popmuted] at (1.5,-0.2) {poli, nuancé};
\node[below, font=\small\itshape, text=popmuted] at (11,-0.2) {catégorique};
\end{tikzpicture}
\end{center}
\legende{L'échelle du désaccord : de la nuance polie au refus radical. Bien choisir sa formule selon l'interlocuteur et le contexte.}
\end{popfigure}

\courssub{Tableau récapitulatif}

\begin{center}
\begin{tabularx}{\linewidth}{|X|X|X|}
\hline
\rowcolor{popblueL} \textbf{Français} & \textbf{Español} & \textbf{Remarque} \\
\hline
Je ne suis pas d'accord (avec) & No estoy de acuerdo (con) & négation simple \\
\hline
Tu n'as pas raison & No tienes razón & accent obligatoire sur \esp{razón} (oxytone en \esp{-n}) \\
\hline
Pas question / jamais & De ninguna manera / de ningún modo & refus catégorique \\
\hline
Pas du tout & En absoluto / para nada & \esp{en absoluto} = faux ami (« pas du tout ») \\
\hline
Loin de là & Ni mucho menos & renforcement négatif \\
\hline
Mais non ! & ¡Qué va! & oral, spontané (Espagne) \\
\hline
Au contraire & Al contrario / todo lo contrario & réfutation directe \\
\hline
En revanche & Por el contrario / en cambio & opposition d'idées \\
\hline
Je m'oppose à (+ nom / inf.) & Me opongo a (+ nom / inf.) & pronominal + \esp{a} \\
\hline
Je m'oppose à ce que tu viennes & Me opongo a que vengas & \esp{a que} + subjonctif \\
\hline
Je suis contre / pour & Estoy en contra de / a favor de & locutions verbales \\
\hline
Je ne partage pas ton avis & No comparto tu opinión & désaccord poli \\
\hline
Je rejette totalement cette idée & Rechazo totalmente esa idea & intensité maximale \\
\hline
\end{tabularx}
\end{center}

\courssub{Exercice résolu et entraînement}

\begin{exoresolu}[Thème : traduisez en espagnol]
Énoncé. Traduisez : « Je m'oppose fermement à cette décision, et je ne suis pas d'accord pour que les élèves sortent plus tôt. — Pas question ! répond le directeur. Au contraire, la situation va s'améliorer. »
\tcblower
Solution détaillée. \esp{Me opongo firmemente a esta decisión y no estoy de acuerdo en que los alumnos salgan más temprano. — ¡De ninguna manera! — responde el director. Al contrario, la situación va a mejorar.}

Points de méthode : « je m'oppose à » exige le pronominal \esp{me opongo a} ; « pour que » + subordonnée de désaccord se rend par \esp{en que} + subjonctif (\esp{salgan}) ; « pas question » = \esp{de ninguna manera} ; « au contraire » (réfutation) = \esp{al contrario} et non \esp{por el contrario}.
\end{exoresolu}

\begin{exercice}[À vous de jouer]
Traduisez en espagnol :
\begin{enumerate}
\item Je ne suis pas très d'accord avec toi : tu n'as pas raison sur ce point.
\item Tu vas vendre ton vélo ? — Pas du tout !
\item Nous nous opposons à ce qu'on détruise ce vieux bâtiment.
\item Je suis contre la violence, mais je suis favorable au dialogue.
\item Je comprends ton point de vue, mais je ne partage pas ton opinion.
\end{enumerate}
\end{exercice}

\begin{corrige}[Exercice « À vous de jouer »]
\begin{enumerate}
\item \esp{No estoy muy de acuerdo contigo: no tienes razón en este punto.}
\item \esp{¿Vas a vender tu bicicleta? — ¡En absoluto! / ¡Para nada!}
\item \esp{Nos oponemos a que destruyan este edificio viejo.} (subjonctif après \esp{a que})
\item \esp{Estoy en contra de la violencia, pero estoy a favor del diálogo.}
\item \esp{Entiendo tu punto de vista, pero no comparto tu opinión.}
\end{enumerate}
\end{corrige}

\begin{savaistu}[Variantes régionales : ¡Qué va! (Espagne) et ¡Para nada! (Amérique latine)]
\esp{¡Qué va!} est la formule de désaccord la plus spontanée en Espagne : on l'entend partout, dans les marchés, les cafés, les cours de récréation. En Amérique latine — et en Guinée équatoriale, le seul pays hispanophone d'Afrique — on entendra plutôt \esp{¡para nada!} ou \esp{¡de ningún modo!}. Connaître ces variantes régionales est un atout pour la compréhension orale : le sens reste le même, seule la couleur régionale change.
\end{savaistu}

\begin{aretenir}[L'essentiel de la section]
\begin{itemize}
\item Négation simple : \esp{no estar de acuerdo (con / en que)}, \esp{no tener razón}.
\item Refus catégorique : \esp{de ninguna manera}, \esp{en absoluto} (faux ami : « pas du tout »), \esp{ni mucho menos}, \esp{¡qué va!}.
\item Opposition : \esp{al contrario} (réfutation) / \esp{por el contrario} (opposition d'idées).
\item \esp{Oponerse a} + nom/infinitif ; \esp{oponerse a que} + subjonctif ; conjugaison : \esp{me opongo, te opones, se opone, nos oponemos, os oponéis, se oponen}.
\item Politesse : \esp{no estoy muy de acuerdo}, \esp{no comparto tu opinión} ; intensité : \esp{me opongo firmemente a}, \esp{rechazo totalmente}.
\end{itemize}
\end{aretenir}

\coursec{Modalisation : atténuer et concéder}

Dans la section précédente, nous avons appris à dire \esp{no estoy de acuerdo}, \esp{de ninguna manera} ou \esp{me opongo}. Mais dans la vie citoyenne comme à l'examen, un désaccord exprimé de façon brute peut sembler agressif, et un accord total peut paraître naïf. Le bon locuteur sait \cle{nuancer} : dire « oui, mais… », « peut-être », « jusqu'à un certain point ». C'est exactement l'objet de la \cle{modalisation}, compétence centrale du Módulo 3 et outil indispensable du commentaire de texte et de la traduction en Terminale A.

\courssub{Qu'est-ce que la modalisation ?}

\begin{definition}[La modalisation]
La \cle{modalisation} est l'ensemble des procédés linguistiques qui permettent d'exprimer le \cle{degré de certitude} ou d'\cle{adhésion} du locuteur à son propos : certitude, doute, probabilité, réserve, concession. En espagnol comme en français, on module une opinion grâce à des verbes (\esp{creo que}, \esp{me parece que}), des adverbes (\esp{quizás}, \esp{posiblemente}) et des tournures concessives (\esp{aunque}, \esp{es verdad que… pero}).
\end{definition}

Observons d'abord ces procédés dans un dialogue authentique entre deux élèves de Terminale à Yaoundé.

\begin{textebox}[Diálogo : ¿Deberían los alumnos usar el móvil en clase?]
— Yo creo que el móvil debería estar prohibido en el aula. Es una distracción constante.

— Hasta cierto punto tienes razón, aunque no comparto tu conclusión. Me parece que el problema no es el móvil en sí, sino el uso que hacemos de él.

— ¿De verdad? A mí me parece que los alumnos no saben controlarse.

— Quizás tengas razón en algunos casos. Puede que muchos compañeros lo usen para chatear, es verdad. Pero, aun reconociendo que existe ese riesgo, creo que el móvil también puede ser una herramienta educativa: se puede consultar un diccionario, buscar información, grabar una explicación del profesor…

— Es cierto que es útil para buscar palabras; sin embargo, la mayoría pierde el tiempo en las redes sociales.

— De acuerdo, pero entonces la solución no es prohibir, sino educar. Aunque sea difícil, los profesores deberían enseñarnos un uso responsable.

— Bueno, entiendo lo que dices, pero sigo pensando que, de alguna manera, la prohibición es más eficaz.

— Por mucho que insistas, no me convencerás: ¡prohibir no es educar!
\end{textebox}

\textbf{Glossaire :} \esp{el aula} = la salle de classe ; \esp{sino} = mais (après négation) ; \esp{controlarse} = se contrôler ; \esp{chatear} = tchater ; \esp{la herramienta} = l'outil ; \esp{grabar} = enregistrer ; \esp{sin embargo} = cependant ; \esp{prohibir} = interdire ; \esp{convencer} = convaincre.

\textbf{Questions de compréhension :}
\begin{enumerate}
\item Relevez trois procédés qui \emph{atténuent} une opinion et trois procédés de \emph{concession}.
\item Quelle est la thèse du premier locuteur ? Et celle du second ?
\item Traduisez : \esp{Por mucho que insistas, no me convencerás.}
\end{enumerate}

\courssub{Atténuer avec les verbes d'opinion}

Pour ne pas affirmer de façon catégorique, l'espagnol utilise, comme le français, des verbes d'opinion à la première personne : \esp{creo que} (je crois que), \esp{pienso que} (je pense que), \esp{me parece que} (il me semble que), \esp{opino que} (j'estime que), \esp{diría que} (je dirais que). Suivis de \esp{que} + \cle{indicatif}, ils présentent l'idée comme une opinion personnelle, donc discutable.

\begin{exemplebox}[Verbes d'opinion à l'affirmatif]
\esp{Creo que el deporte une a los pueblos.} « Je crois que le sport unit les peuples. »

\esp{Me parece que la propuesta es interesante.} « Il me semble que la proposition est intéressante. »

\esp{Diría que es una buena idea, pero hay que estudiarla.} « Je dirais que c'est une bonne idée, mais il faut l'étudier. » (le conditionnel \esp{diría} atténue encore davantage)
\end{exemplebox}

\begin{propriete}[Indicatif ou subjonctif après les verbes d'opinion]
Règle fondamentale :
\begin{itemize}
\item Après \esp{creo que}, \esp{me parece que}, \esp{pienso que}, \esp{opino que} \textbf{affirmatifs} $\rightarrow$ \cle{indicatif} : \esp{Creo que \textbf{tiene} razón.} « Je crois qu'il a raison. »
\item Après \esp{no creo que}, \esp{no me parece que}, \esp{no pienso que}, \esp{no opino que} \textbf{négatifs} $\rightarrow$ \cle{subjonctif} : \esp{No creo que \textbf{tenga} razón.} « Je ne crois pas qu'il ait raison. »
\end{itemize}
La négation introduit le doute, et le doute exige le subjonctif en espagnol.
\end{propriete}

\begin{attention}[Piège du francophone]
En français, on dit « je ne crois pas qu'il \emph{ait} raison » avec un subjonctif qui passe souvent inaperçu. En espagnol, l'oubli du subjonctif est une faute grave et très visible : ne dites jamais \esp{no creo que tiene razón}. Dites \esp{no creo que \textbf{tenga} razón}. De même : \esp{No me parece que sea justo.} « Il ne me semble pas que ce soit juste. »
\end{attention}

\courssub{Atténuer avec les adverbes et \esp{puede que} + subjonctif}

Les adverbes de doute permettent de présenter un fait comme simplement possible :

\begin{center}
\begin{tabularx}{\linewidth}{|X|X|X|}
\hline
\rowcolor{popblueL} Español & Français & Remarque \\
\hline
\esp{quizás / quizá} & peut-être & + subjonctif (doute) / indicatif (probabilité forte) \\
\hline
\esp{tal vez} & peut-être & + subjonctif généralement \\
\hline
\esp{posiblemente} & possiblement / peut-être & registre soutenu \\
\hline
\esp{en cierto modo} & d'une certaine manière & atténue un jugement \\
\hline
\esp{de alguna manera} & d'une certaine façon & idem \\
\hline
\esp{más o menos} & plus ou moins & atténue une qualité \\
\hline
\esp{hasta cierto punto} & jusqu'à un certain point & accord partiel \\
\hline
\end{tabularx}
\end{center}

\begin{exemplebox}[Adverbes d'atténuation]
\esp{Quizás tengas razón.} « Peut-être as-tu raison. » (subjonctif : doute)

\esp{En cierto modo, estoy de acuerdo contigo.} « D'une certaine manière, je suis d'accord avec toi. »

\esp{La película es más o menos interesante.} « Le film est plus ou moins intéressant. »
\end{exemplebox}

\begin{propriete}[\esp{Puede que} + subjonctif]
La tournure \esp{puede que} + \cle{subjonctif} traduit « il se peut que » :
\esp{Puede que tengas razón.} « Il se peut que tu aies raison. »
\esp{Puede que esté equivocado.} « Il se peut qu'il se trompe. »
Le subjonctif est \textbf{obligatoire}, car on exprime une simple possibilité. Ne confondez pas avec \esp{puede ser que}, équivalent synonyme : \esp{Puede ser que llueva mañana.} « Il se peut qu'il pleuve demain. »
\end{propriete}

\courssub{La concession avec \esp{aunque} : indicatif ou subjonctif}

La concession consiste à \textbf{reconnaître un argument adverse} avant de défendre sa propre position : c'est l'arme rhétorique du débat citoyen et du commentaire de texte. La conjonction reine est \esp{aunque} (bien que, quoique, même si).

\begin{attention}[\esp{Aunque} + indicatif ou + subjonctif : le choix qui change le sens]
\begin{itemize}
\item \esp{aunque} + \cle{indicatif} = fait \textbf{réel, constaté, admis par les deux locuteurs} : \esp{Aunque llueve, salgo.} « Bien qu'il pleuve (c'est un fait, je le vois), je sors. »
\item \esp{aunque} + \cle{subjonctif} = fait \textbf{hypothétique, incertain, futur, non réalisé} : \esp{Aunque llueva, saldré.} « Même s'il pleut (hypothèse pour demain), je sortirai. »
\end{itemize}
Le français utilise le subjonctif dans les deux cas ; l'espagnol, lui, distingue la réalité de l'hypothèse par le mode. C'est l'un des points les plus rentables — et les plus sanctionnés — en thème.
\end{attention}

\begin{exemplebox}[Les deux emplois d'\esp{aunque}]
\esp{Aunque tienes razón, no comparto tu conclusión.} « Bien que tu aies raison (fait admis), je ne partage pas ta conclusion. » — indicatif \esp{tienes}.

\esp{Aunque sea difícil, lo intentaré.} « Même si c'est difficile (hypothèse), je l'essaierai. » — subjonctif \esp{sea}.

\esp{Aunque no me escuches, seguiré hablando.} « Même si tu ne m'écoutes pas (hypothèse), je continuerai à parler. » — subjonctif.
\end{exemplebox}

\begin{propriete}[\esp{Por mucho que} + subjonctif = « avoir beau »]
La tournure \esp{por mucho que} + subjonctif traduit le français « avoir beau + infinitif » ou « pour autant que » :
\esp{Por mucho que insistas, no cambiaré de opinión.} « Tu as beau insister, je ne changerai pas d'avis. »
\esp{Por mucho que estudie, no aprueba.} « Il a beau étudier, il ne réussit pas. »
Variantes : \esp{por más que}, \esp{por muchos que} (accord en genre et en nombre selon le nom qui suit \esp{que} : \esp{por muchos errores que cometa}).
\end{propriete}

\courssub{Les tournures concessives de la discussion}

Au-delà d'\esp{aunque}, l'espagnol dispose d'un riche répertoire pour concéder avant de répliquer. Ces tournures sont l'ossature d'un débat poli et d'un essai équilibré.

\begin{center}
\begin{tabularx}{\linewidth}{|X|X|}
\hline
\rowcolor{popgreenL} Español & Français \\
\hline
\esp{Es verdad que…, pero…} & Il est vrai que…, mais… \\
\hline
\esp{Es cierto que…, sin embargo…} & Il est certain que…, cependant… \\
\hline
\esp{Tienes razón, pero…} & Tu as raison, mais… \\
\hline
\esp{Entiendo lo que dices, pero…} & Je comprends ce que tu dis, mais… \\
\hline
\esp{Aun reconociendo que…, creo que…} & Tout en reconnaissant que…, je crois que… \\
\hline
\esp{Hasta cierto punto tienes razón, aunque…} & Jusqu'à un certain point tu as raison, bien que… \\
\hline
\esp{Estoy de acuerdo en que…, pero…} & Je suis d'accord sur le fait que…, mais… \\
\hline
\end{tabularx}
\end{center}

\begin{exemplebox}[Concession en contexte]
\esp{Es cierto que cuesta mucho, pero vale la pena.} « Il est vrai que ça coûte cher, mais ça en vaut la peine. »

\esp{Entiendo lo que dices, pero no lo comparto.} « Je comprends ce que tu dis, mais je ne le partage pas. »

\esp{Aun reconociendo que el proyecto es caro, pienso que es necesario.} « Tout en reconnaissant que le projet est coûteux, je pense qu'il est nécessaire. »

\esp{Hasta cierto punto tienes razón, aunque exageras un poco.} « Jusqu'à un certain point tu as raison, bien que tu exagères un peu. »
\end{exemplebox}

\courssub{Contrastif français/espagnol et pièges de traduction (synthèse)}

Ce tableau récapitule les correspondances essentielles et les erreurs à ne pas commettre pour l'ensemble de la leçon (Accord, Désaccord, Modalisation).

\begin{center}
\begin{tabularx}{\linewidth}{|X|X|X|}
\hline
\rowcolor{poppurpleL} Français & Español correct & \textbf{Piège à éviter} \\
\hline
Je suis d'accord & \esp{Estoy de acuerdo} & \esp{*Soy de acuerdo} (confusion \esp{ser}/\esp{estar}) \\
\hline
Tu as raison & \esp{Tienes razón} & \esp{*Eres razón} (verbe \esp{tener}, pas \esp{ser}) \\
\hline
Je ne suis pas d'accord & \esp{No estoy de acuerdo} & \esp{*No soy de acuerdo} \\
\hline
Absolument pas / De aucune façon & \esp{De ninguna manera} / \esp{En absoluto} & \esp{*De ningún modo} (possible mais moins idiomatique pour un refus catégorique) \\
\hline
Au contraire & \esp{Al contrario} & \esp{*En el contrario} \\
\hline
Je m'y oppose & \esp{Me opongo} & \esp{*Me opongo a eso} (redondant si contexte clair) \\
\hline
Je pense que c'est vrai & \esp{Creo que es verdad} \textbf{(Indicatif)} & \esp{*Creo que sea verdad} \\
\hline
Je ne pense pas que ce soit vrai & \esp{No creo que sea verdad} \textbf{(Subjonctif)} & \esp{*No creo que es verdad} \\
\hline
Bien qu'il soit fatigué (fait réel) & \esp{Aunque \textbf{está} cansado} \textbf{(Indicatif)} & \esp{*Aunque esté cansado} (si le fait est certain) \\
\hline
Même s'il est fatigué (hypothèse) & \esp{Aunque \textbf{esté} cansado} \textbf{(Subjonctif)} & \esp{*Aunque está cansado} (si c'est une hypothèse) \\
\hline
Tu as beau dire & \esp{Por mucho que digas} \textbf{(Subjonctif)} & \esp{*Por mucho que dices} \\
\hline
\end{tabularx}
\end{center}

\begin{methode}[Traduire « je suis d'accord, mais… »]
\begin{enumerate}
\item \textbf{Repérez} la structure française : accord + restriction (« mais », « cependant », « toutefois »).
\item \textbf{Choisissez} une tournure espagnole complète : \esp{Estoy de acuerdo, pero…} ou mieux, la concession \esp{Tienes razón, pero…} / \esp{Es verdad que…, pero…}.
\item \textbf{Vérifiez le mode} : verbe d'opinion affirmatif $\rightarrow$ indicatif ; négatif ou doute $\rightarrow$ subjonctif ; \esp{aunque} : réel $\rightarrow$ indicatif, hypothétique $\rightarrow$ subjonctif.
\item \textbf{Évitez} la juxtaposition sèche de deux phrases sans connecteur : un texte sans \esp{pero}, \esp{sin embargo} ou \esp{aunque} est jugé pauvre en expression écrite.
\end{enumerate}
\end{methode}

\begin{experience}[Débat citoyen : « La technologie en classe »]
\textbf{Objectif :} Pratiquer l'accord, le désaccord et la modalisation à l'oral.
\textbf{Déroulement :}
\begin{enumerate}
\item Former des binômes. L'un défend l'interdiction du téléphone (rôle A), l'autre son usage encadré (rôle B).
\item Utiliser \textbf{au moins} : 2 formules d'accord, 2 de désaccord, 1 atténuation (\esp{quizás}, \esp{me parece}), 1 concession (\esp{aunque} + indicatif ou subjonctif), 1 \esp{por mucho que}.
\item Exemple de début : \esp{A: « Creo que el móvil distrae. » B: « Hasta cierto punto tienes razón, aunque puede ser una herramienta. »}
\item Changement de rôles après 3 minutes. Le professeur note les structures utilisées au tableau pour correction collective.
\end{enumerate}
\end{experience}

\begin{popfigure}
\begin{tikzpicture}[
  mindmap, grow cyclic,
  every node/.style={concept, align=center, font=\small},
  root concept/.append style={concept color=popblue, font=\small\bfseries},
  level 1 concept/.append style={sibling angle=120, level distance=3.4cm},
  level 2 concept/.append style={level distance=2.6cm, font=\scriptsize},
]
\node[root concept]{Nuancer une opinion}
  child[concept color=poporange]{ node {ATENUAR}
    child[concept color=poporangeL]{ node[align=center]{me parece que\\ creo que\\ + indicativo} }
    child[concept color=poporangeL]{ node[align=center]{quizás, tal vez\\ puede que + subj.} }
  }
  child[concept color=popgreen]{ node {CONCEDER}
    child[concept color=popgreenL]{ node[align=center]{aunque + ind.\\ (hecho real)\\ aunque + subj.\\ (hipótesis)} }
    child[concept color=popgreenL]{ node[align=center]{es verdad que\ldots\\ pero\ldots\\ por mucho que\\ + subj.} }
  };
\end{tikzpicture}
\legende{Carte mentale : les deux grandes stratégies pour nuancer une opinion en espagnol.}
\end{popfigure}

\courssub{Exercice résolu et entraînement}

\begin{exoresolu}[Thème : nuancer un désaccord]
Traduisez en espagnol : « Je comprends ton point de vue, mais je ne crois pas que la solution soit si simple. Il se peut que tu aies raison sur un point : bien que le projet coûte cher, il est nécessaire. Tu as beau critiquer, le travail avance. »

\tcblower

\textbf{Solution détaillée :}

\esp{Entiendo tu punto de vista, pero no creo que la solución sea tan sencilla. Puede que tengas razón en un punto: aunque el proyecto cuesta caro, es necesario. Por mucho que critiques, el trabajo avanza.}

\textbf{Commentaires :}
\begin{itemize}
\item « Je comprends ton point de vue » $\rightarrow$ \esp{Entiendo tu punto de vista} : tournure concessive d'ouverture, suivie de \esp{pero}.
\item « je ne crois pas que la solution soit » $\rightarrow$ verbe d'opinion \textbf{négatif} : subjonctif obligatoire, \esp{no creo que… sea}.
\item « il se peut que tu aies raison » $\rightarrow$ \esp{puede que} + subjonctif : \esp{puede que tengas razón}.
\item « bien que le projet coûte cher » $\rightarrow$ fait \textbf{réel, constaté} : \esp{aunque} + indicatif \esp{cuesta}.
\item « tu as beau critiquer » $\rightarrow$ \esp{por mucho que} + subjonctif : \esp{por mucho que critiques}.
\end{itemize}
\end{exoresolu}

\begin{exercice}[À vous de jouer]
Traduisez en espagnol :
\begin{enumerate}
\item D'une certaine manière, tu as raison, mais je ne partage pas ta conclusion.
\item Même s'il pleut demain, nous irons au marché de Mokolo.
\item Il se peut qu'elle soit fatiguée ; elle travaille beaucoup.
\item Il est vrai que le français est utile, cependant l'espagnol ouvre aussi des portes.
\item Vous avez beau protester, la décision est prise.
\end{enumerate}
\end{exercice}

\begin{corrige}[Exercice « À vous de jouer »]
\begin{enumerate}
\item \esp{En cierto modo tienes razón, pero no comparto tu conclusión.}
\item \esp{Aunque llueva mañana, iremos al mercado de Mokolo.} (hypothèse $\rightarrow$ subjonctif)
\item \esp{Puede que esté cansada; trabaja mucho.}
\item \esp{Es verdad que el francés es útil; sin embargo, el español también abre puertas.}
\item \esp{Por mucho que protestéis, la decisión está tomada.} (Variante AmLat : \esp{Por mucho que protesten…})
\end{enumerate}
\end{corrige}

\begin{aretenir}[L'essentiel de la section]
\begin{itemize}
\item Modaliser, c'est doser son adhésion : \textbf{atténuer} (\esp{me parece que}, \esp{quizás}, \esp{puede que} + subjonctif) ou \textbf{concéder} (\esp{aunque}, \esp{es verdad que… pero}).
\item Opinion affirmative $\rightarrow$ indicatif ; opinion niée $\rightarrow$ subjonctif : \esp{no creo que tenga razón}.
\item \esp{Aunque} + indicatif = fait réel ; \esp{aunque} + subjonctif = hypothèse.
\item \esp{Por mucho que} + subjonctif = « avoir beau ».
\item En traduction comme au débat, jamais de juxtaposition sèche : reliez toujours par \esp{pero}, \esp{sin embargo}, \esp{aunque}.
\end{itemize}
\end{aretenir}

\coursec{Contrastif français/espagnol et pièges de traduction}

Dans la section précédente, nous avons appris à \cle{atténuer} et à \cle{concéder} pour nuancer une opinion (\esp{hasta cierto punto}, \esp{aunque}, \esp{es verdad que... pero}). Mais toute cette finesse est perdue si la construction grammaticale est fautive. Or, c'est précisément ici que le candidat francophone trébuche : il connaît le vocabulaire, mais il traduit \cle{mot à mot} et produit \esp{*soy de acuerdo} ou \esp{*tengo razón de}. Cette section est donc consacrée à un travail systématique de \cle{comparaison} entre le français et l'espagnol : nous allons démonter chaque piège, un par un, pour que la construction correcte devienne un réflexe.

\courssub{Tableau contrastif des constructions de base}

Commençons par observer les quatre constructions fondamentales de l'accord et du désaccord, côte à côte dans les deux langues.

\begin{center}
\begin{tabularx}{\linewidth}{|X|X|X|}
\hline
\rowcolor{popblueL} \textbf{Français} & \textbf{Español} & \textbf{Remarque} \\
\hline
être d'accord avec qqn & \esp{estar de acuerdo con algn} & verbe \esp{estar}, préposition \esp{con} \\
\hline
être d'accord sur le fait que & \esp{estar de acuerdo en que} & préposition \esp{en} + \esp{que} \\
\hline
s'opposer à qqch / à ce que & \esp{oponerse a algo / a que} & verbe pronominal + \esp{a} \\
\hline
être contre qqch & \esp{estar en contra de algo} & locution figée avec \esp{de} \\
\hline
être pour qqch & \esp{estar a favor de algo} & locution figée avec \esp{de} \\
\hline
avoir raison / tort & \esp{tener razón / no tener razón} & verbe \esp{tener}, sans article \\
\hline
\end{tabularx}
\end{center}

\begin{exemplebox}[Observation en contexte]
\esp{Estoy de acuerdo con María en que el proyecto es demasiado caro.} « Je suis d'accord avec María sur le fait que le projet est trop cher. »

\esp{Los vecinos se oponen a la construcción de la fábrica.} « Les voisins s'opposent à la construction de l'usine. »

\esp{Estamos a favor de la paz, pero en contra de esa ley.} « Nous sommes pour la paix, mais contre cette loi. »

\esp{Tú tienes razón: el transporte en Yaoundé es un problema serio.} « Tu as raison : le transport à Yaoundé est un problème sérieux. »
\end{exemplebox}

On remarque immédiatement que le français et l'espagnol ne choisissent \textbf{ni les mêmes verbes} (\emph{être}/\emph{avoir} contre \esp{estar}/\esp{tener}) \textbf{ni toujours les mêmes prépositions}. C'est de ce double décalage que naissent toutes les erreurs.

\begin{propriete}[Deux verbes à ne jamais confondre : \esp{estar} et \esp{tener}]
\begin{itemize}
\item \textbf{Être d'accord} se traduit par \esp{ESTAR de acuerdo} : \esp{Estoy de acuerdo contigo.} « Je suis d'accord avec toi. » Le verbe \esp{ser} est ici \textbf{impossible} : \esp{*soy de acuerdo} est une faute grave.
\item \textbf{Avoir raison / tort} se traduit par \esp{TENER razón / no tener razón} : \esp{Tienes razón.} « Tu as raison. » ; \esp{No tengo razón.} « J'ai tort. » On ne dit jamais \esp{*soy razón} ni \esp{*estoy razón}.
\end{itemize}
Explication : pour l'espagnol, l'accord est un \textbf{état} (d'où \esp{estar}, verbe d'état), et la raison est quelque chose qu'on \textbf{possède} (d'où \esp{tener}, verbe de possession). Le français, lui, utilise \emph{être} et \emph{avoir} : la logique est différente, il faut l'accepter telle quelle.
\end{propriete}

\courssub{Le piège des prépositions}

Deuxième source d'erreurs : les prépositions. Le français dit \emph{d'accord \textbf{avec}}, \emph{s'opposer \textbf{à}}, \emph{être \textbf{contre}}, \emph{être \textbf{pour}}. L'espagnol répond : \esp{con}, \esp{a}, \esp{en contra de}, \esp{a favor de}. Observons chaque cas.

\begin{propriete}[Les quatre prépositions à mémoriser]
\begin{enumerate}
\item \esp{estar de acuerdo \textbf{con} alguien} : \esp{Estoy de acuerdo con el director.} « Je suis d'accord avec le directeur. »
\item \esp{estar de acuerdo \textbf{en que} + proposition} : \esp{Estamos de acuerdo en que hay que actuar.} « Nous sommes d'accord sur le fait qu'il faut agir. » Devant une subordonnée, \esp{con} devient \esp{en que}.
\item \esp{oponerse \textbf{a} algo / \textbf{a que} + subjonctif} : \esp{Me opongo a esa decisión.} « Je m'oppose à cette décision. » ; \esp{Se oponen a que se cierre el mercado.} « Ils s'opposent à ce que le marché soit fermé. »
\item \esp{estar \textbf{en contra de} / \textbf{a favor de} algo} : \esp{Estoy en contra de la guerra y a favor del diálogo.} « Je suis contre la guerre et pour le dialogue. »
\end{enumerate}
\end{propriete}

\begin{attention}[La préposition qui disparaît ou qui change]
\begin{itemize}
\item \esp{*Tengo razón de...} : faute par calque de « avoir raison \emph{de} ». En espagnol, \esp{tener razón} ne prend \textbf{aucune préposition} : \esp{Tienes razón, ese camino es peligroso.} « Tu as raison, ce chemin est dangereux. »
\item \esp{*Me opongo la propuesta} : faute par oubli de \esp{a}. Le verbe \esp{oponerse} exige \esp{a} : \esp{Me opongo a la propuesta.} « Je m'oppose à la proposition. »
\item \esp{*Estoy de acuerdo con que...} : devant une subordonnée, on dit \esp{en que}, pas \esp{con que} : \esp{Estoy de acuerdo en que el precio es alto.} « Je suis d'accord sur le fait que le prix est élevé. »
\item Ne pas confondre \esp{en contra de} (contre) avec \esp{contra} seul, qui existe aussi : \esp{un voto contra la ley} « un vote contre la loi ». Mais après \esp{estar}, la locution complète est \esp{en contra de}.
\end{itemize}
\end{attention}

\courssub{Le piège du subjonctif après les verbes d'opinion}

C'est le point le plus décisif pour le thème au baccalauréat. En français, « je ne pense pas qu'il \emph{a} raison » emploie l'indicatif dans le langage courant (le subjonctif soutenu « qu'il \emph{ait} raison » est souvent négligé). En espagnol, \textbf{aucune hésitation n'est permise} : le subjonctif est obligatoire.

\begin{propriete}[Subjonctif obligatoire après l'opinion niée ou interrogée]
Après les verbes d'opinion à la \textbf{forme négative} (\esp{no creer que}, \esp{no pensar que}, \esp{no parecer que}) ou \textbf{interrogative} (\esp{¿crees que...?}), ainsi qu'après \esp{oponerse a que}, \esp{estar en contra de que} et \esp{puede que}, la subordonnée est au \textbf{subjonctif} :
\begin{itemize}
\item \esp{No creo que tenga razón.} « Je ne crois pas qu'il ait raison. »
\item \esp{¿Crees que sea verdad?} « Tu crois que c'est vrai ? »
\item \esp{Me opongo a que se construya la fábrica.} « Je m'oppose à ce que l'usine soit construite. »
\item \esp{Puede que esté de acuerdo.} « Il se peut qu'elle soit d'accord. »
\end{itemize}
À l'inverse, à la \textbf{forme affirmative}, ces verbes sont suivis de l'\textbf{indicatif} : \esp{Creo que tiene razón.} « Je crois qu'il a raison. » ; \esp{Pienso que es verdad.} « Je pense que c'est vrai. »
\end{propriete}

\begin{exemplebox}[Les trois phrases modèles du Bac]
\esp{Puede que esté de acuerdo.} « Il se peut qu'elle soit d'accord. » (\esp{puede que} + subjonctif \esp{esté})

\esp{No pienso que esté equivocado.} « Je ne pense pas qu'il ait tort. » (\esp{no pensar que} + subjonctif \esp{esté} ; noter : « avoir tort » = \esp{estar equivocado} ou \esp{no tener razón})

\esp{Nos oponemos a que se vayan.} « Nous nous opposons à ce qu'ils partent. » (\esp{oponerse a que} + subjonctif \esp{se vayan}, verbe \esp{irse})
\end{exemplebox}

\begin{attention}[Erreur n°1 des candidats francophones]
Toutes les fautes suivantes proviennent d'une \textbf{traduction mot à mot du français}. Apprends à les repérer et à les corriger :
\begin{itemize}
\item \esp{*Soy de acuerdo.} → \esp{Estoy de acuerdo.} (calque de « je \emph{suis} d'accord » : verbe \esp{estar})
\item \esp{*Tengo razón de que venga.} → \esp{Tengo razón: tiene que venir.} (\esp{tener razón} ne se construit pas avec \esp{de})
\item \esp{*Me opongo la propuesta.} → \esp{Me opongo a la propuesta.} (préposition \esp{a} obligatoire)
\item \esp{*No creo que tiene razón.} → \esp{No creo que tenga razón.} (subjonctif \esp{tenga} après \esp{no creer que})
\item \esp{*Estoy de acuerdo con que es caro.} → \esp{Estoy de acuerdo en que es caro.} (\esp{en que} devant une subordonnée à l'indicatif)
\end{itemize}
\end{attention}

\courssub{Faux amis et mots pièges}

Le lexique de l'opinion contient plusieurs \cle{faux amis} redoutables, car ils ressemblent à des mots français mais n'ont pas le même sens.

\begin{definition}[Faux amis de l'opinion]
\begin{itemize}
\item \esp{en absoluto} = « pas du tout, absolument pas » (sens \textbf{négatif} uniquement) : \esp{¿Estás de acuerdo? — En absoluto.} « Tu es d'accord ? — Absolument pas. » Pour dire « absolument » au sens affirmatif, on utilise \esp{absolutamente}, \esp{desde luego} ou \esp{por supuesto}.
\item \esp{actualmente} = « actuellement, en ce moment » (sens \textbf{temporel}) : \esp{Actualmente vivo en Douala.} « J'habite actuellement à Douala. » Ce n'est jamais « en fait, en réalité », qui se traduit par \esp{en realidad} ou \esp{de hecho}.
\item \esp{sin duda} = « sans aucun doute, certainement » (sens \textbf{affirmatif fort}) : \esp{Sin duda, es el mejor jugador del equipo.} « C'est sans aucun doute le meilleur joueur de l'équipe. »
\end{itemize}
\end{definition}

\begin{attention}[« Sans doute » : le piège de l'ambiguïté française]
Le français « sans doute » a \textbf{deux sens} opposés selon le contexte, et l'espagnol les distingue nettement :
\begin{itemize}
\item Sens \textbf{certain} (« très probablement, certes ») → \esp{sin duda} : \esp{Sin duda vendrá mañana.} « Il viendra sans doute demain » (= il est quasiment sûr qu'il viendra).
\item Sens \textbf{douteux} (« peut-être ») → \esp{quizás / tal vez} (+ subjonctif) : \esp{Quizás venga mañana.} « Il viendra peut-être demain » (= rien n'est sûr).
\end{itemize}
Avant de traduire « sans doute », demande-toi toujours : la phrase exprime-t-elle la quasi-certitude ou la simple possibilité ?
\end{attention}

\courssub{L'ordre des mots dans les exclamations de désaccord}

En espagnol, pour renforcer un désaccord, on peut \textbf{placer la formule négative en tête de phrase}, ce qui provoque l'\cle{inversion} du verbe avant le sujet. Cette construction, très fréquente à l'oral comme à l'écrit, donne de l'emphase.

\begin{exemplebox}[L'inversion emphatique]
Ordre neutre : \esp{No acepto eso de ninguna manera.} « Je n'accepte absolument pas cela. »

Ordre emphatique : \esp{¡De ninguna manera acepto eso!} « En aucune façon je n'accepte cela ! »

Autres exemples :
\begin{itemize}
\item \esp{¡En absoluto estoy de acuerdo contigo!} « Je ne suis absolument pas d'accord avec toi ! »
\item \esp{¡Nunca he dicho tal cosa!} « Je n'ai jamais dit une telle chose ! »
\item \esp{¡Tampoco yo lo creo!} « Moi non plus je ne le crois pas ! »
\end{itemize}
Remarque : quand l'expression négative précède le verbe, le \esp{no} \textbf{disparaît} : \esp{De ninguna manera acepto} (et non \esp{*de ninguna manera no acepto}). Si l'expression suit le verbe, le \esp{no} est obligatoire : \esp{No acepto eso de ninguna manera}.
\end{exemplebox}

\begin{popfigure}
\begin{center}
\begin{tikzpicture}[
fr/.style={rectangle, rounded corners, draw=popblue, fill=popblueL, align=center, minimum height=1.1cm, text width=3.6cm, font=\small},
es/.style={rectangle, rounded corners, draw=popgreen, fill=popgreenL, align=center, minimum height=1.1cm, text width=4.2cm, font=\small},
er/.style={rectangle, rounded corners, draw=poppink, fill=poppinkL, align=center, minimum height=1.1cm, text width=4.2cm, font=\small},
hd/.style={font=\bfseries\small, text=popdark}
]
\node[hd] (h1) at (0,0) {Français};
\node[hd] (h2) at (5.4,0) {Espagnol correct};
\node[hd] (h3) at (10.8,0) {Erreur fréquente};
\node[fr] at (0,-1.6) {être d'accord avec qqn};
\node[es] at (5.4,-1.6) {\esp{estar de acuerdo con algn}};
\node[er] at (10.8,-1.6) {\esp{*ser de acuerdo a algn}};
\node[fr] at (0,-3.1) {avoir raison};
\node[es] at (5.4,-3.1) {\esp{tener razón}};
\node[er] at (10.8,-3.1) {\esp{*ser razón / *tener razón de}};
\node[fr] at (0,-4.6) {s'opposer à qqch};
\node[es] at (5.4,-4.6) {\esp{oponerse a algo}};
\node[er] at (10.8,-4.6) {\esp{*oponerse algo}};
\node[fr] at (0,-6.1) {je ne crois pas qu'il ait raison};
\node[es] at (5.4,-6.1) {\esp{no creo que tenga razón}};
\node[er] at (10.8,-6.1) {\esp{*no creo que tiene razón}};
\draw[-{Stealth}, thick, popgreen] (2.0,-1.6) -- (3.2,-1.6);
\draw[-{Stealth}, thick, popgreen] (2.0,-3.1) -- (3.2,-3.1);
\draw[-{Stealth}, thick, popgreen] (2.0,-4.6) -- (3.2,-4.6);
\draw[-{Stealth}, thick, popgreen] (2.0,-6.1) -- (3.2,-6.1);
\draw[-{Stealth}, thick, poppink, dashed] (7.7,-1.6) -- (8.6,-1.6);
\draw[-{Stealth}, thick, poppink, dashed] (7.7,-3.1) -- (8.6,-3.1);
\draw[-{Stealth}, thick, poppink, dashed] (7.7,-4.6) -- (8.6,-4.6);
\draw[-{Stealth}, thick, poppink, dashed] (7.7,-6.1) -- (8.6,-6.1);
\end{tikzpicture}
\end{center}
\legende{Tableau contrastif : la construction correcte (vert) et l'erreur de calque à bannir (rose, barrée par la flèche pointillée).}
\end{popfigure}

\courssub{Stratégie de traduction et entraînement}

\begin{methode}[Traduire une phrase d'opinion en 4 temps]
\begin{enumerate}
\item \textbf{Identifier} la formule d'accord ou de désaccord dans la phrase française (être d'accord, avoir tort, s'opposer, être contre, sans doute...).
\item \textbf{Choisir} l'équivalent espagnol exact : \esp{estar de acuerdo}, \esp{tener razón}, \esp{oponerse a}, \esp{estar en contra de}, \esp{sin duda} ou \esp{quizás}.
\item \textbf{Vérifier le couple verbe + préposition} : \esp{estar} (jamais \esp{ser}), \esp{tener} (jamais \esp{ser}), \esp{con} / \esp{a} / \esp{en que} / \esp{de}.
\item \textbf{Vérifier le mode de la subordonnée} : opinion niée, interrogative ou \esp{a que} / \esp{puede que} → subjonctif ; opinion affirmée → indicatif.
\end{enumerate}
\end{methode}

\begin{exoresolu}[Thème guidé : cinq phrases d'opinion]
Traduire en espagnol :
\begin{enumerate}
\item Je suis d'accord avec toi : ce marché est trop cher.
\item Tu crois qu'il a raison ?
\item Nous nous opposons à ce que le stade soit détruit.
\item Elle n'est absolument pas d'accord avec cette décision.
\item Il viendra sans doute (= certainement) à la réunion.
\end{enumerate}
\tcblower
\textbf{Solutions détaillées (méthode en 4 temps) :}
\begin{enumerate}
\item Formule : « être d'accord avec » → \esp{estar de acuerdo con}. Pas de subordonnée problématique. \esp{Estoy de acuerdo contigo: este mercado es demasiado caro.} (Noter la forme \esp{contigo} = « avec toi ».)
\item Opinion à l'interrogative → subjonctif : \esp{¿Crees que tenga razón?} (et non \esp{*tiene}).
\item \esp{Oponerse a que} + subjonctif : \esp{Nos oponemos a que el estadio sea destruido.} (subjonctif passif : \esp{sea destruido}.)
\item \esp{En absoluto} = « absolument pas » ; verbe \esp{estar} : \esp{Ella no está en absoluto de acuerdo con esta decisión.} ou, avec inversion emphatique : \esp{¡En absoluto está ella de acuerdo con esta decisión!}
\item « Sans doute » au sens certain → \esp{sin duda} : \esp{Sin duda vendrá a la reunión.} (Si le sens était « peut-être » : \esp{Quizás venga a la reunión.})
\end{enumerate}
\end{exoresolu}

\begin{exercice}[À toi de jouer]
Traduis en espagnol en appliquant la méthode en 4 temps :
\begin{enumerate}
\item Nous ne pensons pas que cette loi soit juste.
\item Il se peut qu'ils soient contre le projet.
\item Tu as raison : les bendskins sont dangereux la nuit.
\item Les élèves s'opposent à la suppression de la récréation.
\item En aucune façon je n'accepte cette injustice !
\end{enumerate}
\end{exercice}

\begin{corrige}[Exercice « À toi de jouer »]
\begin{enumerate}
\item \esp{No pensamos que esta ley sea justa.} (opinion niée → subjonctif \esp{sea})
\item \esp{Puede que estén en contra del proyecto.} (\esp{puede que} + subjonctif \esp{estén} ; \esp{del} = \esp{de} + \esp{el})
\item \esp{Tienes razón: los bendskins son peligrosos por la noche.} (\esp{tener razón}, sans préposition)
\item \esp{Los alumnos se oponen a la supresión del recreo.} (\esp{oponerse a} + nom)
\item \esp{¡De ninguna manera acepto esta injusticia!} (inversion emphatique, sans \esp{no})
\end{enumerate}
\end{corrige}

\begin{savaistu}[Où se gagnent les points au Bac ?]
Au baccalauréat, les phrases de thème portant sur l'opinion combinent presque toujours \textbf{deux difficultés à la fois} : une formule figée (\esp{estar de acuerdo con}, \esp{oponerse a}) \emph{plus} un subjonctif ou une préposition délicate. Les correcteurs le savent : les points se gagnent sur la \textbf{construction}, pas sur le vocabulaire. Un candidat qui écrit \esp{No creo que tenga razón} avec le subjonctif correct marque plus qu'un candidat qui utilise un vocabulaire riche mais écrit \esp{*soy de acuerdo}. En Guinée équatoriale, voisin hispanophone du Cameroun, ces formules rythment les débats publics quotidiens : \esp{estar a favor de}, \esp{oponerse a}, \esp{tener razón} y sont autant de marqueurs d'une expression citoyenne claire et correcte.
\end{savaistu}

\begin{aretenir}[L'essentiel de la section]
\begin{itemize}
\item Être d'accord = \esp{ESTAR de acuerdo CON algn / EN que} ; avoir raison = \esp{TENER razón} (sans préposition).
\item S'opposer à = \esp{oponerse A} ; être contre = \esp{estar EN CONTRA DE} ; être pour = \esp{estar A FAVOR DE}.
\item Opinion niée ou interrogée, \esp{oponerse a que}, \esp{puede que} → \textbf{subjonctif obligatoire}.
\item \esp{En absoluto} = absolument \emph{pas} ; \esp{sin duda} = certainement ; « peut-être » = \esp{quizás / tal vez}.
\item Inversion emphatique : \esp{¡De ninguna manera acepto eso!} (le \esp{no} disparaît devant le verbe).
\end{itemize}
\end{aretenir}

\coursec{Contrastif français/espagnol et pièges de traduction}

Dans les sections précédentes, nous avons observé les formules d'accord et de désaccord, puis les outils de la modalisation (concession, atténuation, doute). Il est maintenant indispensable de confronter les deux langues pour fixer les équivalences exactes et déjouer les pièges qui coûtent des points au baccalauréat : les faux amis (\esp{estar de acuerdo} ne s'écrit jamais *\esp{estar de acordo}, calque du portugais), les prépositions imposées par chaque verbe (\esp{oponerse a}, \esp{interesarse por}, \esp{contar con}), le choix du mode (subjonctif après \esp{no creo que}, \esp{puede que} ; indicatif après \esp{aunque} concessionnel, \esp{es verdad que}). Maîtriser ces contrastes, c'est sécuriser la note de thème et de version.

\begin{center}
\begin{tabularx}{\linewidth}{|X|X|X|}
\hline
\rowcolor{popblueL} \textbf{Formule espagnole} & \textbf{Rendu naturel en français} & \textbf{Piège à éviter} \\
\hline
\esp{Estoy de acuerdo (con)} & je suis d'accord (avec) & *\esp{estar de acordo} (portugais) \\
\hline
\esp{Tienes razón} & tu as raison & *\esp{ser razón} (calque \emph{être raison}) \\
\hline
\esp{Al contrario} & au contraire & ne pas confondre avec « contrairement à » = \esp{a diferencia de} \\
\hline
\esp{De ninguna manera} & en aucun cas, pas question & *\esp{de ningún modo} correct mais plus soutenu \\
\hline
\esp{¡Qué va !} & allons donc ! pas du tout ! & jamais de traduction littérale (*\esp{quoi va}) \\
\hline
\esp{Me opongo a} & je m'oppose à & verbe pronominal obligatoire (\esp{me/te/se...}) \\
\hline
\esp{Puede que + subj.} & il se peut que + subj. & subjonctif correctement conjugué (\esp{pueda, esté, tenga, sea, vaya}) \\
\hline
\esp{Es verdad que... pero...} & il est vrai que... mais... & concession = indicatif des deux côtés \\
\hline
\esp{Interesarse por} & s'intéresser à & *\esp{interesarse en} (anglicisme) \\
\hline
\esp{Votar por} & voter pour (quelqu'un) & *\esp{votar para} \\
\hline
\end{tabularx}
\end{center}

\begin{attention}[Puede que + subjonctif : la bonne nouvelle... et le piège]
Bonne nouvelle : dans « il se peut que + subjonctif », le français et l'espagnol coïncident parfaitement : \esp{puede que} + subjonctif. La structure est la même, le mode est le même. Le seul danger, c'est la \textbf{conjugaison} du subjonctif lui-même : \esp{pueda} (de \esp{poder}), \esp{esté} (de \esp{estar}), \esp{tenga} (de \esp{tener}), \esp{sea} (de \esp{ser}), \esp{vaya} (de \esp{ir}), \esp{haya} (de \esp{haber}). Une forme comme *\esp{puede que está} ou *\esp{puede que tiene} est une faute grave, sanctionnée au Bac. Révisez les six formes du subjonctif présent des verbes irréguliers avant l'examen.
\end{attention}

\coursec{Entraînement à la traduction : thème et version}

C'est en traduisant, phrase après phrase, texte après texte, que ces formules deviendront des réflexes. Cette section vous entraîne aux deux exercices du baccalauréat : le \cle{thème} (français $\rightarrow$ espagnol) et la \cle{version} (espagnol $\rightarrow$ français), avec des corrections commentées qui justifient chaque choix grammatical et lexical.

\courssub{La méthode de traduction en cinq étapes}

Traduire une phrase d'opinion ne consiste pas à remplacer mot à mot. C'est un parcours en cinq étapes, qu'il faut suivre dans l'ordre, au brouillon si nécessaire.

\begin{popfigure}
\begin{tikzpicture}[node distance=0.35cm, every node/.style={font=\small}]
\node[draw=popblue, fill=popblueL, rounded corners, minimum width=2.6cm, minimum height=0.9cm, align=center] (a) {\textbf{1. Lire}\\ la phrase entière};
\node[draw=poporange, fill=poporangeL, rounded corners, minimum width=2.6cm, minimum height=0.9cm, align=center, right=of a] (b) {\textbf{2. Identifier}\\ la formule d'opinion};
\node[draw=poppurple, fill=poppurpleL, rounded corners, minimum width=2.6cm, minimum height=0.9cm, align=center, right=of b] (c) {\textbf{3. Choisir}\\ l'équivalent espagnol};
\node[draw=popgreen, fill=popgreenL, rounded corners, minimum width=2.6cm, minimum height=0.9cm, align=center, right=of c] (d) {\textbf{4. Vérifier}\\ la construction};
\node[draw=poppink, fill=poppinkL, rounded corners, minimum width=2.6cm, minimum height=0.9cm, align=center, right=of d] (e) {\textbf{5. Relire}\\ accents et accords};
\draw[-{Stealth}, thick, popdark] (a) -- (b);
\draw[-{Stealth}, thick, popdark] (b) -- (c);
\draw[-{Stealth}, thick, popdark] (c) -- (d);
\draw[-{Stealth}, thick, popdark] (d) -- (e);
\end{tikzpicture}
\legende{Frise méthodologique : les cinq étapes d'une traduction réussie.}
\end{popfigure}

\begin{methode}[Appliquer la frise, étape par étape]
\begin{enumerate}
\item \textbf{Lire} la phrase française entière, sans commencer à écrire : quel est le sens global ? Qui parle ? À qui ? Quel registre ?
\item \textbf{Identifier} la formule d'opinion : est-ce un accord, un désaccord, une concession, un doute ? Quel degré d'intensité (total, nuancé, catégorique) ?
\item \textbf{Choisir} l'équivalent espagnol dans votre répertoire : \esp{estoy de acuerdo con}, \esp{de ninguna manera}, \esp{aunque} + mode, \esp{puede que} + subjonctif...
\item \textbf{Vérifier la construction} : la préposition exigée par le verbe (\esp{oponerse a}, \esp{estar de acuerdo con}, \esp{interesarse por}), le mode (subjonctif ou indicatif), la concordance des temps, l'accentuation, les accords (genre/nombre).
\item \textbf{Relire} la phrase espagnole à voix haute : sonne-t-elle naturellement ? Un hispanophone la dirait-il ainsi ?
\end{enumerate}
\end{methode}

\courssub{Thème guidé : six phrases d'opinion à traduire}

\begin{exercice}[Thème : la vie citoyenne en débat]
Traduisez en espagnol les phrases suivantes, en appliquant la méthode en cinq étapes.
\begin{enumerate}
\item Je suis tout à fait d'accord avec vous : l'éducation est la clé du développement.
\item Pas question que je vote pour lui !
\item Bien qu'il ait raison, je m'oppose à cette méthode.
\item Il se peut que les jeunes s'intéressent à la politique autrement.
\item Je ne crois pas que les réseaux sociaux soient dangereux pour la démocratie.
\item Tu as raison sur un point, mais je ne suis pas d'accord avec ta conclusion.
\end{enumerate}
\end{exercice}

\begin{exoresolu}[Correction commentée du thème, phrase par phrase]
Énoncé : voir les six phrases ci-dessus.
\tcblower
\textbf{Phrase 1.} « Je suis tout à fait d'accord avec vous : l'éducation est la clé du développement. »\\
\esp{Estoy totalmente de acuerdo con usted : la educación es la clave del desarrollo.}\\
Justification : « tout à fait » se rend par \esp{totalmente} (ou \esp{completamente}), placé devant \esp{de acuerdo}. « Avec vous » = \esp{con usted} : la préposition \esp{con} est obligatoire avec \esp{estar de acuerdo}. « La clé » = \esp{la clave} : attention au faux ami, \esp{clave} est féminin (\esp{la clave del desarrollo}).

\textbf{Phrase 2.} « Pas question que je vote pour lui ! »\\
\esp{¡De ninguna manera votaré por él !} ou \esp{¡Ni mucho menos !} ou \esp{¡Que yo vote por él, de ninguna manera !}\\
Justification : « pas question » est un désaccord catégorique ; l'équivalent naturel est \esp{de ninguna manera} ou \esp{ni mucho menos}. Le futur \esp{votaré} rend bien la fermeté du refus. « Pour lui » = \esp{por él} : c'est bien \esp{por} (vote en faveur de), jamais \esp{para}. Ne pas traduire littéralement *\esp{no hay cuestión}, qui ne veut rien dire ici.

\textbf{Phrase 3.} « Bien qu'il ait raison, je m'oppose à cette méthode. »\\
\esp{Aunque tiene razón, me opongo a ese método.}\\
Justification : « bien que » = \esp{aunque}. Le français exige le subjonctif (« ait ») ; l'espagnol, lui, emploie l'\textbf{indicatif} après \esp{aunque} quand le fait est présenté comme réel, admis : \esp{tiene razón}. C'est une concession effective, pas hypothétique. « S'opposer à » = \esp{oponerse a} : verbe pronominal obligatoire (\esp{me opongo}) et préposition \esp{a}.

\textbf{Phrase 4.} « Il se peut que les jeunes s'intéressent à la politique autrement. »\\
\esp{Puede que los jóvenes se interesen por la política de otra manera.}\\
Justification : « il se peut que » = \esp{puede que} + \textbf{subjonctif} : \esp{se interesen} (subjonctif présent de \esp{interesarse}, verbe à changement e $\rightarrow$ ie : \esp{interese, intereses, interese, interesemos, intereséis, interesen}). « S'intéresser à » = \esp{interesarse por} : préposition \esp{por}, pas \esp{en}.

\textbf{Phrase 5.} « Je ne crois pas que les réseaux sociaux soient dangereux pour la démocratie. »\\
\esp{No creo que las redes sociales sean peligrosas para la democracia.}\\
Justification : verbe d'opinion nié, \esp{no creo que}, donc subjonctif : \esp{sean} (subjonctif de \esp{ser}, irrégulier : \esp{sea, seas, sea, seamos, seáis, sean}). « Dangereux pour » = \esp{peligrosas para} : \esp{para} marque ici la destination, l'objet du danger. Accord féminin pluriel : \esp{peligrosas}.

\textbf{Phrase 6.} « Tu as raison sur un point, mais je ne suis pas d'accord avec ta conclusion. »\\
\esp{Tienes razón en un punto, pero no estoy de acuerdo con tu conclusión.}\\
Justification : « avoir raison » = \esp{tener razón} (jamais *\esp{ser razón}) ; « sur un point » = \esp{en un punto}. Le désaccord nuancé se construit avec \esp{pero} + \esp{no estar de acuerdo con}. Notez l'accent : \esp{conclusión}.
\end{exoresolu}

\courssub{Version : un débat sur la vie citoyenne}

Lisez d'abord le texte en entier, sans dictionnaire, pour en saisir le sens global. Repérez ensuite les formules d'accord et de désaccord, les concessions et les modalisateurs.

\begin{textebox}[¿Se interesan los jóvenes por la política?]
No estoy de acuerdo con quienes dicen que los jóvenes no se interesan por la política. Al contrario, me parece que participan de otra manera. Es verdad que votan poco, pero se movilizan en las redes sociales, firman peticiones en línea y defienden causas concretas: la educación, el medio ambiente, la justicia. Tienen razón quienes afirman que el voto es un deber cívico; sin embargo, no creo que sea la única forma de compromiso. Puede que los partidos tradicionales no les representen, pero eso no significa que sean indiferentes. En mi opinión, la juventud no está alejada de la vida pública: simplemente la vive con otras herramientas y con otro lenguaje.
\end{textebox}

\textbf{Glossaire} : \esp{quienes dicen} = ceux qui disent ; \esp{se movilizan} = ils se mobilisent ; \esp{firman peticiones en línea} = ils signent des pétitions en ligne ; \esp{defienden causas} = ils défendent des causes ; \esp{el medio ambiente} = l'environnement ; \esp{el voto} = le vote ; \esp{un deber cívico} = un devoir civique ; \esp{sin embargo} = cependant ; \esp{el compromiso} = l'engagement ; \esp{no les representen} = ne les représentent pas (subjonctif) ; \esp{indiferentes} = indifférents ; \esp{alejada de} = éloignée de ; \esp{las herramientas} = les outils.

\begin{exoresolu}[Traduction commentée de la version]
Énoncé : traduisez le texte « ¿Se interesan los jóvenes por la política? » en français.
\tcblower
\textbf{Traduction proposée :}

« Je ne suis pas d'accord avec ceux qui disent que les jeunes ne s'intéressent pas à la politique. Au contraire, il me semble qu'ils participent autrement. Il est vrai qu'ils votent peu, mais ils se mobilisent sur les réseaux sociaux, signent des pétitions en ligne et défendent des causes concrètes : l'éducation, l'environnement, la justice. Ceux qui affirment que le vote est un devoir civique ont raison ; cependant, je ne crois pas que ce soit la seule forme d'engagement. Il se peut que les partis traditionnels ne les représentent pas, mais cela ne signifie pas qu'ils soient indifférents. À mon avis, la jeunesse n'est pas éloignée de la vie publique : elle la vit simplement avec d'autres outils et avec un autre langage. »

\textbf{Commentaires :}
\begin{itemize}
\item \esp{Al contrario} se rend naturellement par « au contraire » : équivalent direct, registre identique.
\item \esp{me parece que} = « il me semble que » : plus naturel que « il me paraît que », soutenu en français.
\item \esp{Es verdad que... pero...} = « il est vrai que... mais... » : la concession espagnole se rend mot à mot ici, car le français possède la même structure (indicatif des deux côtés).
\item \esp{Tienen razón quienes afirman...} : inversion espagnole (VSU) pour la focalisation ; le français rétablit l'ordre canonique SVO : « ceux qui affirment... ont raison ».
\item \esp{no creo que sea} : subjonctif espagnol rendu par le subjonctif français « que ce soit » : les deux langues coïncident.
\item \esp{Puede que no les representen} = « il se peut que... ne les représentent pas » : subjonctif dans les deux langues.
\item \esp{no significa que sean} : \esp{significar que} à la forme négative appelle le subjonctif (doute/négation) : \esp{sean} = « qu'ils soient ».
\end{itemize}
\end{exoresolu}

\begin{methode}[En version, ne jamais traduire les formules mot à mot]
Une formule d'opinion est un bloc figé : elle se traduit par un bloc figé de l'autre langue, choisi selon le \textbf{ton} du texte. Ainsi \esp{¡qué va !} se rendra par « allons donc ! » ou « pas du tout ! » selon le contexte, jamais par *« quoi va ! ». De même : \esp{¡ni mucho menos !} = « loin de là ! » ; \esp{¡de ninguna manera !} = « en aucun cas ! » ; \esp{¡por supuesto !} = « bien sûr ! » ; \esp{¡efectivamente !} = « en effet ! » ou « exactement ! ». Demandez-vous toujours : « comment un Français dirait-il cela, dans cette situation ? »
\end{methode}

\courssub{Reformulation : doser l'accord et le désaccord}

Le Bac et la vie citoyenne exigent de savoir \cle{nuancer}. Deux transformations sont à maîtriser pour passer du registre familier/brutal au registre soutenu/diplomatique.

\textbf{1. D'un accord total à un accord nuancé.} On ajoute un limiteur (\esp{en parte}, \esp{hasta cierto punto}, \esp{en gran parte}) ou une restriction introduite par \esp{pero}, \esp{aunque}, \esp{sin embargo} :
\begin{itemize}
\item Accord total : \esp{Estoy totalmente de acuerdo contigo.} « Je suis tout à fait d'accord avec toi. »
\item Accord nuancé : \esp{Estoy de acuerdo contigo hasta cierto punto, pero no comparto tu conclusión.} « Je suis d'accord avec toi jusqu'à un certain point, mais je ne partage pas ta conclusion. »
\item Autre possibilité : \esp{En parte tienes razón, aunque hay matices.} « En partie, tu as raison, bien qu'il y ait des nuances. »
\end{itemize}

\textbf{2. D'un désaccord brutal à un désaccord poli.} On remplace la formule catégorique par une concession suivie d'une réserve, ou par un verbe d'opinion atténué (\esp{comprender}, \esp{no creer que}, \esp{no estar seguro de}) :
\begin{itemize}
\item Brutal : \esp{¡De ninguna manera ! Te equivocas.} « Pas question ! Tu te trompes. »
\item Poli : \esp{Entiendo tu punto de vista, pero no estoy de acuerdo.} « Je comprends ton point de vue, mais je ne suis pas d'accord. »
\item Très diplomatique : \esp{No sé si comparto esa idea; puede que haya otras soluciones.} « Je ne sais pas si je partage cette idée ; il se peut qu'il y ait d'autres solutions. »
\end{itemize}

\begin{exercice}[Reformulation]
Transformez les phrases suivantes selon l'indication.
\begin{enumerate}
\item Transformez en accord nuancé : \esp{Por supuesto, tienes razón en todo.}
\item Transformez en désaccord poli : \esp{¡Qué va ! Eso es falso.}
\end{enumerate}
\end{exercice}

\begin{corrige}[Reformulation]
\begin{enumerate}
\item \esp{Tienes razón en gran parte, aunque no en todo.} « Tu as raison en grande partie, bien que pas en tout. » Le limiteur \esp{en gran parte} remplace l'accord total \esp{en todo} ; \esp{aunque} introduit la restriction.
\item \esp{Comprendo lo que dices, pero no creo que sea exactamente así.} « Je comprends ce que tu dis, mais je ne crois pas que ce soit exactement ainsi. » La concession \esp{comprendo} adoucit le refus, et \esp{no creo que} + subjonctif (\esp{sea}) remplace l'affirmation brutale \esp{es falso}.
\end{enumerate}
\end{corrige}

\courssub{Complément Bac : élargir l'expression de l'opinion personnelle}

\begin{aretenir}[Complément Bac : les introducteurs d'opinion pour l'essai argumenté]
Pour enrichir votre essai argumenté (expression écrite) et varier vos thèmes, ajoutez à votre répertoire ces introducteurs d'opinion personnelle. Ils sont suivis de l'\textbf{indicatif} (opinion affirmée comme telle). Placés en début de paragraphe, ils structurent l'essai et montrent au correcteur une maîtrise de la modalisation.
\begin{itemize}
\item \esp{A mi juicio, la educación cívica es indispensable.} « À mon sens, l'éducation civique est indispensable. »
\item \esp{Desde mi punto de vista, el voto es un derecho y un deber.} « De mon point de vue, le vote est un droit et un devoir. »
\item \esp{En mi opinión, los jóvenes deben participar más.} « À mon avis, les jeunes doivent participer davantage. »
\item \esp{Según mi criterio / A mi modo de ver, la participación es clave.} « Selon mon jugement / À ma façon de voir, la participation est clé. »
\end{itemize}
\end{aretenir}

\begin{savaistu}[Ces formules valent des points au Bac]
Les sujets de Bac camerounais proposent régulièrement des phrases de thème portant sur la vie citoyenne : le vote, l'éducation, les réseaux sociaux, l'engagement des jeunes, parfois en lien avec l'actualité de la Guinée équatoriale, seul pays hispanophone d'Afrique. Maîtriser les formules de cette leçon (\esp{estar de acuerdo con}, \esp{tener razón}, \esp{oponerse a}, \esp{interesarse por}, \esp{puede que} + subjonctif, \esp{aunque} + indicatif concessionnel), c'est sécuriser deux à trois points de grammaire au thème — et autant en expression écrite, où un désaccord poliment argumenté est toujours mieux noté qu'un refus brutal.
\end{savaistu}

\courssub{Bilan : grille d'auto-évaluation}

Évaluez-vous honnêtement après cette leçon. Cochez mentalement chaque ligne ; toute case non validée doit vous renvoyer à la section correspondante.

\begin{center}
\begin{tabularx}{\linewidth}{|X|X|}
\hline
\rowcolor{popgreenL} \textbf{Compétence} & \textbf{Je sais...} \\
\hline
Je sais dire & exprimer un accord simple (\esp{estoy de acuerdo con}), un accord total (\esp{totalmente de acuerdo}), un désaccord (\esp{no estoy de acuerdo}), un refus catégorique (\esp{de ninguna manera}) \\
\hline
Je sais construire & une phrase avec \esp{tener razón}, \esp{oponerse a} + nom, \esp{interesarse por}, \esp{votar por}, avec la bonne préposition et le bon auxiliaire \\
\hline
Je sais nuancer & atténuer par \esp{hasta cierto punto}, \esp{en parte} ; concéder par \esp{aunque} + indicatif et \esp{es verdad que... pero...} ; douter par \esp{puede que} et \esp{no creo que} + subjonctif \\
\hline
Je sais traduire & une formule figée par son équivalent naturel, sans mot à mot, en adaptant le ton (tableau contrastif) \\
\hline
Je sais rédiger & un paragraphe d'opinion introduit par \esp{a mi juicio}, \esp{desde mi punto de vista}, \esp{en mi opinión} \\
\hline
\end{tabularx}
\end{center}

Si toutes les cases sont validées, vous êtes prêt pour le thème et la version du baccalauréat sur ce point de grammaire. Sinon, reprenez la frise méthodologique : lire, identifier, choisir, vérifier, relire. \esp{¡Seguro que lo conseguirás!} « Tu y arriveras, c'est sûr ! »

\newpage

\coursec{Activité d'intégration}

\coursec{Leçon E15 : Traduction — Expression de l'accord et du désaccord}

\courssub{Observation : l'accord et le désaccord en contexte}

\begin{methode}[Comment aborder un texte de débat citoyen]
\begin{enumerate}
\item \textbf{Identifier le thème et la motion} : repérer le sujet (ex. vote à 16 ans) et la question posée (\esp{¿Debería ser obligatorio...?}).
\item \textbf{Repérer les acteurs et les arguments} : qui parle ? Quels sont les arguments \esp{A FAVOR} (pour) et \esp{EN CONTRA} (contre) ?
\item \textbf{Surligner les marqueurs d'opinion} : formules d'accord (\esp{estoy de acuerdo}), de désaccord (\esp{no estoy de acuerdo}), de concession (\esp{es cierto que... pero}), d'atténuation (\esp{creo que, quizás}).
\item \textbf{Noter les structures grammaticales associées} : \esp{estar de acuerdo con}, \esp{oponerse a} + infinitif/\esp{que} + subjonctif, \esp{aunque} + indicatif/subjonctif.
\end{enumerate}
\end{methode}

\begin{textebox}[Motion du débat — Club bilingüe, Yaoundé]
\textbf{¿Debería ser obligatorio el voto a partir de los 16 años en los países africanos e hispanohablantes?}

\emph{Contexto:} En varios países latinoamericanos, como Argentina, Ecuador o Nicaragua, los jóvenes pueden votar desde los 16 años. Sus defensores afirman que votar joven crea ciudadanos más comprometidos. Sus detractores sostienen que a los 16 años falta madurez política y que el voto debe seguir siendo un derecho reservado a los mayores de edad. ¿Y tú, qué opinas?
\end{textebox}

\textbf{Traduction de la motion :} « Le vote devrait-il être obligatoire à partir de 16 ans dans les pays africains et hispanophones ? » En Argentine, en Équateur ou au Nicaragua, les jeunes peuvent voter dès 16 ans. Ses défenseurs affirment que voter jeune crée des citoyens plus engagés ; ses détracteurs soutiennent qu'à 16 ans il manque la maturité politique et que le vote doit rester un droit réservé aux majeurs.

\begin{savaistu}[Le vote à 16 ans dans l'espace hispanophone et en Guinée équatoriale]
En \textbf{Guinée équatoriale}, pays africain hispanophone, l'âge du vote est fixé à \textbf{18 ans} (Constitution, art. 15). En \textbf{Espagne}, c'est aussi 18 ans. En \textbf{Amérique latine}, plusieurs pays ont abaissé l'âge à 16 ans pour les élections nationales : \textbf{Argentine} (2012), \textbf{Équateur}, \textbf{Nicaragua}, \textbf{Brésil} (luso-phones, mais contexte latino), \textbf{Cuba} (élections municipales dès 16 ans). L'\textbf{Autriche} (européen) l'a fait en 2007. Au \textbf{Cameroun}, l'âge électoral est de 20 ans pour les législatives et municipales, 18 ans pour le référendum. Le débat sur la majorité civique à 16 ans traverse donc tout l'espace atlantique.
\end{savaistu}

\courssub{Les formules d'accord : formes et emplois}

\begin{aretenir}[Répertoire : Exprimer l'accord]
\begin{itemize}
\item \textbf{Accord simple / standard :} \esp{Estoy de acuerdo (con...)} ; \esp{Tienes razón} ; \esp{Opino lo mismo (que tú/él...)} ; \esp{Así es}.
\item \textbf{Accord renforcé / enthousiaste :} \esp{Estoy \textbf{totalmente} de acuerdo} ; \esp{\textbf{Por supuesto}} ; \esp{\textbf{Desde luego}} ; \esp{\textbf{Claro que sí}} ; \esp{\textbf{Efectivamente}} ; \esp{\textbf{Exacto}}.
\item \textbf{Accord partiel / nuancé :} \esp{Estoy de acuerdo \textbf{hasta cierto punto}} ; \esp{Comparto tu opinión \textbf{en parte}}.
\end{itemize}
\end{aretenir}

\begin{propriete}[Grammaire de l'accord : constructions obligatoires]
\begin{center}
\begin{tabularx}{\linewidth}{|X|X|X|}
\hline
\rowcolor{popblueL} \textbf{Formule} & \textbf{Construction} & \textbf{Exemple traduit} \\
\hline
\esp{Estar de acuerdo} & \esp{Estar} (conjugué) \textbf{de acuerdo} \esp{con} + \textbf{nom / pronom / infinitif} & \esp{Estoy de acuerdo con tu propuesta.} (Je suis d'accord avec ta proposition.) \\
\hline
\esp{Tener razón} & \esp{Tener} (conjugué) \textbf{razón} \textbf{(sans article)} & \esp{Tienes razón.} (Tu as raison.) \textbf{Pas} \esp{*tienes la razón.} \\
\hline
\esp{Opinar lo mismo} & \esp{Opinar} (conjugué) \textbf{lo mismo} \esp{que} + \textbf{indicatif} & \esp{Opino lo mismo que tú.} (J'ai le même avis que toi.) \\
\hline
\end{tabularx}
\end{center}

\end{propriete}

\begin{attention}[Piège majeur : \esp{*Soy de acuerdo}]
\esp{Estar de acuerdo} se construit avec le verbe \textbf{\esp{estar}} (état, résultat), jamais avec \esp{ser}.
\begin{itemize}
\item \checkmark \esp{Estoy de acuerdo.}
\item \cross \esp{*Soy de acuerdo.} (Calque du français « Je \textbf{suis} d'accord ».)
\end{itemize}
De même : \esp{Estar en contra de} (être contre), \esp{Estar a favor de} (être pour).
\end{attention}

\courssub{Les formules de désaccord : formes et emplois}

\begin{aretenir}[Répertoire : Exprimer le désaccord]
\begin{itemize}
\item \textbf{Désaccord poli / standard :} \esp{No estoy de acuerdo (con...)} ; \esp{No comparto tu opinión} ; \esp{No lo veo así} ; \esp{No estoy seguro de eso}.
\item \textbf{Désaccord ferme / catégorique :} \esp{De ninguna manera} ; \esp{De ningún modo} ; \esp{Ni hablar} (familier) ; \esp{Qué va} (très familier, Am. Latine) ; \esp{Me opongo (a esa idea / a que...)}.
\item \textbf{Désaccord argumentatif :} \esp{Al contrario} ; \esp{Por el contrario} ; \esp{Discrepo (de / en)}.
\end{itemize}
\end{aretenir}

\begin{propriete}[Grammaire du désaccord : constructions et prépositions]
\begin{center}
\begin{tabularx}{\linewidth}{|X|X|X|}
\hline
\rowcolor{poporangeL} \textbf{Formule} & \textbf{Construction} & \textbf{Exemple traduit} \\
\hline
\esp{No estar de acuerdo} & \esp{No estar} \textbf{de acuerdo} \esp{con} + \textbf{nom / pronom / infinitif} & \esp{No estoy de acuerdo con esa medida.} (Je ne suis pas d'accord avec cette mesure.) \\
\hline
\esp{Opinarse a} & \esp{Opinarse} \textbf{a} + \textbf{infinitif / nom} \newline \esp{Opinarse} \textbf{a que} + \textbf{subjonctif} & \esp{Me opongo a la ley.} / \esp{Me opongo a que sea obligatorio.} (Je m'oppose à la loi / à ce qu'il soit obligatoire.) \\
\hline
\esp{Discrepar} & \esp{Discrepar} \textbf{de} (personne) / \textbf{en} (point précis) & \esp{Discrepo de tu análisis.} / \esp{Discrepo en ese punto.} \\
\hline
\end{tabularx}
\end{center}

\end{propriete}

\begin{attention}[Pièges du désaccord]
\begin{itemize}
\item \textbf{\esp{Tienes razón}} : pas d'article défini. \cross \esp{*Tienes la razón}. Pour « tu as raison \textbf{de} + infinitif », on n'utilise pas \esp{razón} : \esp{Haces bien en estudiar.} (Tu as raison d'étudier.)
\item \textbf{\esp{Al contrario}} : sans article (\cross \esp{*al el contrario}). Équivaut à « Au contraire ».
\item \textbf{\esp{De ninguna manera} vs \esp{De ningún modo}} : tous deux = « En aucun cas / Absolument pas ». \esp{De todas formas / De todos modos} = « De toute façon / Quoi qu'il en soit ». Ne pas confondre.
\item \textbf{\esp{Me opongo}} : \textbf{préposition \esp{a} obligatoire}. \esp{Me opongo \textbf{a} esa idea} / \esp{Me opongo \textbf{a} que \textbf{vengas}} (subjonctif).
\end{itemize}
\end{attention}

\courssub{Modalisation : atténuer et concéder}

\begin{aretenir}[Répertoire : Nuancer son opinion (Modalisation)]
\begin{itemize}
\item \textbf{Concession (admettre un point avant de réfuter) :}
      \esp{Es cierto que... , \textbf{pero / sin embargo / no obstante}...} ;
      \esp{Aunque \textbf{tengas} razón en parte, ...} (subjonctif : concession hypothétique) ;
      \esp{Aunque \textbf{tienes} razón en eso, ...} (indicatif : fait admis, certain) ;
      \esp{Entiendo lo que dices, \textbf{sin embargo / no obstante}...} ;
      \esp{Reconozco que... , \textbf{pero}...}
\item \textbf{Atténuation (exprimer un doute, une probabilité, subjectivité) :}
      \esp{Creo que / Pienso que / Opino que} + \textbf{indicatif} (opinion personnelle) ;
      \esp{Me parece que} + \textbf{indicatif / subjonctif} (selon certitude) ;
      \esp{Puede que / Es posible que / Quizás / Tal vez} + \textbf{subjonctif} ;
      \esp{A mi juicio / En mi opinión / Desde mi punto de vista} + \textbf{indicatif} ;
      \esp{Hasta cierto punto} (jusqu'à un certain point).
\end{itemize}
\end{aretenir}

\begin{propriete}[Le choix du mode après \esp{Aunque} et verbes d'opinion]
\begin{center}
\begin{tabularx}{\linewidth}{|X|X|X|}
\hline
\rowcolor{poppurpleL} \textbf{Structure} & \textbf{Mode} & \textbf{Valeur / Exemple} \\
\hline
\esp{Aunque} + \textbf{indicatif} & Indicatif & Fait \textbf{connu, admis, certain}. \esp{Aunque \textbf{tienes} razón, no estoy de acuerdo.} (Bien que tu aies raison [fait], je ne suis pas d'accord.) \\
\hline
\esp{Aunque} + \textbf{subjonctif} & Subjonctif & Fait \textbf{hypothétique, non réalisé, mis en doute}. \esp{Aunque \textbf{tengas} razón, no cambiaré de opinión.} (Même si tu as raison [hypothèse], je ne changerai pas d'avis.) \\
\hline
\esp{Creo que / Pienso que / Opino que} & \textbf{Indicatif} & Affirmation subjective mais présentée comme vraie. \esp{Creo que \textbf{es} buena idea.} \\
\hline
\esp{No creo que / No pienso que} & \textbf{Subjonctif} & Négation de l'opinion $\rightarrow$ doute. \esp{No creo que \textbf{sea} buena idea.} \\
\hline
\esp{Me parece que} & Indicatif / Subjonctif & Indicatif si certitude (\esp{Me parece que \textbf{viene}}), Subjonctif si doute (\esp{No me parece que \textbf{venga}}). \\
\hline
\esp{Puede que / Quizás / Tal vez} & \textbf{Subjonctif} (standard) & Probabilité. \esp{Quizás \textbf{venga} mañana.} (Indicatif possible en Am. Lat. pour \esp{Quizás/Tal vez} si forte probabilité, mais Subjonctif plus sûr au Bac.) \\
\hline
\end{tabularx}
\end{center}

\end{propriete}

\courssub{Contrastif français/espagnol et pièges de traduction (Thème / Version)}

\begin{attention}[Tableau récapitulatif des faux-amis et calques fréquents]
\begin{center}
\begin{tabularx}{\linewidth}{|X|X|X|}
\hline
\rowcolor{popredL} \textbf{Français} & \textbf{Espagnol correct} & \textbf{Erreur fréquente (Calque / Faux ami)} \\
\hline
Je suis d'accord & \esp{Estoy de acuerdo} & \cross \esp{*Soy de acuerdo} (Ser vs Estar) \\
\hline
Tu as raison & \esp{Tienes razón} & \cross \esp{*Tienes la razón} / \cross \esp{*Tienes razón de...} \\
\hline
Au contraire & \esp{Al contrario} & \cross \esp{*Au contraire} / \cross \esp{*En el contrario} \\
\hline
En aucun cas / Absolument pas & \esp{De ninguna manera} / \esp{De ningún modo} & \cross \esp{*De ninguna forma de todas formas} (mélange) \\
\hline
Je m'oppose à cela & \esp{Me opongo a eso} & \cross \esp{*Me opongo eso} (préposition \esp{a} oubliée) \\
\hline
C'est un droit, pas un devoir & \esp{Es un derecho, no una obligación} & \cross \esp{*No un deber} (\esp{deber} = devoir moral / devoirs scolaires ; \esp{obligación} = obligation légale) \\
\hline
Point de vue & \esp{Punto de vista} & \cross \esp{*Punto de visión} (\esp{visión} = vision physiologique / visionnaire) \\
\hline
Rester (demeurer) & \esp{Seguir siendo} / \esp{Permanecer} & \cross \esp{*Rester} (n'existe pas) / \cross \esp{*Quedar} (rester dans un lieu) \\
\hline
Obligatoire & \esp{Obligatorio} & \cross \esp{*Obligatoire} (orthographe fr.) \\
\hline
\end{tabularx}
\end{center}
\end{attention}

\courssub{Entraînement à la traduction : thème et version}

\begin{methode}[Méthode de traduction Thème / Version (Phrases d'opinion)]
\begin{enumerate}
\item \textbf{Analyser la structure française} : identifier le verbe d'opinion (\esp{penser, croire, être d'accord}), la négation, les connecteurs (\esp{mais, cependant, bien que}).
\item \textbf{Choisir le verbe espagnol adapté} : \esp{Estar de acuerdo} (accord), \esp{Tener razón} (avoir raison), \esp{Opinar} (opiner), \esp{Creer/Pensar} (croire/penser).
\item \textbf{Appliquer la concordance des temps et des modes} :
      \begin{itemize}
      \item \esp{Bien que / Aunque} : Indicatif (fait réel) vs Subjonctif (hypothétique).
      \item \esp{Il faut que / Es necesario que} : Subjonctif.
      \item \esp{Je ne pense pas que / No creo que} : Subjonctif.
      \end{itemize}
\item \textbf{Vérifier les prépositions} : \esp{de acuerdo \textbf{con}}, \esp{oponerse \textbf{a}}, \esp{razón} (sans article).
\item \textbf{Relire à voix haute} pour contrôler la prosodie et les accords.
\end{enumerate}
\end{methode}

\begin{exoresolu}[Compte rendu modèle du débat (8 phrases — Cahier des charges : 3 accords, 3 désaccords, 2 concessions)]
Rédiger un compte rendu structuré du débat citoyen en mobilisant le répertoire lexical et grammatical.
\tcblower
\textbf{Production modèle :}

\esp{El debate sobre el voto a los 16 años fue muy animado. \textbf{Estoy totalmente de acuerdo con} los que dicen que votar joven crea ciudadanos comprometidos. \textbf{Es cierto que} a los 16 años \textbf{falta a veces madurez política}, \textbf{pero} la escuela puede preparar a los jóvenes. \textbf{Sin embargo, me opongo a que} el voto \textbf{sea} obligatorio: debe seguir siendo un derecho, no una obligación. \textbf{De ninguna manera} se puede obligar a un adolescente a elegir a sus representantes. Mi compañera Amina \textbf{tiene razón} cuando afirma que en Argentina los jóvenes votan desde los 16 años. \textbf{Aunque entiendo} los argumentos del otro grupo, \textbf{no estoy de acuerdo con} su visión negativa de la juventud. En conclusión, \textbf{efectivamente}, el voto joven es positivo, \textbf{sólo} si es voluntario.}

\textbf{Traduction (Version) :}
« Le débat sur le vote à 16 ans fut très animé. \textbf{Je suis totalement d'accord avec} ceux qui disent que voter jeune crée des citoyens engagés. \textbf{Il est vrai qu'}à 16 ans \textbf{il manque parfois la maturité politique}, \textbf{mais} l'école peut préparer les jeunes. \textbf{Cependant, je m'oppose à ce que} le vote \textbf{soit} obligatoire : il doit rester un droit, non une obligation. \textbf{En aucun cas} on ne peut obliger un adolescent à choisir ses représentants. Ma camarade Amina \textbf{a raison} quand elle affirme qu'en Argentine les jeunes votent dès 16 ans. \textbf{Bien que je comprenne} les arguments de l'autre groupe, \textbf{je ne suis pas d'accord avec} sa vision négative de la jeunesse. En conclusion, \textbf{effectivement}, le vote jeune est positif, \textbf{seulement} s'il est volontaire. »

\textbf{Commentaire linguistique détaillé :}
\begin{itemize}
\item \esp{\textbf{Estoy totalmente de acuerdo con}} : Accord renforcé (\esp{totalmente}), construction \esp{estar de acuerdo} + \esp{con} + nom. \checkmark
\item \esp{\textbf{Es cierto que... pero...}} : Concession standard (indicatif \esp{falta} car fait constaté) + opposition (\esp{pero}). \checkmark
\item \esp{\textbf{Me opongo a que} el voto \textbf{sea}} : \esp{Opinarse a que} + \textbf{Subjonctif présent} (\esp{sea} de \esp{ser}) car rejet/volonté négative. \checkmark
\item \esp{\textbf{De ninguna manera}} : Désaccord catégorique en tête de phrase (topicalisation emphatique). Structure impersonnelle \esp{se puede}. \checkmark
\item \esp{\textbf{Tiene razón}} : Sans article. Verbe \esp{tener}. \checkmark
\item \esp{\textbf{Aunque entiendo}... \textbf{no estoy de acuerdo}} : \esp{Aunque} + \textbf{Indicatif} (\esp{entiendo}) car le locuteur admet la validité des arguments (fait réel). Désaccord poli \esp{no estoy de acuerdo con}. \checkmark
\item \esp{\textbf{Efectivamente}} / \esp{\textbf{Sólo}} : Reprise variée (accord final). \textbf{Accent sur \esp{sólo}} (adverbe = seulement) pour distinguer de l'adjectif \esp{solo} (seul), norme exigeante au Bac. \checkmark
\end{itemize}
\end{exoresolu}

\begin{exoresolu}[Thème guidé : 2 phrases à traduire en espagnol]
Traduire en espagnol en justifiant les choix grammaticaux.
\tcblower
\textbf{1. « Tu as raison : le vote doit rester un droit, pas une obligation. »}
\par\esp{Tienes razón: el voto debe \textbf{seguir siendo} un derecho, no una obligación.}
\begin{itemize}
\item \esp{Tienes razón} : pas d'article.
\item \esp{Debe seguir siendo} : \esp{Deber} (obligation) + infinitif composé \esp{seguir siendo} (periphrase aspectuelle « continuer à être / rester »). \esp{Quedar} ou \esp{permanecer} moins idiomatiques ici.
\end{itemize}

\textbf{2. « Je ne suis pas d'accord, bien que je comprenne ton point de vue. »}
\par\esp{No estoy de acuerdo, \textbf{aunque entiendo} tu \textbf{punto de vista}.}
\begin{itemize}
\item \esp{No estoy de acuerdo} : \esp{Estar} + \esp{de acuerdo} (sans \esp{con} si contexte clair).
\item \esp{Aunque entiendo} : \esp{Aunque} + \textbf{Indicatif} (\esp{entiendo}) car « je comprends » est un fait réel, admis par le locuteur au moment de parler.
\item \esp{Punto de vista} : Locution figée. \cross \esp{*Punto de visión}.
\end{itemize}
\end{exoresolu}

\begin{exercice}[Entraînement individuel — Thème et Version]
\textbf{A. Version (Espagnol $\rightarrow$ Français) :} Traduire les interventions suivantes d'élèves lors du débat.
\begin{enumerate}
\item \esp{— Por supuesto que los jóvenes de 16 años están preparados. Desde luego, en mi país, Ecuador, votan y nada malo pasa.}
\item \esp{— No comparto tu opinión. Me parece que a esa edad falta madurez. De ningún modo apoyo el voto obligatorio.}
\item \esp{— Es cierto que algunos son maduros, aunque no todos. Puede que sea mejor esperar a los 18 años.}
\item \esp{— Entiendo tu postura, sin embargo, me opongo a que se obligue a nadie. El voto es un derecho, no un deber.}
\end{enumerate}

\textbf{B. Thème (Français $\rightarrow$ Espagnol) :} Traduire en espagnol en soignant les structures d'opinion.
\begin{enumerate}
\item « Je pense que voter à 16 ans est une bonne idée, même si certains manquent de maturité. »
\item « Nous ne sommes pas d'accord avec l'obligation. Au contraire, le vote doit être volontaire. »
\item « Bien que tu aies raison sur l'engagement civique, je m'oppose à la contrainte légale. »
\item « À mon avis, l'école devrait éduquer à la citoyenneté avant d'imposer le vote. »
\end{enumerate}
\end{exercice}

\begin{corrige}[Exercice Entraînement individuel]
\textbf{A. Version}
\begin{enumerate}
\item « Bien sûr que les jeunes de 16 ans sont préparés. Certainement, dans mon pays, l'Équateur, ils votent et il ne se passe rien de mal. » (\esp{Por supuesto / Desde luego} = accord fort ; \esp{nada malo pasa} = rien de mal ne se passe).
\item « Je ne partage pas ton avis. Il me semble qu'à cet âge manque la maturité. D'aucune manière / En aucun cas je ne soutiens le vote obligatoire. » (\esp{No comparto} = désaccord poli ; \esp{Me parece que} + indicatif = opinion ; \esp{De ningún modo} = désaccord ferme).
\item « Il est vrai que certains sont mûrs, bien que pas tous. Il se peut qu'il soit mieux d'attendre 18 ans. » (\esp{Es cierto que} concession ; \esp{aunque no todos} = ellipsis de \esp{sean maduros}, indicatif car fait admis ; \esp{Puede que} + subjonctif \esp{sea} = probabilité).
\item « Je comprends ta position, toutefois, je m'oppose à ce qu'on oblige qui que ce soit. Le vote est un droit, pas un devoir. » (\esp{Entiendo... sin embargo} concession polie ; \esp{Me opongo a que} + subjonctif \esp{se obligue} (passif réfléchi) ; \esp{nadie} = personne / qui que ce soit).
\end{enumerate}

\textbf{B. Thème}
\begin{enumerate}
\item \esp{Creo que votar a los 16 años es una buena idea, \textbf{aunque algunos falten} a la madurez.} (\esp{Creo que} + indicatif \esp{es} ; \esp{aunque} + \textbf{subjonctif} \esp{falten} car « certains manquent » est vu comme une généralité/hypothèse non identifiée, ou indicatif \esp{faltan} si fait constaté. \textbf{Au Bac : subjonctif plus sûr après \esp{aunque} si nuance de concession hypothétique}). \textit{Alternative acceptée :} \esp{aunque algunos \textbf{falten} / \textbf{falten}...}
\item \esp{No estamos de acuerdo con la obligación. \textbf{Al contrario}, el voto debe ser voluntario.} (\esp{Estar de acuerdo con} ; \esp{Al contrario} sans article ; \esp{debe ser} obligation externe).
\item \esp{\textbf{Aunque tengas} razón sobre el compromiso cívico, \textbf{me opongo a} la obligación legal.} (\esp{Aunque} + \textbf{subjonctif} \esp{tengas} : « bien que tu aies raison » = concession hypothétique/polie, standard à l'écrit formel ; \esp{Me opongo a} + nom).
\item \esp{\textbf{A mi juicio / En mi opinión}, la escuela \textbf{debería educar} para la ciudadanía \textbf{antes de imponer} el voto.} (\esp{Debería} + infinitif = conseil / opinion conditionnelle ; \esp{antes de} + infinitif = antériorité).
\end{enumerate}
\end{corrige}

\courssub{Bilan de la leçon}

Cette leçon a permis de couvrir l'intégralité des attendus du programme officiel (Module 3, Terminale A) pour l'expression de l'accord et du désaccord :
\begin{itemize}
\item \textbf{Lexique :} Maîtrise d'un répertoire varié (\esp{estoy de acuerdo, tienes razón, por supuesto, efectivamente, no estoy de acuerdo, de ninguna manera, al contrario, me opongo}) et de leurs variantes de registre (formel / familier).
\item \textbf{Grammaire :} Consolidation des constructions syntaxiques obligatoires (\esp{estar de acuerdo con}, \esp{tener razón} sans article, \esp{oponerse a} + inf./\esp{que}+subj., \esp{discrepar de/en}) et de la modalisation (\esp{Aunque} + Ind./Subj., \esp{Creer/No creer} + Ind./Subj., \esp{Puede que} + Subj.).
\item \textbf{Traduction (Thème/Version) :} Entraînement intensif aux pièges contrastifs (calques \esp{*soy de acuerdo}, faux-amis \esp{*punto de visión}, prépositions, accents diacritiques \esp{sólo}).
\item \textbf{Compétences communicatives :} Interaction orale en temps réel (débat citoyen), médiation linguistique (rôle de traducteur), production écrite structurée (compte rendu), réception écrite (motion).
\end{itemize}
L'élève est désormais outillé pour réussir les épreuves du Baccalauréat (commentaire de texte, thème, version, expression écrite) et pour participer activement à la vie citoyenne dans un contexte bilingue camerounais ou hispanophone.

\newpage

\coursec{Méthodes et exercices résolus}

\coursec{Leçon E15 : Traduction — Expression de l'accord et du désaccord}

\begin{savaistu}[Le Cameroun et l'espace hispanophone : un atout majeur]
Le Cameroun partage une frontière avec la \textbf{Guinée équatoriale} (pays hispanophone depuis 1778, membre de la CEMAC et de l'Union africaine). L'espagnol y est langue officielle aux côtés du français et du portugais. Des milliers de Camerounais y travaillent (bois, pêche, commerce frontalier à Kye-Ossi, Ebebiyín). Maîtriser l'expression de l'accord et du désaccord en espagnol, c'est pouvoir négocier, débattre et travailler efficacement dans cet espace de proximité. Au Bac, les correcteurs valorisent les exemples ancrés dans ce contexte réel (ex. : \esp{Guinea Ecuatorial}, \esp{Kye-Ossi}, \esp{Malabo}, \esp{CEMAC}).
\end{savaistu}

\courssub{Observation : l'accord et le désaccord en contexte}

\begin{textebox}[Dialogue au marché de Mokolo (Yaoundé) — Contexte authentique]
— \esp{Buenos días, señor. ¿Está de acuerdo con el precio de este cacao?} \\
— \esp{No, de ninguna manera. Es demasiado bajo. Mi cosecha vale más.} \\
— \esp{Entiendo su punto de vista. Pero, efectivamente, el mercado internacional ha bajado.} \\
— \esp{Al contrario, la demanda sube. Por supuesto, si me paga bien, vendo todo.} \\
— \esp{Vale. Tienes razón. Estoy de acuerdo: 2.500 FCFA el kilo.} \\
— \esp{¡Trato hecho!}
\end{textebox}
\textbf{Traduction et glossaire} : « Bonjour monsieur. Êtes-vous d'accord avec le prix de ce cacao ? — Non, absolument pas. Il est trop bas. Ma récolte vaut plus. — Je comprends votre point de vue. Mais, effectivement, le marché international a baissé. — Au contraire, la demande monte. Bien sûr, si vous me payez bien, je vends tout. — D'accord. Tu as raison. Je suis d'accord : 2 500 FCFA le kilo. — Marché conclu ! » \textbf{Glossaire} : \esp{cosecha} (récolte), \esp{mercado internacional} (marché international), \esp{demanda} (demande), \esp{trato hecho} (marché conclu / c'est entendu).

\begin{propriete}[Constat linguistique : deux verbes d'état pour « être »]
En espagnol, l'accord et le désaccord relèvent de \textbf{\esp{estar}} (état, situation), jamais de \esp{ser} (identité, essence).
\begin{center}
\begin{tabularx}{\linewidth}{|X|X|X|}
\hline
\rowcolor{popblueL} \textbf{Français} & \textbf{Espagnol (Estar)} & \textbf{Remarque} \\
\hline
Je suis d'accord & \esp{Estoy de acuerdo} & \esp{de acuerdo} invariable (locution adjectivale) \\
Tu as raison & \esp{Tienes razón} & Verbe \esp{tener} (avoir) + nom fém. \esp{razón} (accent !) \\
Il est contre & \esp{Está en contra} & \esp{estar en contra de} + nom/infinitif \\
\hline
\end{tabularx}
\end{center}
\end{propriete}

\courssub{Les formules d'accord : formes et emplois}

\begin{propriete}[Échelle de l'accord : du familier au formel]
\begin{center}
\begin{tabularx}{\linewidth}{|>{\bfseries}l|X|X|}
\hline
\rowcolor{popgreenL} \textbf{Degré} & \textbf{Formule espagnole} & \textbf{Équivalent français / Registre} \\
\hline
Total & \esp{Estoy totalmente de acuerdo (con usted/contigo)} & Je suis tout à fait d'accord (standard) \\
\hline
Total & \esp{Por supuesto (que sí)} & Bien sûr / Évidemment (standard) \\
\hline
Approbation & \esp{Tienes toda la razón} / \esp{Llevas razón} & Tu as tout à fait raison / Tu as raison (standard) \\
\hline
Confirmation & \esp{Efectivamente} / \esp{En efecto} & Effectivement / En effet (soutien à un fait) \\
\hline
Évidence & \esp{Claro (que sí)} / \esp{Desde luego} & Clairement / Bien sûr (courant) \\
\hline
Accord partiel & \esp{Estoy de acuerdo en parte} / \esp{Hasta cierto punto} & Je suis d'accord en partie / Jusqu'à un certain point \\
\hline
Familier & \esp{¡Claro!} / \esp{¡Cómo no!} / \esp{Vale} & Ça marche / OK / Bien sûr (oral, amis) \\
\hline
\end{tabularx}
\end{center}
\end{propriete}

\begin{exemplebox}[Exemples d'accord contextualisés (Cameroun / Hispanophone)]
\begin{enumerate}
\item \esp{— ¿Vamos al partido de los Leones Indomables? — ¡Por supuesto! Me encanta el fútbol.} « On va au match des Lions ? — Bien sûr ! J'adore le foot. »
\item \esp{Estoy de acuerdo con la nueva ley de tráfico: hay muchos accidentes en la carretera Douala--Yaoundé.} « Je suis d'accord avec la nouvelle loi sur la circulation : il y a beaucoup d'accidents sur la route Douala-Yaoundé. »
\item \esp{Efectivamente, el español es oficial en Guinea Ecuatorial, nuestro vecino.} « Effectivement, l'espagnol est officiel en Guinée équatoriale, notre voisin. »
\end{enumerate}
\end{exemplebox}

\courssub{Les formules de désaccord : formes et emplois}

\begin{propriete}[Échelle du désaccord : de la politesse à la fermeté]
\begin{center}
\begin{tabularx}{\linewidth}{|>{\bfseries}l|X|X|}
\hline
\rowcolor{poporangeL} \textbf{Degré} & \textbf{Formule espagnole} & \textbf{Équivalent français / Registre} \\
\hline
Poli / Standard & \esp{No estoy de acuerdo (con usted/contigo)} & Je ne suis pas d'accord \\
\hline
Argumenté & \esp{No estoy de acuerdo, porque...} / \esp{Discrepo (de esa opinión)} & Je ne suis pas d'accord, car... / Je diverge (formel) \\
\hline
Opposition & \esp{Al contrario, pienso que...} & Au contraire, je pense que... \\
\hline
Catégorique & \esp{De ninguna manera} / \esp{En absoluto} & En aucune façon / Absolument pas \\
\hline
Refus ferme & \esp{Me opongo (a esa propuesta)} / \esp{Me niego a aceptar...} & Je m'oppose (à cette proposition) / Je refuse d'accepter... \\
\hline
Correction & \esp{No es cierto} / \esp{No es verdad} / \esp{Te equivocas} & Ce n'est pas vrai / Tu te trompes \\
\hline
Familier & \esp{¡Qué va!} / \esp{¡Ni hablar!} / \esp{¡Anda ya!} & Mais non ! / Pas question ! (très oral) \\
\hline
\end{tabularx}
\end{center}
\end{propriete}

\begin{attention}[Piège majeur : \esp{Ser} vs \esp{Estar} et \esp{Haber} vs \esp{Tener}]
\begin{enumerate}
\item \textbf{Jamais} \esp{ser de acuerdo}. Toujours \textbf{\esp{estar de acuerdo}}. \esp{*Soy de acuerdo} $\to$ \textbf{INCORRECT}.
\item \textbf{Jamais} \esp{haber razón}. Toujours \textbf{\esp{tener razón}}. \esp{*Tengo razón} (correct) vs \esp{*He razón} (incorrect).
\item \esp{De acuerdo} ne s'accorde \textbf{jamais} en genre/nom : \esp{Ella está de acuerdo}, \esp{Ellas están de acuerdo}.
\item Préposition obligatoire : \esp{estar de acuerdo \textbf{con} alguien / \textbf{en} algo}. \esp{Estoy de acuerdo contigo / en ese punto}.
\end{enumerate}
\end{attention}

\courssub{Modalisation : atténuer et concéder (la nuance, marque du niveau Terminale)}

\begin{definition}[Modalisation]
La \cle{modalisation} est l'ensemble des procédés linguistiques (adverbes, verbes, connecteurs, modes) qui permettent au locuteur de \textbf{nuancer} son engagement : atténuer une affirmation, exprimer un doute, concéder un point à l'adversaire avant de mieux le réfuter. C'est la marque de la maturité langagière attendue au Baccalauréat.
\end{definition}

\begin{propriete}[Boîte à outils de la nuance : Grammaire et Lexique]
\begin{center}
\begin{tabularx}{\linewidth}{|>{\bfseries}l|X|X|}
\hline
\rowcolor{poppurpleL} \textbf{Fonction} & \textbf{Outils espagnols} & \textbf{Règle grammaticale clé} \\
\hline
\textbf{Atténuer} (doute) & \esp{Quizás / Tal vez} + \textbf{Subjonctif} & \esp{Quizás tenga razón.} (Peut-être a-t-il raison) \\
& \esp{Creo que... un poco / bastante} & Indicatif (opinion affirmative) \\
& \esp{Me parece que...} / \esp{Me da la impresión de que...} & Indicatif (moins fort que \esp{creo}) \\
& \esp{Hasta cierto punto} / \esp{En parte} & Invariables, suivis de \esp{que} + indicatif \\
\hline
\textbf{Concéder} (fait réel) & \esp{Es verdad que...} / \esp{Es cierto que...} + \textbf{Indicatif} & \esp{Es verdad que cansan, pero...} \\
& \esp{Admito que...} + \textbf{Indicatif} & Reconnaissance explicite \\
\hline
\textbf{Concéder} (hypothèse) & \esp{Aunque} + \textbf{Subjonctif} & \esp{Aunque sea difícil, lo haré.} (Même si c'est difficile...) \\
\hline
\textbf{Opposer / Réfuter} & \esp{... pero / sin embargo / no obstante / en cambio} & \esp{pero} (coord.), \esp{sin embargo/no obstante} (adverbes, virgule) \\
\hline
\end{tabularx}
\end{center}
\end{propriete}

\begin{methode}[Choisir le mode après \esp{Aunque} et \esp{Quizás/Tal vez}]
\begin{enumerate}
\item \textbf{\esp{Aunque} + INDICATIF} = Fait réel, connu, admis par les deux. \esp{Aunque \textbf{llueve}, salgo.} (Il pleut, c'est un fait, je sors quand même).
\item \textbf{\esp{Aunque} + SUBJONCTIF} = Hypothèse, éventualité, non réalisé. \esp{Aunque \textbf{llueva}, saldré.} (Même s'il pleut [hypothétique], je sortirai).
\item \textbf{\esp{Quizás / Tal vez} + SUBJONCTIF} = Doute, probabilité faible/moyenne. \esp{Quizás \textbf{venga} mañana.}
\item \textbf{\esp{Quizás / Tal vez} + INDICATIF} = Probabilité forte, quasi-certitude (rare à l'oral, plus littéraire). \esp{Quizás \textbf{viene} mañana.} (Il viendra probablement).
\end{enumerate}
\textbf{Astuce Bac} : En situation d'examen, par défaut, après \esp{quizás/tal vez} exprimant une opinion personnelle incertaine $\to$ \textbf{Subjonctif}.
\end{methode}

\begin{exoresolu}[Transformer une opinion tranchée en opinion nuancée]
Énoncé. Réécrivez les phrases suivantes en les nuancant selon la consigne.
\begin{enumerate}
\item \esp{El español es difícil.} (atténuer : \esp{me parece} + \esp{un poco})
\item \esp{Los deberes son inútiles.} (concéder : \esp{es verdad que... pero...})
\item \esp{Tiene razón.} (atténuer : \esp{quizás} + subjonctif)
\item \esp{No voy a la fiesta.} (concéder hypothèse : \esp{aunque} + subjonctif + \esp{pero})
\end{enumerate}
\tcblower
\textbf{Solutions détaillées} :
\begin{enumerate}
\item \esp{Me parece que el español es un poco difícil.} (Opinion personnelle atténuée, indicatif après \esp{me parece que} affirmatif).
\item \esp{Es verdad que los deberes cansan a los alumnos, pero son necesarios para consolidar lo aprendido.} (Concession fait réel $\to$ indicatif \esp{cansan}, opposition \esp{pero}).
\item \esp{Quizás tenga razón.} (Doute $\to$ subjonctif présent \esp{tenga} de \esp{tener}). \textbf{Erreur fréquente} : \esp{*Quizás tiene razón} (change le sens en « il a probablement raison »).
\item \esp{Aunque no tenga ganas, iré a la fiesta.} (Hypothèse \esp{no tenga ganas} $\to$ subjonctif ; conséquence future \esp{iré}).
\end{enumerate}
\end{exoresolu}

\courssub{Contrastif français/espagnol et pièges de traduction (Thème / Version)}

\begin{propriete}[Tableau récapitulatif des correspondances et faux-amis]
\begin{center}
\begin{tabularx}{\linewidth}{|X|X|X|}
\hline
\rowcolor{popblueL} \textbf{Français} & \textbf{Espagnol correct} & \textbf{Piège / Faux ami à éviter} \\
\hline
Être d'accord (avec) & \esp{Estar de acuerdo (con)} & \esp{*Ser de acuerdo} ; \esp{*Estar de acuerdo \textbf{a}} \\
Avoir raison / tort & \esp{Tener razón / No tener razón} & \esp{*Haber razón} ; \esp{*Estar razón} \\
Être pour / contre & \esp{Estar a favor de / Estar en contra de} & \esp{*Ser para / contra} ; \esp{oponerse \textbf{a}} (pas \esp{de}) \\
Effectivement / En effet & \esp{Efectivamente / En efecto} & \esp{*Actualmente} = « actuellement » (temps) \\
Bien sûr / Évidemment & \esp{Por supuesto / Claro que sí / Desde luego} & \esp{*Por supuesto que no} = « Bien sûr que non » \\
Peut-être (doute) & \esp{Quizás / Tal vez} + \textbf{Subjonctif} & \esp{*Quizás + indicatif} (sauf quasi-certitude) \\
Il me semble que & \esp{Me parece que} + \textbf{Indicatif} & \esp{*Me parece \textbf{de} que} (queísmo) \\
\hline
\end{tabularx}
\end{center}
\end{propriete}

\begin{attention}[Les 5 erreurs qui coûtent des points au Bac (Check-list finale)]
\begin{enumerate}
\item \textbf{Accents oubliés} : \esp{razón}, \esp{opinión}, \esp{también}, \esp{además}, \esp{está}, \esp{prohibición}, \esp{móvil}, \esp{Camerún}, \esp{sí} (affirmation vs \esp{si} « si »), \esp{tú} (sujet vs \esp{tu} « ton »), \esp{él} (sujet vs \esp{el} « le »).
\item \textbf{Calque « être d'accord »} : \esp{\textbf{estar} de acuerdo \textbf{con}} (jamais \esp{ser}, jamais \esp{a/de}).
\item \textbf{Subjonctif ignoré} après opinion \textbf{niée} : \esp{No creo que \textbf{sea} verdad}, \esp{No pienso que \textbf{tenga} razón}, \esp{No me parece que \textbf{sea} justo}.
\item \textbf{Faux amis} : \esp{actualmente} = « actuellement » (pas « en effet ») ; \esp{la violencia} = « la violence » (pas \esp{violación} = « viol ») ; \esp{regular} = « réglementer/encadrer » (pas « régler » un problème) ; \esp{éxito} = « succès » (pas « exit »).
\item \textbf{Prépositions verbes d'opinion} : \esp{Estar a favor \textbf{de}} / \esp{Estar en contra \textbf{de}} / \esp{Opinión \textbf{sobre}} / \esp{Discrepar \textbf{de}}.
\end{enumerate}
\end{attention}

\courssub{Entraînement à la traduction : Thème (Français $\to$ Espagnol)}

\begin{methode}[Méthode du Thème : 5 étapes pour sécuriser la note]
\begin{enumerate}
\item \textbf{Analyser} : repérer le verbe d'opinion, la polarité (affirmatif/négatif), le registre.
\item \textbf{Transposer le verbe} : « Je pense que » $\to$ \esp{Creo que / Pienso que / Opino que} ; « Il me semble » $\to$ \esp{Me parece que} ; « Je suis pour/contre » $\to$ \esp{Estoy a favor de / en contra de}.
\item \textbf{Gérer le mode} : Affirmatif $\to$ \textbf{Indicatif} ; Négatif (\esp{no creo}, \esp{no pienso}, \esp{no opino}) $\to$ \textbf{Subjonctif}.
\item \textbf{Vérifier les collocations figées} : \esp{estar de acuerdo con}, \esp{tener razón}, \esp{estar a favor/en contra de}, \esp{oponerse a}.
\item \textbf{Relire} : accords (genre/nombre), accents, ponctuation \esp{¿ ? ¡ !}, ordre des mots (sujet souvent postposé après verbe d'opinion : \esp{Creo yo que...} possible mais \esp{Creo que...} standard).
\end{enumerate}
\end{methode}

\begin{exoresolu}[Thème : phrases d'opinion variées]
Énoncé. Traduisez en espagnol.
\begin{enumerate}
\item Je ne suis pas d'accord avec toi : le portable est utile en classe.
\item Nous pensons que cette mesure n'est pas juste.
\item Tu as raison, cependant il me semble que c'est compliqué.
\item Ils sont contre l'interdiction totale des réseaux sociaux.
\item Peut-être as-tu raison, mais j'en doute.
\end{enumerate}
\tcblower
\textbf{Raisonnement et solutions} :
\begin{enumerate}
\item Négation + opinion $\to$ \esp{No estoy de acuerdo contigo: el móvil es útil en clase.} (Indicatif \esp{es} car affirmation dans la proposition principale).
\item Opinion négative \esp{No pensamos que} $\to$ Subjonctif \esp{sea}. \esp{No pensamos que esta medida \textbf{sea} justa.}
\item Accord + concession. \esp{Tienes razón, sin embargo me parece que es más complicado.} (Indicatif après \esp{me parece que} affirmatif).
\item \esp{Están en contra de la prohibición total de las redes sociales.} (\esp{prohibición} fém., accent).
\item \esp{Quizás \textbf{tengas} razón, pero lo \textbf{dudo}.} (\esp{quizás} + subjonctif \esp{tengas} ; \esp{dudar} + indicatif car certitude du locuteur sur son propre doute).
\end{enumerate}
\end{exoresolu}

\courssub{Entraînement à la traduction : Version (Espagnol $\to$ Français)}

\begin{methode}[Méthode de la Version : 4 étapes pour un français naturel]
\begin{enumerate}
\item \textbf{Lecture globale} : identifier les locuteurs, le sujet, la tension (accord/désaccord/concession).
\item \textbf{Décodage lexical} : traduire les \textbf{formules figées} par leurs équivalents idiomatiques (pas mot à mot). \esp{Por supuesto} = « Bien sûr » ; \esp{De ninguna manera} = « Absolument pas » ; \esp{Sin embargo} = « Cependant ».
\item \textbf{Gestion de la syntaxe} : éviter le calque. \esp{Me parece que} $\to$ « Il me semble que » (pas « Il m'apparaît que »). \esp{Se distraen} $\to$ « Ils se laissent distraire / Ils sont distraits » (pas « Ils se distraient » sens réflexif rare).
\item \textbf{Relecture finale} : le français doit être fluide, sans trace d'hispanisme (ex. : « Il a raison » et non « Il a la raison »).
\end{enumerate}
\end{methode}

\begin{exoresolu}[Version : Débat citoyen (Radio Yaoundé)]
Énoncé. Traduisez en français le texte suivant.

\begin{textebox}[Debate ciudadano: ¿Uniforme obligatorio?]
— \esp{¿Qué opina de la obligatoriedad del uniforme en los liceos?}
— \esp{Estoy a favor. En primer lugar, fomenta la igualdad social. Además, evita la competencia por la ropa de marca.}
— \esp{Yo no estoy de acuerdo. Al contrario, creo que limita la libertad de expresión de los jóvenes.}
— \esp{Es verdad que la libertad es importante; sin embargo, la disciplina escolar es prioritaria.}
— \esp{Quizás tenga usted razón, pero hay que escuchar a los alumnos.}
\end{textebox}
\tcblower
\textbf{Analyse segmentée} :
\begin{itemize}
\item \esp{¿Qué opina de...?} = « Que pensez-vous de... ? » (formel \esp{usted} implicite).
\item \esp{Estoy a favor} = « Je suis pour / favorable ». \esp{Fomenta} (indicatif) = « favorise ». \esp{La igualdad social} = « l'égalité sociale ».
\item \esp{Evita la competencia} = « évite la concurrence/rivalité ». \esp{Ropa de marca} = « vêtements de marque ».
\item \esp{Al contrario} = « Au contraire ». \esp{Limita la libertad de expresión} = « limite la liberté d'expression ».
\item \esp{Es verdad que... ; sin embargo...} = « Il est vrai que... ; cependant... » (Concession + opposition).
\item \esp{Quizás tenga usted razón} = « Peut-être avez-vous raison » (Subjonctif \esp{tenga} après \esp{quizás} doute). \esp{Hay que + infinitif} = « Il faut / Il est nécessaire de ».
\end{itemize}
\textbf{Traduction rédigée (français naturel)} :
« — Que pensez-vous de l'obligation de l'uniforme dans les lycées ?
— Je suis favorable. D'abord, cela favorise l'égalité sociale. De plus, cela évite la rivalité pour les vêtements de marque.
— Moi, je ne suis pas d'accord. Au contraire, je pense que cela limite la liberté d'expression des jeunes.
— Il est vrai que la liberté est importante ; cependant, la discipline scolaire est prioritaire.
— Peut-être avez-vous raison, mais il faut écouter les élèves. »
\end{exoresolu}

\courssub{Savoir-faire : Rédiger un paragraphe d'opinion argumenté (Expression écrite)}

\begin{methode}[Structure type du paragraphe argumenté (Niveau Terminale)]
\begin{enumerate}
\item \textbf{Thèse (Position claire)} : \esp{En mi opinión, estoy a favor / en contra de...} / \esp{Desde mi punto de vista...} / \esp{Considero que...}
\item \textbf{Argument 1 + Connecteur} : \esp{En primer lugar / Primeramente / Para empezar,} + idée force + \esp{porque / ya que / puesto que} + explication.
\item \textbf{Argument 2 + Exemple concret (Contexte Cameroun/Hispanophone)} : \esp{Además / Por otra parte / Asimismo,} + idée + \esp{Por ejemplo,} + référence locale (\esp{en los liceos de Yaundé}, \esp{en Guinea Ecuatorial}, \esp{en mi barrio}).
\item \textbf{Concession + Réfutation (La touche excellence)} : \esp{Es cierto que... / Admito que... / Aunque...,} + \textbf{pero / sin embargo / no obstante} + contre-argument solide.
\item \textbf{Conclusion} : \esp{En conclusión / Por todo ello / Por eso,} + réaffirmation de la thèse + ouverture.
\end{enumerate}
\end{methode}

\begin{exoresolu}[Rédaction guidée : Sujet de type Bac]
Énoncé. \esp{« ¿Estás a favor o en contra del uso del teléfono móvil en el aula? Escribe un párrafo de 80--100 palabras. »}
\tcblower
\textbf{Plan appliqué} : Contre (distraction) $\to$ Arg 1 (baisse résultats) $\to$ Arg 2 + Exemple (lycées Yaoundé) $\to$ Concession (outil pédagogique possible) $\to$ Réfutation (encadrement strict nécessaire) $\to$ Conclusion.
\textbf{Production rédigée (Modèle d'excellence)} :
\esp{En mi opinión, estoy en contra del uso libre del móvil en el aula. En primer lugar, distrae a los alumnos y baja su rendimiento académico: muchos chicos chatean en vez de seguir la clase. Además, en los liceos de Yaundé se han multiplicado los casos de trampa en los exámenes por móvil. Es cierto que puede ser una herramienta pedagógica útil; sin embargo, sin una regulación estricta, el riesgo de dispersión es demasiado grande. Por eso creo que su uso debe estar prohibido durante las horas de clase, salvo autorización expresa del profesor.}
\textbf{Traduction} : « À mon avis, je suis contre l'usage libre du portable en classe. D'abord, il distrait les élèves et baisse leur rendement scolaire : beaucoup de gamins chattent au lieu de suivre le cours. De plus, dans les lycées de Yaoundé, les cas de triche à l'examen par portable se sont multipliés. Il est vrai que cela peut être un outil pédagogique utile ; cependant, sans régulation stricte, le risque de dispersion est trop grand. C'est pourquoi je pense que son usage doit être interdit pendant les heures de cours, sauf autorisation expresse du prof. »
\textbf{Analyse des critères de réussite} : Position nette (\esp{en contra}), connecteurs variés (\esp{en primer lugar, además, es cierto que... sin embargo, por eso}), subjonctif maîtrisé (\esp{sea} implicite dans \esp{autorización expresa}, \esp{haya} si passé), vocabulaire précis (\esp{rendimiento, trampa, regulación, dispersión}), exemple localisé (\esp{Yaundé}), longueur respectée.
\end{exoresolu}

\begin{experience}[Activité orale : « El Debate Relámpago » (15 min)]
\textbf{Objectif} : Réagir spontanément avec formules d'accord/désaccord + justification.
\textbf{Déroulement} :
\begin{enumerate}
\item La classe en cercle. Le prof lance une affirmation polémique (ex. \esp{« La tarea para casa debería abolirse. »}).
\item L'élève à droite \textbf{doit} accorder (\esp{Estoy de acuerdo porque...}) ; l'élève suivant \textbf{doit} désaccorder (\esp{No estoy de acuerdo, al contrario...}).
\item On alterne autour du cercle. Interdiction de répéter une justification déjà dite.
\item Variante « Niveau Bac » : imposer \esp{Quizás... subjuntivo} ou \esp{Aunque... indicativo/subjuntivo} pour la concession.
\end{enumerate}
\textbf{Évaluation} : Fluidité, correction grammaticale (subjonctif !), richesse lexicale, pertinence de l'argument.
\end{experience}

\courssub{Exercices d'entraînement (Devoir / Travail personnel)}

\begin{exercice}[Exercice 1 : Complétion grammaticale (Mode et Prépositions)]
Complétez par la forme verbale correcte (Indicatif ou Subjonctif) et la préposition si nécessaire.
\begin{enumerate}
\item No creo que esta ley \_\_\_\_\_\_ (ser) justa.
\item Estoy de acuerdo \_\_\_\_\_\_ ti: el español \_\_\_\_\_\_ (ser) muy útil.
\item Quizás el profesor \_\_\_\_\_\_ (tener) razón, pero yo \_\_\_\_\_\_ (dudar).
\item Aunque la tarea \_\_\_\_\_\_ (ser) aburrida, \_\_\_\_\_\_ (hacer) la necesito para aprender.
\item Me opongo \_\_\_\_\_\_ la violencia \_\_\_\_\_\_ los estadios.
\end{enumerate}

\end{exercice}

\begin{exercice}[Exercice 2 : Thème ciblé (Pièges classiques)]
Traduisez en espagnol en évitant les 5 erreurs du Bac.
\begin{enumerate}
\item Je suis d'accord avec elle : actuellement, l'espagnol est très important.
\item Ils n'ont pas raison. Effectivement, c'est difficile, mais pas impossible.
\item Nous sommes contre l'interdiction. Au contraire, c'est un droit.
\item Peut-être vient-il demain ? (Doute).
\item Il me semble qu'ils ont tort.
\end{enumerate}

\end{exercice}

\begin{exercice}[Exercice 3 : Version littéraire / Journalistique]
Traduisez en français le paragraphe suivant.
\end{exercice}

\begin{textebox}[Artículo de opinión: La voz de la juventud]
\esp{Los jóvenes de hoy no son apáticos, como a veces se dice. Al contrario, estamos muy comprometidos con el clima y la justicia social. Es verdad que usamos mucho las redes, pero eso no significa que no leamos. De ninguna manera aceptamos que nos llamen "generación de cristal". Quizás algunos adultos no entiendan nuestra forma de luchar, pero tenemos razón: el futuro es nuestro.}
\end{textebox}

\begin{exercice}[Exercice 4 : Expression écrite (Sujet type Bac)]
\esp{« ¿Estás a favor o en contra de que los alumnos elijan sus propias asignaturas en Terminale? Redacta un párrafo argumentado de 90--110 palabras. Utiliza al menos: una concesión (aunque + subjuntivo / es verdad que... pero), dos conectores de adición, un ejemplo concreto (Camerún o Guinea Ecuatorial) y una conclusión clara. »}
\end{exercice}

\begin{corrige}[Exercices 1 à 4]
\textbf{Exercice 1} :
\begin{enumerate}
\item \esp{sea} (Subjonctif : \esp{no creo que}).
\item \esp{con} / \esp{es} (Indicatif : affirmation).
\item \esp{tenga} (Subjonctif : \esp{quizás} doute) / \esp{dudo} (Indicatif : certitude du locuteur).
\item \esp{sea} (Subjonctif : \esp{aunque} hypothétique/futur) / \esp{hago} ou \esp{haré} (Indicatif conséquence). \textit{Note : \esp{Aunque sea aburrida, la hago / la haré}.}
\item \esp{a} / \esp{en} (\esp{oponerse a}, \esp{en los estadios}).
\end{enumerate}

\textbf{Exercice 2} :
\begin{enumerate}
\item \esp{Estoy de acuerdo con ella: actualmente el español es muy importante.} (\textbf{Piège} : \esp{actualmente} = actuellement ; \esp{con ella}).
\item \esp{No tienen razón. Efectivamente, es difícil, pero no imposible.} (\textbf{Piège} : \esp{tienen razón}, \esp{efectivamente} = effectivement).
\item \esp{Estamos en contra de la prohibición. Al contrario, es un derecho.} (\textbf{Piège} : \esp{estar en contra de}, \esp{al contrario}).
\item \esp{¿Quizás \textbf{venga} mañana?} (\textbf{Piège} : \esp{quizás} + subjonctif \esp{venga} de \esp{venir}).
\item \esp{Me parece que \textbf{tienen} tort / \textbf{no tienen} razón.} (\textbf{Piège} : \esp{me parece que} + indicatif ; \esp{tener razón/tort}).
\end{enumerate}

\textbf{Exercice 3 (Version)} :
« Les jeunes d'aujourd'hui ne sont pas apathiques, comme on le dit parfois. Au contraire, nous sommes très engagés pour le climat et la justice sociale. Il est vrai que nous utilisons beaucoup les réseaux (sociaux), mais cela ne veut pas dire que nous ne lisions pas. En aucune façon nous n'acceptons qu'on nous appelle "génération de cristal". Peut-être que certains adultes ne comprennent pas notre façon de lutter, mais nous avons raison : l'avenir est à nous. »
\textit{Notes : \esp{se dice} = impersonnel « on dit » ; \esp{generación de cristal} = terme médiatique réel (génération de cristal / snowflake generation) ; \esp{tenemos razón} = nous avons raison.}

\textbf{Exercice 4 (Production écrite - Proposition de corrigé type)} :
\esp{Desde mi punto de vista, estoy en contra de que los alumnos elijan libremente todas sus asignaturas en Terminale. En primer lugar, a los diecisiete años no siempre se tiene la madurez para diseñar un currículo equilibrado: se corre el riesgo de evitar las materias difíciles pero necesarias, como las matemáticas o la filosofía. Además, en el sistema camerunés, como en el de Guinea Ecuatorial, el bachillerato debe garantizar una cultura general común. Es cierto que la motivación aumenta cuando se elige; sin embargo, una elección total crearía lagunas graves para la universidad. Aunque algunos pidan más autonomía, pienso que un tronc commun solide con algunas optativas es la mejor solución. En conclusión, la libertad de elegir debe ser guiada, no absoluta.}
\textit{Vérification items obligatoires : Concession \esp{Es cierto que... sin embargo} / \esp{Aunque algunos pidan (subj.)}, Additions \esp{En primer lugar, Además}, Exemple \esp{Camerún, Guinea Ecuatorial}, Conclusion \esp{En conclusión}.}
\end{corrige}

\newpage

\coursec{Exercices}

\courssub{Je vérifie mes connaissances}

\begin{exercice}[QCM : la bonne formule]
Pour chaque situation, choisis la réponse correcte.
\begin{enumerate}
\item Pour exprimer un accord total, on dit :
\begin{enumerate}
\item \esp{de ninguna manera}
\item \esp{por supuesto}
\item \esp{me opongo}
\end{enumerate}
\item \esp{No creo que Juan \ldots{} razón.}
\begin{enumerate}
\item \esp{tiene}
\item \esp{tenga}
\item \esp{tendrá}
\end{enumerate}
\item \esp{Me opongo \ldots{} esta reforma.}
\begin{enumerate}
\item \esp{de}
\item \esp{a}
\item \esp{con}
\end{enumerate}
\item Pour atténuer une opinion, on peut dire :
\begin{enumerate}
\item \esp{hasta cierto punto}
\item \esp{al contrario}
\item \esp{efectivamente}
\end{enumerate}
\end{enumerate}
\end{exercice}

\begin{exercice}[Vrai ou faux ?]
Indique si chaque affirmation est vraie (V) ou fausse (F) et justifie ta réponse par une règle ou un exemple.
\begin{enumerate}
\item \esp{Estoy de acuerdo} se construit avec le verbe \esp{ser}.
\item Après \esp{no creo que}, le verbe de la subordonnée se met au subjonctif.
\item \esp{Tienes razón} s'emploie pour exprimer le désaccord poli.
\item \esp{Aunque} peut être suivi de l'indicatif ou du subjonctif selon le sens.
\item \esp{Me opongo a que vengan} est une phrase correcte.
\end{enumerate}
\end{exercice}

\begin{exercice}[Vocabulaire : accord et désaccord]
Complète chaque formule avec le mot manquant choisi dans la liste : \esp{razón -- acuerdo -- contrario -- supuesto -- manera -- verdad}.
\begin{enumerate}
\item Estoy de \ldots{} contigo.
\item Tienes \ldots{}, es una buena idea.
\item Por \ldots{}, puedes contar conmigo.
\item De ninguna \ldots{} aceptaré esa propuesta.
\item Al \ldots{}, yo pienso que es un error.
\item Es \ldots{} que el tema es difícil, pero hay que intentarlo.
\end{enumerate}
\end{exercice}

\begin{exercice}[Conjugaison : indicatif ou subjonctif ?]
Conjugue le verbe entre parenthèses au présent de l'indicatif ou du subjonctif selon la formule d'opinion.
\begin{enumerate}
\item Creo que el debate \ldots{} (ser) muy interesante.
\item No creo que ellos \ldots{} (tener) razón.
\item Es verdad que \ldots{} (haber) muchos jóvenes en la manifestación.
\item Me opongo a que se \ldots{} (cerrar) la biblioteca del barrio.
\item Estoy seguro de que ustedes \ldots{} (estar) de acuerdo conmigo.
\item Dudo que el gobierno \ldots{} (escuchar) a los estudiantes.
\end{enumerate}
\end{exercice}

\courssub{Je m'entraîne}

\begin{exercice}[Thème : phrases d'opinion]
Traduis en espagnol les phrases suivantes.
\begin{enumerate}
\item Je suis tout à fait d'accord avec toi.
\item Tu as raison, mais je pense qu'il faut réfléchir davantage.
\item Je ne crois pas qu'il ait raison de s'opposer à cette loi.
\item Bien sûr, je suis d'accord, mais seulement jusqu'à un certain point.
\item Je m'oppose à ce que les élèves utilisent le téléphone en classe.
\end{enumerate}
\end{exercice}

\begin{exercice}[Correction d'erreurs]
Chaque phrase contient une erreur typique de francophone. Repère-la et corrige-la.
\begin{enumerate}
\item \esp{Soy de acuerdo con la propuesta del director.}
\item \esp{Me opongo la idea de cerrar el mercado.}
\item \esp{No creo que tiene razón.}
\item \esp{Estoy de acuerdo a venir mañana.}
\item \esp{Tienes la razón, efectivamente.}
\item \esp{De ninguna manera yo estoy de acuerdo con eso.}
\item \esp{Aunque es caro, pero lo compro.}
\item \esp{Pienso que no es verdad que él sabe la respuesta.}
\end{enumerate}
\end{exercice}

\begin{exercice}[Texte à trous : un débat au lycée]
Complète le dialogue avec la formule appropriée selon le degré d'accord voulu : \esp{estoy totalmente de acuerdo -- de ninguna manera -- hasta cierto punto -- tienes razón -- me opongo -- es verdad que}.
\begin{textebox}[Un débat au lycée de Yaoundé]
\textbf{Ana :} Creo que los deberes son demasiado numerosos. ¿Qué piensas, Paul?\\
\textbf{Paul :} \ldots{} contigo : no tenemos tiempo para descansar.\\
\textbf{Ana :} Y además, pienso que deberían suprimir los exámenes.\\
\textbf{Paul :} ¡\ldots{}! Los exámenes son necesarios para evaluar a los alumnos.\\
\textbf{Ana :} \ldots{} los exámenes son útiles, pero hay demasiados.\\
\textbf{Paul :} \ldots{}, tienes razón en eso. Pero \ldots{} a que se supriman todos los controles.\\
\textbf{Ana :} Bueno, \ldots{} : podemos reducir el número de exámenes sin eliminarlos.
\end{textebox}
\end{exercice}

\begin{exercice}[Choix modal : atténuer et concéder]
Réécris chaque opinion brute en deux versions : d'abord une version \cle{atténuée} (avec \esp{creo que}, \esp{me parece que}, \esp{hasta cierto punto}...), puis une version \cle{concessive} (avec \esp{aunque} ou \esp{es verdad que... pero}).
\begin{enumerate}
\item \esp{El transporte público de Douala es malo.}
\item \esp{Los jóvenes no se interesan por la política.}
\item \esp{Esta ley es injusta.}
\item \esp{El fútbol es el único deporte importante del Camerún.}
\end{enumerate}
\end{exercice}

\courssub{Je me prépare au Bac}

\begin{exercice}[Compréhension et langue]
Lis le texte suivant, puis réponds aux questions.
\begin{textebox}[La participación de los jóvenes]
En muchos países, se dice que los jóvenes no se interesan por la vida pública. Sin embargo, es verdad que participan, pero de otra manera : al contrario de lo que piensan algunos adultos, prefieren actuar en asociaciones, en las redes sociales o en proyectos locales antes que votar en las elecciones.

« Me opongo a que se diga que la juventud es pasiva », declara una estudiante de Yaoundé. « Hasta cierto punto, entiendo las críticas, pero estoy segura de que mi generación quiere cambiar las cosas. No estoy de acuerdo con quienes afirman que no nos importa nada. »

Los especialistas creen que los gobiernos deben escuchar más a los jóvenes. Por supuesto, no todos están de acuerdo con esta idea, pero de ninguna manera se puede ignorar su voz.
\end{textebox}
\textbf{Questions de compréhension} (réponds en espagnol, avec tes propres mots) :
\begin{enumerate}
\item Según el texto, ¿cómo participan los jóvenes en la vida pública ?
\item ¿Qué opina la estudiante de Yaoundé sobre su generación ?
\item ¿Qué deben hacer los gobiernos, según los especialistas ?
\end{enumerate}
\textbf{Questions de langue} :
\begin{enumerate}
\item Relève dans le texte deux formules d'accord et deux formules de désaccord.
\item Justifie l'emploi du subjonctif dans \esp{me opongo a que se diga}.
\item Traduis en français la phrase : \esp{Al contrario de lo que piensan algunos adultos, prefieren actuar en asociaciones.}
\end{enumerate}
\end{exercice}

\begin{exercice}[Thème et version type Bac]
\textbf{A. Thème.} Traduis en espagnol les phrases suivantes.
\begin{enumerate}
\item Je ne crois pas qu'il ait raison de s'opposer à cette loi.
\item Il est vrai que le projet est coûteux, mais je suis d'accord avec le maire.
\item Nous nous opposons à ce que l'on ferme cette école.
\item Tu as raison : les jeunes doivent participer davantage à la vie du quartier.
\item Jusqu'à un certain point, je comprends ton opinion, mais je ne la partage pas.
\end{enumerate}
\textbf{B. Version.} Traduis en français l'extrait suivant.
\begin{textebox}[Para traducir]
Muchos creen que los jóvenes cameruneses solo piensan en el fútbol y en la música. Al contrario, es verdad que les interesan los estudios, pero también quieren participar en las decisiones de su país. Una alumna de Douala declara : « Me opongo a que los adultos decidan todo sin consultarnos. Hasta cierto punto, comprendo su prudencia, pero de ninguna manera acepto que ignoren nuestras ideas. Por supuesto, no todos los jóvenes piensan igual ; sin embargo, estoy de acuerdo con quienes piden más debates en los liceos. »
\end{textebox}
\end{exercice}

\begin{exercice}[Expression écrite type Bac]
Sujet : \esp{¿Estás de acuerdo con el uso del móvil en el aula?}

Rédige en espagnol un paragraphe argumenté de \textbf{60 à 80 mots} dans lequel tu :
\begin{enumerate}
\item exprimes clairement ton accord ou ton désaccord avec une formule adaptée (\esp{estoy de acuerdo}, \esp{me opongo a que}, \esp{no creo que}...) ;
\item fais au moins une \cle{concession} (\esp{aunque}, \esp{es verdad que... pero}) ;
\item utilises une formule de désaccord \cle{poli} envers l'opinion contraire ;
\item emploies au moins un verbe au \cle{subjonctif} correctement introduit.
\end{enumerate}
\end{exercice}

\newpage

\coursec{Corrigés des exercices}

\begin{corrige}[Exercice 1 — QCM : la bonne formule]
\begin{enumerate}
\item \textbf{b)} \esp{por supuesto}. C'est la formule d'accord total. \esp{De ninguna manera} exprime un refus catégorique et \esp{me opongo} un désaccord ferme.
\item \textbf{b)} \esp{tenga}. Après \esp{no creo que} (opinion niée), la subordonnée exige le \cle{subjonctif} : \esp{No creo que Juan tenga razón.} « Je ne crois pas que Juan ait raison. »
\item \textbf{b)} \esp{a}. Le verbe \esp{oponerse} se construit toujours avec la préposition \esp{a} : \esp{Me opongo a esta reforma.} « Je m'oppose à cette réforme. »
\item \textbf{a)} \esp{hasta cierto punto} « jusqu'à un certain point » : formule d'\cle{atténuation}. \esp{Al contrario} marque l'opposition, \esp{efectivamente} confirme.
\end{enumerate}
\end{corrige}

\begin{corrige}[Exercice 2 — Vrai ou faux ?]
\begin{enumerate}
\item \textbf{Faux.} \esp{Estar de acuerdo} se construit avec \esp{estar}, jamais avec \esp{ser} : \esp{Estoy de acuerdo contigo.} « Je suis d'accord avec toi. » (Calque interdit : *\esp{soy de acuerdo}.)
\item \textbf{Vrai.} \esp{No creo que} exprime un doute ou une opinion niée, donc la subordonnée est au subjonctif : \esp{No creo que venga.} « Je ne crois pas qu'il vienne. »
\item \textbf{Faux.} \esp{Tienes razón} « tu as raison » exprime l'\cle{accord}, non le désaccord. Pour un désaccord poli : \esp{No estoy de acuerdo}, \esp{entiendo tu punto de vista, pero...}
\item \textbf{Vrai.} \esp{Aunque} + indicatif = fait réel, admis : \esp{Aunque llueve, salgo.} « Bien qu'il pleuve, je sors. » \esp{Aunque} + subjonctif = fait hypothétique, non constaté : \esp{Aunque llueva, saldré.} « Même s'il pleuvait, je sortirais. »
\item \textbf{Vrai.} \esp{Oponerse a que} + subjonctif est la construction correcte : \esp{Me opongo a que vengan.} « Je m'oppose à ce qu'ils viennent. »
\end{enumerate}
\end{corrige}

\begin{corrige}[Exercice 3 — Vocabulaire : accord et désaccord]
\begin{enumerate}
\item Estoy de \textbf{acuerdo} contigo. « Je suis d'accord avec toi. »
\item Tienes \textbf{razón}, es una buena idea. « Tu as raison, c'est une bonne idée. »
\item Por \textbf{supuesto}, puedes contar conmigo. « Bien sûr, tu peux compter sur moi. »
\item De ninguna \textbf{manera} aceptaré esa propuesta. « En aucune façon je n'accepterai cette proposition. »
\item Al \textbf{contrario}, yo pienso que es un error. « Au contraire, je pense que c'est une erreur. »
\item Es \textbf{verdad} que el tema es difícil, pero hay que intentarlo. « Il est vrai que le sujet est difficile, mais il faut essayer. »
\end{enumerate}
\end{corrige}

\begin{corrige}[Exercice 4 — Conjugaison : indicatif ou subjonctif ?]
\begin{enumerate}
\item Creo que el debate \textbf{es} muy interesante. (Indicatif : \esp{creo que} affirme une opinion positive.)
\item No creo que ellos \textbf{tengan} razón. (Subjonctif : opinion niée.)
\item Es verdad que \textbf{hay} muchos jóvenes en la manifestación. (Indicatif : \esp{es verdad que} constate un fait réel.)
\item Me opongo a que se \textbf{cierre} la biblioteca del barrio. (Subjonctif après \esp{oponerse a que} ; \esp{cerrar} : e $\rightarrow$ ie, \esp{cierre}.)
\item Estoy seguro de que ustedes \textbf{están} de acuerdo conmigo. (Indicatif : certitude affirmée.)
\item Dudo que el gobierno \textbf{escuche} a los estudiantes. (Subjonctif : \esp{dudar que} exprime le doute.)
\end{enumerate}
\textbf{Règle récapitulative :} affirmation et certitude (\esp{creo que}, \esp{es verdad que}, \esp{estoy seguro de que}) $\rightarrow$ indicatif ; doute, négation de l'opinion, volonté ou refus (\esp{no creo que}, \esp{dudo que}, \esp{me opongo a que}) $\rightarrow$ subjonctif.
\end{corrige}

\begin{corrige}[Exercice 5 — Thème : phrases d'opinion]
\begin{enumerate}
\item \esp{Estoy totalmente de acuerdo contigo.} (« tout à fait » = \esp{totalmente} ; jamais *\esp{muy de acuerdo}.)
\item \esp{Tienes razón, pero pienso que hay que reflexionar más.} (\esp{hay que} + infinitif = « il faut ».)
\item \esp{No creo que tenga razón al oponerse a esta ley.} (Subjonctif \esp{tenga} après \esp{no creo que} ; « avoir raison de faire » = \esp{tener razón al} + infinitif ; \esp{oponerse a} + nom.)
\item \esp{Por supuesto, estoy de acuerdo, pero solo hasta cierto punto.}
\item \esp{Me opongo a que los alumnos usen el teléfono en clase.} (Subjonctif \esp{usen} après \esp{me opongo a que} ; « les élèves » = \esp{los alumnos} ou \esp{los estudiantes}.)
\end{enumerate}
\end{corrige}

\begin{corrige}[Exercice 6 — Correction d'erreurs]
\begin{enumerate}
\item \esp{Soy de acuerdo} $\rightarrow$ \esp{\textbf{Estoy} de acuerdo con la propuesta del director.} (On dit \esp{estar de acuerdo}, jamais *\esp{ser de acuerdo} : calque du français « être d'accord » avec le mauvais auxiliaire.)
\item \esp{Me opongo la idea} $\rightarrow$ \esp{Me opongo \textbf{a} la idea de cerrar el mercado.} (\esp{Oponerse} exige la préposition \esp{a}.)
\item \esp{No creo que tiene razón} $\rightarrow$ \esp{No creo que \textbf{tenga} razón.} (Subjonctif obligatoire après \esp{no creo que}.)
\item \esp{Estoy de acuerdo a venir} $\rightarrow$ \esp{Estoy de acuerdo \textbf{en} venir mañana} ou \esp{\textbf{De acuerdo}, vendré mañana.} (« D'accord pour + infinitif » se rend par \esp{estar de acuerdo en} + infinitif.)
\item \esp{Tienes la razón} $\rightarrow$ \esp{Tienes razón.} (Pas d'article dans \esp{tener razón} : calque de « avoir la raison ».)
\item \esp{De ninguna manera yo estoy de acuerdo} $\rightarrow$ \esp{De ninguna manera estoy de acuerdo con eso.} (Après une expression négative en tête de phrase, le sujet se place après le verbe : inversion.)
\item \esp{Aunque es caro, pero lo compro} $\rightarrow$ \esp{Aunque es caro, lo compro.} (On ne cumule pas \esp{aunque} et \esp{pero} : calque de « bien que... mais ».)
\item \esp{Pienso que no es verdad que él sabe} $\rightarrow$ \esp{Pienso que no es verdad que él \textbf{sepa} la respuesta.} (Après \esp{no es verdad que}, fait nié, le subjonctif s'impose.)
\end{enumerate}
\end{corrige}

\begin{corrige}[Exercice 7 — Texte à trous : un débat au lycée]
\begin{enumerate}
\item \textbf{Estoy totalmente de acuerdo} contigo : no tenemos tiempo para descansar. (Accord total : Paul renforce l'opinion d'Ana.)
\item ¡\textbf{De ninguna manera}! Los exámenes son necesarios... (Refus catégorique : Paul rejette la suppression des examens.)
\item \textbf{Es verdad que} los exámenes son útiles, pero hay demasiados. (Concession : Ana admet un fait avant de nuancer.)
\item \textbf{Hasta cierto punto}, tienes razón en eso. (Accord partiel, atténué.)
\item Pero \textbf{me opongo} a que se supriman todos los controles. (Désaccord ferme + subjonctif \esp{se supriman}.)
\item Bueno, \textbf{tienes razón} : podemos reducir el número de exámenes. (Accord final, marque d'ouverture.)
\end{enumerate}
\textbf{Traduction du dialogue complété :} « Ana : Je crois que les devoirs sont trop nombreux. Qu'en penses-tu, Paul ? — Paul : Je suis tout à fait d'accord avec toi : nous n'avons pas le temps de nous reposer. — Ana : Et en plus, je pense qu'ils devraient supprimer les examens. — Paul : En aucune façon ! Les examens sont nécessaires pour évaluer les élèves. — Ana : Il est vrai que les examens sont utiles, mais il y en a trop. — Paul : Jusqu'à un certain point, tu as raison là-dessus. Mais je m'oppose à ce qu'on supprime tous les contrôles. — Ana : Bon, tu as raison : nous pouvons réduire le nombre d'examens sans les éliminer. »
\end{corrige}

\begin{corrige}[Exercice 8 — Choix modal : atténuer et concéder]
Propositions (d'autres formulations correctes sont possibles) :
\begin{enumerate}
\item Atténuée : \esp{Me parece que el transporte público de Douala es bastante malo.} Concessive : \esp{Aunque el transporte público de Douala tiene algunos problemas, funciona todos los días.}
\item Atténuée : \esp{Creo que, hasta cierto punto, los jóvenes se interesan poco por la política.} Concessive : \esp{Es verdad que muchos jóvenes no votan, pero participan en asociaciones.}
\item Atténuée : \esp{En mi opinión, esta ley puede parecer injusta.} Concessive : \esp{Aunque esta ley tiene aspectos positivos, me parece injusta para los más pobres.}
\item Atténuée : \esp{Pienso que el fútbol es quizá el deporte más importante de Camerún.} Concessive : \esp{Es verdad que el fútbol es muy popular, pero no es el único deporte importante de Camerún.}
\end{enumerate}
\textbf{Méthode :} pour atténuer, on introduit un verbe d'opinion à la première personne (\esp{creo que}, \esp{me parece que}, \esp{pienso que}) ou un limiteur (\esp{hasta cierto punto}, \esp{bastante}, \esp{quizá}) ; pour concéder, on reconnaît d'abord une part de vérité (\esp{aunque} + indicatif, \esp{es verdad que... pero}), puis on maintient son point de vue.
\end{corrige}

\begin{corrige}[Exercice 9 — Compréhension et langue]
\textbf{Questions de compréhension} (réponses modèles) :
\begin{enumerate}
\item \esp{Según el texto, los jóvenes participan de otra manera : actúan en asociaciones, en las redes sociales o en proyectos locales, en lugar de votar.} « Selon le texte, les jeunes participent autrement : ils agissent dans des associations, sur les réseaux sociaux ou dans des projets locaux, au lieu de voter. »
\item \esp{La estudiante opina que su generación no es pasiva y que quiere cambiar las cosas.} « L'étudiante estime que sa génération n'est pas passive et qu'elle veut changer les choses. »
\item \esp{Los especialistas creen que los gobiernos deben escuchar más a los jóvenes.} « Les spécialistes pensent que les gouvernements doivent davantage écouter les jeunes. »
\end{enumerate}
\textbf{Questions de langue} :
\begin{enumerate}
\item Formules d'accord : \esp{es verdad que} (concession accordée), \esp{estoy de acuerdo con}, \esp{por supuesto}. Formules de désaccord : \esp{me opongo a que}, \esp{no estoy de acuerdo con}, \esp{de ninguna manera}, \esp{al contrario de}.
\item Dans \esp{me opongo a que se diga}, le subjonctif \esp{se diga} s'explique par la construction \esp{oponerse a que} + subjonctif : le refus porte sur un fait non admis, envisagé comme irrecevable. Toute volonté, tout refus ou sentiment visant le fait d'autrui entraîne le subjonctif.
\item Traduction : « Contrairement à ce que pensent certains adultes, ils préfèrent agir dans des associations. » (Attention : \esp{al contrario de lo que} = « contrairement à ce que », et non « au contraire de ce qui ».)
\end{enumerate}
\end{corrige}

\begin{corrige}[Exercice 10 — Thème et version type Bac]
\textbf{A. Thème} (barème indicatif : 1 point par phrase, dont 0,5 pour la structure et 0,5 pour la grammaire et le lexique) :
\begin{enumerate}
\item \esp{No creo que tenga razón al oponerse a esta ley.} (Subjonctif \esp{tenga} ; « avoir raison de faire » = \esp{tener razón al} + infinitif.)
\item \esp{Es verdad que el proyecto es caro, pero estoy de acuerdo con el alcalde.} (« Coûteux » = \esp{caro} ; « le maire » = \esp{el alcalde}, mot à retenir avec son article masculin.)
\item \esp{Nos oponemos a que cierren esta escuela} ou \esp{a que se cierre esta escuela.} (Subjonctif après \esp{oponerse a que} ; \esp{cerrar} : e $\rightarrow$ ie.)
\item \esp{Tienes razón : los jóvenes deben participar más en la vida del barrio.} (\esp{deber} + infinitif = « devoir » ; « davantage » = \esp{más}.)
\item \esp{Hasta cierto punto, comprendo tu opinión, pero no la comparto.} (« Partager une opinion » = \esp{compartir}, pas *\esp{partir}.)
\end{enumerate}
\textbf{B. Version} (barème indicatif : 5 points ; 1 point par segment de sens, moins 0,25 par erreur de grammaire ou de contresens) :
« Beaucoup croient que les jeunes Camerounais ne pensent qu'au football et à la musique. Au contraire, il est vrai que les études les intéressent, mais ils veulent aussi participer aux décisions de leur pays. Une élève de Douala déclare : “Je m'oppose à ce que les adultes décident de tout sans nous consulter. Jusqu'à un certain point, je comprends leur prudence, mais je n'accepte en aucune façon qu'ils ignorent nos idées. Bien sûr, tous les jeunes ne pensent pas de la même manière ; cependant, je suis d'accord avec ceux qui demandent plus de débats dans les lycées.” »

\textbf{Points de vigilance :} \esp{solo... en} = « ne... que » ; \esp{les interesan} : le verbe s'accorde avec \esp{los estudios} (« les études les intéressent ») ; \esp{decidan} et \esp{ignoren} sont des subjonctifs à rendre par « décident » / « qu'ils ignorent » ; \esp{liceo} = « lycée » (espagnol d'Amérique et de Guinée équatoriale, fréquent aussi au Cameroun).
\end{corrige}

\begin{corrige}[Exercice 11 — Expression écrite type Bac]
\textbf{Proposition de rédaction modèle} (environ 75 mots) :
\begin{textebox}[Production modèle]
Me opongo a que los alumnos usen el móvil en el aula. Es verdad que el teléfono puede ser útil para buscar información, pero creo que distrae a los estudiantes y perjudica el trabajo en clase. Entiendo la opinión de quienes lo defienden ; sin embargo, no estoy de acuerdo con ellos, porque no creo que sea posible concentrarse con un móvil en la mano. En mi opinión, hay que dejarlo en la mochila.
\end{textebox}
\textbf{Traduction :} « Je m'oppose à ce que les élèves utilisent le portable en classe. Il est vrai que le téléphone peut être utile pour chercher des informations, mais je crois qu'il distrait les étudiants et nuit au travail en classe. Je comprends l'opinion de ceux qui le défendent ; cependant, je ne suis pas d'accord avec eux, parce que je ne crois pas qu'il soit possible de se concentrer avec un portable à la main. À mon avis, il faut le laisser dans le sac. »

\textbf{Barème indicatif} (sur 5 points) : prise de position claire avec formule adaptée (1 pt) ; concession correctement construite (1 pt) ; désaccord poli envers l'opinion contraire (1 pt) ; subjonctif correctement introduit et conjugué (1 pt) ; correction linguistique globale et respect du nombre de mots (1 pt).

\textbf{Points de vigilance :}
\begin{itemize}
\item Vérifier la construction : \esp{me opongo a que} + subjonctif (\esp{usen}), \esp{no creo que} + subjonctif (\esp{sea}).
\item Ne pas cumuler \esp{aunque} et \esp{pero} dans la même phrase.
\item Employer \esp{estar de acuerdo} (et non *\esp{ser de acuerdo}) et \esp{tener razón} sans article.
\item Compter les mots : entre 60 et 80 ; au-delà ou en deçà, on perd des points.
\end{itemize}
\end{corrige}

\newpage

\coursec{Fiche bilan}

\begin{popfigure}
\begin{tikzpicture}[mindmap, grow cyclic, every node/.style={concept, text=white, font=\small\bfseries},
root concept/.append style={concept color=popdark, font=\bfseries},
level 1/.append style={level distance=3.6cm, sibling angle=51.4},
level 2/.append style={level distance=2.2cm, sibling angle=45, font=\scriptsize}]
\node[root concept]{Acuerdo y desacuerdo}
  child[concept color=popblue]{ node{Formules d'accord}
    child{ node{estoy de acuerdo} }
    child{ node{tienes razón} }
    child{ node{por supuesto} }
  }
  child[concept color=poporange]{ node{Formules de désaccord}
    child{ node{no estoy de acuerdo} }
    child{ node{de ninguna manera} }
    child{ node{me opongo a} }
  }
  child[concept color=popgreen]{ node{Modalisation}
    child{ node{quizás, tal vez} }
    child{ node{creo que + ind.} }
  }
  child[concept color=poppurple]{ node{Concession}
    child{ node{aunque + ind./subj.} }
    child{ node{es cierto que... pero} }
  }
  child[concept color=poppink]{ node{Subjonctif}
    child{ node{no creo que + subj.} }
    child{ node{dudo que + subj.} }
  }
  child[concept color=popgold]{ node{Pièges}
    child{ node{estar (pas ser) de acuerdo} }
    child{ node{razón : tener, pas être} }
  }
  child[concept color=popblue!70]{ node{Méthode}
    child{ node{repérer l'opinion} }
    child{ node{traduire puis nuancer} }
  };
\end{tikzpicture}
\legende{Carte mentale de la leçon : exprimer l'accord et le désaccord en espagnol}
\end{popfigure}

\begin{bilan}
\textbf{1. Lexique clé (à connaître par cœur)}
\begin{itemize}
  \item \esp{Estoy (totalmente) de acuerdo (contigo).} : je suis (tout à fait) d'accord (avec toi).
  \item \esp{Tienes razón. / No te falta razón.} : tu as raison. / tu n'as pas tort.
  \item \esp{Por supuesto. / Efectivamente. / Desde luego.} : bien sûr. / effectivement. / évidemment.
  \item \esp{No estoy de acuerdo. / De ninguna manera. / Al contrario.} : je ne suis pas d'accord. / en aucun cas. / au contraire.
  \item \esp{Me opongo a esta idea.} : je m'oppose à cette idée.
  \item \esp{En mi opinión / a mi juicio / desde mi punto de vista} : à mon avis.
\end{itemize}

\textbf{2. Règles de grammaire à retenir}
\begin{center}
\begin{tabularx}{\linewidth}{|X|X|}
\hline
\rowcolor{popblueL} Règle & Exemple \\
\hline
On dit \esp{estar de acuerdo} (jamais \esp{ser}) : accord = état. & \esp{Estoy de acuerdo con Ana.} \\
\hline
\esp{Tener razón} : « avoir raison » se construit avec \esp{tener}. & \esp{Los alumnos tienen razón.} \\
\hline
Verbe d'opinion affirmatif + \textbf{indicatif}. & \esp{Creo que es justo.} \\
\hline
Verbe d'opinion à la forme négative + \textbf{subjonctif}. & \esp{No creo que sea justo.} \\
\hline
\esp{Aunque} + indicatif (fait réel) ; + subjonctif (hypothèse). & \esp{Aunque llueve, salgo. / Aunque llueva, salgo.} \\
\hline
\esp{Oponerse a} + nom ou infinitif. & \esp{Me opongo a cambiar la ley.} \\
\hline
\end{tabularx}
\end{center}

\textbf{3. Modalisation express}
\begin{itemize}
  \item Atténuer : \esp{quizás, tal vez, puede que + subj., me parece que, diria que}.
  \item Concéder : \esp{es cierto que... pero..., aunque..., sin embargo, no obstante}.
  \item Renforcer : \esp{sin duda, desde luego, está claro que}.
\end{itemize}

\textbf{4. Méthode de traduction (thème/version)}
\begin{enumerate}
  \item Repérer qui parle et le degré d'opinion (accord, désaccord, doute).
  \item Choisir la formule adaptée au registre (familier : \esp{tienes razón} ; soutenu : \esp{me opongo a}).
  \item Vérifier le mode : indicatif après opinion affirmative, subjonctif après négation ou doute.
  \item Relire : verbe \esp{estar}, préposition \esp{con/a}, accents et ponctuation ¿ ¡.
\end{enumerate}
\end{bilan}

\begin{aretenir}[Checklist avant le Bac]
\begin{itemize}
  \item $\square$ Je sais dire « je suis d'accord » avec \esp{estar de acuerdo} (et jamais \esp{ser}).
  \item $\square$ Je traduis « tu as raison » par \esp{tienes razón} sans calquer « être ».
  \item $\square$ Je connais au moins trois formules d'accord et trois de désaccord.
  \item $\square$ Je sais employer \esp{de ninguna manera} et \esp{al contrario} à bon escient.
  \item $\square$ Je mets l'indicatif après \esp{creo que} et le subjonctif après \esp{no creo que}.
  \item $\square$ Je sais atténuer une opinion avec \esp{quizás}, \esp{tal vez}, \esp{me parece que}.
  \item $\square$ Je sais concéder avec \esp{aunque} et \esp{es cierto que... pero}.
  \item $\square$ Je conjugue correctement \esp{oponerse a} + infinitif.
  \item $\square$ Je distingue les registres : familier, courant, soutenu.
  \item $\square$ Je n'oublie pas les signes ¿ et ¡ dans une phrase exclamative ou interrogative.
  \item $\square$ Je vérifie les accents : \esp{opinión, razón, está, además}.
  \item $\square$ Je relis toujours ma traduction en comparant avec la phrase française.
\end{itemize}
\end{aretenir}

\findecours

\end{document}
